\section{Tomita--Takesaki Theory}\label{sctn:tomita-takesaki}

Let $\mathcal{M}$ be a von Neumann algebra on a Hilbert space with a vector $\Omega$ that is cyclic and separating for $\mathcal{M}$. The map $a\Omega \mapsto a^*\Omega$ on $\mathcal{M}\Omega$ is antilinear, densely defined and closable, and the polar decomposition of its closure $S = J\Delta^{1/2}$ produces the \emph{modular operator} $\Delta$, a positive, self-adjoint and injective operator, and the \emph{modular conjugation} $J$, an antiunitary involution with $J\Delta J = \Delta^{-1}$. Tomita's theorem states that $\Delta^{it}\mathcal{M}\Delta^{-it} = \mathcal{M}$ for every real $t$ and that $J\mathcal{M}J = \mathcal{M}'$. The section is pure operator algebra; it involves no net.

The section first constructs the Tomita operator $S$ (Subsection~\ref{sbsctn:tomita-operator}) and then $\Delta$ and $J$ (Subsection~\ref{sbsctn:modular-operator-conjugation}). It continues with the auxiliary results that J.~Sorce's short proof of Tomita's theorem (\emph{J. Funct. Anal.}, 2024) cites, Carlson's theorem from complex analysis, Sorce's proof of $\Delta^{it}\mathcal{M}\Delta^{-it} = \mathcal{M}$, and the proof of $J\mathcal{M}J = \mathcal{M}'$ following Bratteli and Robinson, \emph{Operator Algebras and Quantum Statistical Mechanics~1}, \S2.5.2.

Throughout, $\mathbf{H}$ denotes a complex Hilbert space, and its inner product $\left< \cdot, \cdot \right>$ is conjugate-linear in the first slot and linear in the second. Section~\ref{sctn:unbounded-spectral-theorems} assumes $\mathbf{H}$ separable, but the results of that section used here hold without separability, as their Lean forms do. The theory of unbounded operators in that section concerns linear operators. The Tomita operator is antilinear, and the section reduces it to a linear operator by composing with a fixed antiunitary involution $C$, a \emph{reference conjugation}. The reference conjugation is an implementation device: the Tomita operator, the modular operator and the modular conjugation do not depend on it (Lemma~\ref{lmm:tomita-operator-independent-of-c} and Theorem~\ref{thrm:modular-objects-independent-of-c}).

\subsection{The Tomita Operator}\label{sbsctn:tomita-operator}

\subsubsection{Antiunitary Operators and a Reference Conjugation}

\begin{definition}[Antilinear and Antiunitary Operators]
    \label{def:antiunitary-operator}
    \lean{LinearIsometryEquiv}
    \uses{def:unitary-operator}
    \leanok
    A map $W : \mathbf{H} \to \mathbf{H}$ is \emph{antilinear} if $W(\alpha\phi + \psi) = \overline{\alpha}\, W\phi + W\psi$ for all $\alpha \in \mathbb{C}$ and $\phi, \psi \in \mathbf{H}$. An \emph{antiunitary operator} is an antilinear bijection $W$ of $\mathbf{H}$ onto $\mathbf{H}$ with
    \begin{align}
        \left< W\phi, W\psi \right> = \overline{\left< \phi, \psi \right>} = \left< \psi, \phi \right> \qquad \text{for all } \phi, \psi \in \mathbf{H}.
    \end{align}
    An \emph{antilinear operator} on $\mathbf{H}$ is a pair of a subspace $\text{Dom}(A) \subset \mathbf{H}$ and an antilinear map $A : \text{Dom}(A) \to \mathbf{H}$; its graph $\Gamma(A) = \{ (\psi, A\psi) \mid \psi \in \text{Dom}(A) \}$ is a real subspace of $\mathbf{H} \times \mathbf{H}$.

    In the Lean, an antiunitary operator is Mathlib's conjugate-linear isometric equivalence $\mathbf{H} \simeq_{\ell i \star}[\mathbb{C}]\, \mathbf{H}$ (\texttt{LinearIsometryEquiv (starRingEnd \ensuremath{\mathbb{C}})}), and a conjugate-linear norm-preserving map whose bijectivity is not yet known is $\mathbf{H} \to_{\ell i \star}[\mathbb{C}]\, \mathbf{H}$. Mathlib supplies conjugate-linearity, norm preservation, continuity and the inverse (\texttt{symm}); the inner-product identity above is Lemma~\ref{lmm:conjugate-linear-isometry-inner}, and Lemma~\ref{lmm:antilinear-isometry-antiunitary} identifies the two notions. Identities such as $C^2 = \mathbf{1}$ and $J^2 = \mathbf{1}$ are stated pointwise, $C(C\phi) = \phi$ for all $\phi$.
\end{definition}

\begin{lemma}[Antiunitary Operators are Isometric]
    \label{lmm:antiunitary-isometric}
    \lean{Physicslib4.Modular.antiunitary_isometric}
    \uses{def:antiunitary-operator}
    \leanok
    An antiunitary operator $W$ is isometric and continuous, and $W^{-1}$ is antiunitary. For a Lean \texttt{LinearIsometryEquiv} these are supplied by Mathlib (\texttt{norm\_map}, \texttt{continuous}, \texttt{symm}).
\end{lemma}
\begin{proof}
    \leanok
    The case $\phi = \psi$ of the defining identity gives $\|W\phi\|^2 = \|\phi\|^2$, and $\|W\phi - W\psi\| = \|W(\phi - \psi)\| = \|\phi - \psi\|$ gives continuity. Applying $W^{-1}$ to $W(\overline{\alpha}\phi + \psi) = \alpha W\phi + W\psi$ shows that $W^{-1}$ is antilinear, and the defining identity for $W^{-1}\phi, W^{-1}\psi$ reads $\left< \phi, \psi \right> = \overline{\left< W^{-1}\phi, W^{-1}\psi \right>}$.
\end{proof}

\begin{lemma}[Composition Rules for Antiunitary Operators]
    \label{lmm:antiunitary-composition}
    \lean{Physicslib4.Modular.antiunitary_composition}
    \uses{def:antiunitary-operator, def:unitary-operator}
    \leanok
    Let $W, W_1, W_2$ be antiunitary and $U$ unitary. Then $WU$ and $UW$ are antiunitary, and $W_1W_2$ is unitary.
\end{lemma}
\begin{proof}
    \leanok
    Each composite is a bijection; a composite of two antilinear maps is linear, and a composite of a linear and an antilinear map is antilinear. The inner-product identities follow by applying the defining identities of the two factors in turn; two conjugations cancel.
\end{proof}

\begin{lemma}[A Conjugate-Linear Isometry Conjugates Inner Products]
    \label{lmm:conjugate-linear-isometry-inner}
    \lean{Physicslib4.Modular.inner_map_map_of_conjLinearIsometry}
    \uses{def:antiunitary-operator}
    \leanok
    Let $W : \mathbf{H} \to \mathbf{H}$ be antilinear and norm-preserving, not necessarily bijective. Then $\left< W\phi, W\psi \right> = \overline{\left< \phi, \psi \right>}$ for all $\phi, \psi \in \mathbf{H}$.
\end{lemma}
\begin{proof}
    \leanok
    \uses{prpstn:polarization-identity}
    The polarization identity (Proposition~\ref{prpstn:polarization-identity}) expresses $\left< \phi, \psi \right>$ through the norms $\|\psi + i^k\phi\|$. Since $W$ sends $\psi + i^k\phi$ to $W\psi + (-i)^kW\phi$ and preserves norms, reindexing the polarization sum gives $\left< W\phi, W\psi \right> = \overline{\left< \phi, \psi \right>}$.
\end{proof}

\begin{lemma}[An Antilinear Isometric Bijection is Antiunitary]
    \label{lmm:antilinear-isometry-antiunitary}
    \lean{Physicslib4.Modular.exists_antiunitary_of_bijective}
    \uses{def:antiunitary-operator}
    \leanok
    An antilinear norm-preserving bijection $W$ of $\mathbf{H}$ is antiunitary. In particular every Lean \texttt{LinearIsometryEquiv (starRingEnd \ensuremath{\mathbb{C}})} is antiunitary in the sense of Definition~\ref{def:antiunitary-operator}.
\end{lemma}
\begin{proof}
    \leanok
    \uses{lmm:conjugate-linear-isometry-inner}
    The inner-product identity is Lemma~\ref{lmm:conjugate-linear-isometry-inner}; bijectivity is assumed.
\end{proof}

\begin{definition}[Reference Conjugation]
    \label{def:reference-conjugation}
    \lean{Physicslib4.Modular.referenceConjugation}
    \uses{def:antiunitary-operator}
    \leanok
    Fix a Hilbert basis $(e_i)_{i \in I}$ of $\mathbf{H}$ (it exists by Mathlib's \texttt{exists\_hilbertBasis}; no separability is needed here). The \emph{reference conjugation} is the map $C : \mathbf{H} \to \mathbf{H}$ that conjugates coordinates,
    \begin{align}
        C\Big( \sum_{i \in I} c_i e_i \Big) = \sum_{i \in I} \overline{c_i}\, e_i, \qquad (c_i) \in \ell^2(I).
    \end{align}
    Equivalently, $C = R^{-1} \circ \kappa \circ R$, where $R : \mathbf{H} \to \ell^2(I)$ is the unitary coordinate map of the basis (\texttt{HilbertBasis.repr}) and $\kappa$ is coordinatewise complex conjugation on $\ell^2(I)$.
\end{definition}

\begin{lemma}[Properties of the Reference Conjugation]
    \label{lmm:reference-conjugation-properties}
    \lean{Physicslib4.Modular.referenceConjugation_properties}
    \uses{def:reference-conjugation, def:antiunitary-operator}
    \leanok
    The reference conjugation $C$ is antilinear and isometric, satisfies $C(C\phi) = \phi$ for all $\phi$, and $\left< C\phi, C\psi \right> = \overline{\left< \phi, \psi \right>}$ for all $\phi, \psi \in \mathbf{H}$. In particular $C$ is an antiunitary involution.
\end{lemma}
\begin{proof}
    \leanok
    \uses{lmm:conjugate-linear-isometry-inner, lmm:antilinear-isometry-antiunitary}
    Coordinatewise conjugation $\kappa$ on $\ell^2(I)$ is antilinear, preserves $\sum_i |c_i|^2$ and satisfies $\kappa(\kappa c) = c$; conjugating by the unitary $R$ transfers these three properties to $C$. Lemma~\ref{lmm:conjugate-linear-isometry-inner} gives the inner-product identity, and $C(C\phi) = \phi$ makes $C$ a bijection, hence antiunitary (Lemma~\ref{lmm:antilinear-isometry-antiunitary}).
\end{proof}

\subsubsection{The Setting}

The setting follows the project's commutant convention (\ref{def:local-von-neumann}): for a set $\mathcal{S} \subset \mathcal{B}(\mathbf{H})$, the commutant $\mathcal{S}'$ is the centralizer $\{ b \in \mathcal{B}(\mathbf{H}) \mid ba = ab \text{ for all } a \in \mathcal{S} \}$.

\begin{definition}[Von Neumann Algebra in the Commutant Convention]
    \label{def:von-neumann-algebra-commutant-convention}
    \lean{Physicslib4.Modular.IsVonNeumannSet}
    \uses{def:local-von-neumann, lmm:bicommutant-of-selfadjoint-is-von-neumann}
    \leanok
    A \emph{von Neumann algebra on $\mathbf{H}$} is a set $\mathcal{M} \subset \mathcal{B}(\mathbf{H})$ that is self-adjoint ($a \in \mathcal{M}$ implies $a^* \in \mathcal{M}$) and equal to its bicommutant, $\mathcal{M} = \mathcal{M}''$. Every local von Neumann algebra $R(\mathbf{B})$ of \ref{def:local-von-neumann} is one, by Lemma~\ref{lmm:bicommutant-of-selfadjoint-is-von-neumann}.
\end{definition}

\begin{lemma}[The Commutant of a Von Neumann Algebra]
    \label{lmm:commutant-von-neumann-algebra}
    \lean{Physicslib4.Modular.isVonNeumannSet_centralizer}
    \uses{def:von-neumann-algebra-commutant-convention}
    \leanok
    If $\mathcal{M}$ is a von Neumann algebra on $\mathbf{H}$, then so is $\mathcal{M}'$, and $(\mathcal{M}')' = \mathcal{M}$. Both $\mathcal{M}$ and $\mathcal{M}'$ are unital $*$-subalgebras of $\mathcal{B}(\mathbf{H})$.
\end{lemma}
\begin{proof}
    \leanok
    \uses{lmm:bicommutant-of-selfadjoint-is-von-neumann}
    $(\mathcal{M}')' = \mathcal{M}'' = \mathcal{M}$ by hypothesis. By Lemma~\ref{lmm:bicommutant-of-selfadjoint-is-von-neumann}, the commutant of a self-adjoint set is a self-adjoint subalgebra, and $\mathcal{M}''' = \mathcal{M}'$; so $\mathcal{M}'$ is self-adjoint and equal to its bicommutant, and $\mathcal{M} = (\mathcal{M}')'$ is a self-adjoint subalgebra. A centralizer contains $\mathbf{1}$.
\end{proof}

\begin{definition}[Separating Vector]
    \label{def:separating-vector}
    \lean{Physicslib4.Modular.IsSeparatingVector}
    \leanok
    A vector $\Omega \in \mathbf{H}$ is \emph{separating} for a set $\mathcal{S} \subset \mathcal{B}(\mathbf{H})$ if $a \in \mathcal{S}$ and $a\Omega = 0$ imply $a = 0$.
\end{definition}

\begin{definition}[Cyclic and Separating Vector]
    \label{def:cyclic-separating-vector}
    \lean{Physicslib4.Modular.IsCyclicSeparatingVector}
    \uses{def:von-neumann-algebra-commutant-convention, def:cyclic-vector}
    \leanok
    Let $\mathcal{M}$ be a von Neumann algebra on $\mathbf{H}$. A vector $\Omega \in \mathbf{H}$ is \emph{cyclic and separating} for $\mathcal{M}$ if $\Omega$ is cyclic (\ref{def:cyclic-vector}, for the inclusion representation) both for $\mathcal{M}$ and for $\mathcal{M}'$: the subspaces $\mathcal{M}\Omega$ and $\mathcal{M}'\Omega$ are dense in $\mathbf{H}$.
\end{definition}

\begin{lemma}[Cyclic for the Commutant Implies Separating]
    \label{lmm:cyclic-commutant-separating}
    \lean{Physicslib4.Modular.isSeparatingVector_of_isCyclicSeparatingVector}
    \uses{def:cyclic-separating-vector, def:separating-vector}
    \leanok
    If $\Omega$ is cyclic and separating for $\mathcal{M}$, then $\Omega$ is separating for $\mathcal{M}$ and for $\mathcal{M}'$.
\end{lemma}
\begin{proof}
    \leanok
    \uses{lmm:commutant-von-neumann-algebra}
    Let $a \in \mathcal{M}$ with $a\Omega = 0$. For $b \in \mathcal{M}'$, $a b\Omega = b a\Omega = 0$, so the bounded operator $a$ vanishes on the dense subspace $\mathcal{M}'\Omega$ and hence $a = 0$; this is the abstract mechanism of Theorem~\ref{thrm:statistical-independence}. The same argument with $\mathcal{M}$ and $\mathcal{M}'$ exchanged, using $(\mathcal{M}')' = \mathcal{M}$ (Lemma~\ref{lmm:commutant-von-neumann-algebra}) and cyclicity for $\mathcal{M}$, gives the second claim.
\end{proof}

For the rest of the section, $\mathcal{M}$ is a von Neumann algebra on $\mathbf{H}$, $\Omega$ is cyclic and separating for $\mathcal{M}$, and $C$ is a reference conjugation.

\subsubsection{The Antilinear Pre-Operators}

\begin{lemma}[The Tomita Pre-Operators are Well Defined]
    \label{lmm:tomita-pre-operators-well-defined}
    \lean{Physicslib4.Modular.tomitaPre_wellDefined}
    \uses{def:cyclic-separating-vector}
    \leanok
    If $a, b \in \mathcal{M}$ and $a\Omega = b\Omega$, then $a^*\Omega = b^*\Omega$. The same holds for $a, b \in \mathcal{M}'$.
\end{lemma}
\begin{proof}
    \leanok
    \uses{lmm:cyclic-commutant-separating, lmm:commutant-von-neumann-algebra}
    $a - b \in \mathcal{M}$ and $(a - b)\Omega = 0$, so $a = b$ because $\Omega$ is separating for $\mathcal{M}$ (Lemma~\ref{lmm:cyclic-commutant-separating}); hence $a^* = b^*$. For $\mathcal{M}'$ the argument is the same.
\end{proof}

\begin{definition}[The Tomita Pre-Operators $S_0$ and $F_0$]
    \label{def:tomita-pre-operators}
    \lean{Physicslib4.Modular.tomitaPre}
    \uses{def:antiunitary-operator, def:cyclic-separating-vector, lmm:tomita-pre-operators-well-defined, lmm:commutant-von-neumann-algebra}
    \leanok
    The antilinear operators $S_0$ and $F_0$ on $\mathbf{H}$ are
    \begin{align}
        \text{Dom}(S_0) = \mathcal{M}\Omega, \quad S_0(a\Omega) = a^*\Omega \quad (a \in \mathcal{M}); \qquad \text{Dom}(F_0) = \mathcal{M}'\Omega, \quad F_0(b\Omega) = b^*\Omega \quad (b \in \mathcal{M}').
    \end{align}
    They are well defined by Lemma~\ref{lmm:tomita-pre-operators-well-defined}, take values in $\mathcal{M}\Omega$ and $\mathcal{M}'\Omega$ because both algebras are self-adjoint (Lemma~\ref{lmm:commutant-von-neumann-algebra}), and are antilinear because $(\alpha a + b)^* = \overline{\alpha} a^* + b^*$.
\end{definition}

\begin{lemma}[Dense Domains and Involutivity of $S_0$ and $F_0$]
    \label{lmm:tomita-pre-operators-dense-involutive}
    \lean{Physicslib4.Modular.tomitaPre_dense_involutive}
    \uses{def:tomita-pre-operators}
    \leanok
    $\text{Dom}(S_0)$ and $\text{Dom}(F_0)$ are dense in $\mathbf{H}$. Moreover $S_0$ maps $\mathcal{M}\Omega$ onto itself with $S_0 S_0 x = x$ for all $x \in \mathcal{M}\Omega$, and likewise for $F_0$ on $\mathcal{M}'\Omega$.
\end{lemma}
\begin{proof}
    \leanok
    \uses{def:cyclic-separating-vector}
    Density is the cyclicity of $\Omega$ for $\mathcal{M}$ and $\mathcal{M}'$ (\ref{def:cyclic-separating-vector}). For $a \in \mathcal{M}$, $S_0 S_0 (a\Omega) = S_0(a^*\Omega) = a^{**}\Omega = a\Omega$, so $S_0$ is its own inverse on $\mathcal{M}\Omega$ and in particular onto.
\end{proof}

\begin{lemma}[The Adjoint Relation between $S_0$ and $F_0$]
    \label{lmm:tomita-adjoint-relation}
    \lean{Physicslib4.Modular.tomitaPre_adjoint_relation}
    \uses{def:tomita-pre-operators}
    \leanok
    For all $x \in \text{Dom}(S_0)$ and $y \in \text{Dom}(F_0)$,
    \begin{align}
        \left< y, S_0 x \right> = \left< x, F_0 y \right>.
    \end{align}
\end{lemma}
\begin{proof}
    \leanok
    Write $x = a\Omega$ with $a \in \mathcal{M}$ and $y = b\Omega$ with $b \in \mathcal{M}'$. Since $ab = ba$,
    \begin{align}
        \left< b\Omega, a^*\Omega \right> = \left< ab\Omega, \Omega \right> = \left< ba\Omega, \Omega \right> = \left< a\Omega, b^*\Omega \right>.
    \end{align}
\end{proof}

\subsubsection{Linearization and Closure}

\begin{definition}[The Linearized Pre-Operator $S_0'$]
    \label{def:linearized-tomita-pre-operator}
    \lean{Physicslib4.Modular.linearizedTomitaPre}
    \uses{def:tomita-pre-operators, def:reference-conjugation, def:hall-3.1}
    \leanok
    The \emph{linearized Tomita pre-operator} is $S_0' = C \circ S_0$, with $\text{Dom}(S_0') = \mathcal{M}\Omega$ and $S_0'(a\Omega) = C(a^*\Omega)$. As a composite of two antilinear maps it is linear, so it is an unbounded operator in the sense of Definition~\ref{def:hall-3.1}; it is densely defined by Lemma~\ref{lmm:tomita-pre-operators-dense-involutive}.
\end{definition}

\begin{lemma}[Values of the Pre-Operators]
    \label{lmm:tomita-pre-operators-apply}
    \lean{Physicslib4.Modular.tomitaPre_apply, Physicslib4.Modular.linearizedTomitaPre_domain, Physicslib4.Modular.linearizedTomitaPre_apply, Physicslib4.Modular.span_eq_of_isVonNeumannSet}
    \uses{def:tomita-pre-operators, def:linearized-tomita-pre-operator, def:von-neumann-algebra-commutant-convention}
    \leanok
    Let $\mathcal{M}$ be a von Neumann algebra and $\Omega$ cyclic and separating. Then $S_0(a\Omega) = a^*\Omega$ for every $a \in \mathcal{M}$ and $F_0(b\Omega) = b^*\Omega$ for every $b \in \mathcal{M}'$; $\text{Dom}(S_0') = \text{Dom}(S_0) = \mathcal{M}\Omega$; and $S_0' x = C(S_0 x)$ for every $x \in \text{Dom}(S_0')$.
\end{lemma}
\begin{proof}
    \leanok
    \uses{lmm:tomita-pre-operators-well-defined, lmm:commutant-von-neumann-algebra}
    $\mathcal{M}$ and $\mathcal{M}'$ are linear subspaces, so their linear spans are themselves and the domains are $\mathcal{M}\Omega$ and $\mathcal{M}'\Omega$. For $a \in \mathcal{M}$ the value $S_0(a\Omega)$ is $c^*\Omega$ for some $c \in \mathcal{M}$ with $c\Omega = a\Omega$, which is $a^*\Omega$ by Lemma~\ref{lmm:tomita-pre-operators-well-defined}; likewise for $F_0$. The graph of $S_0'$ is $\{(a\Omega, C(a^*\Omega)) : a \in \mathcal{M}\}$, which is the graph of $C \circ S_0$.
\end{proof}

\begin{lemma}[The Adjoint of $S_0'$ Extends $F_0 C$]
    \label{lmm:linearized-adjoint-extends}
    \lean{Physicslib4.Modular.linearizedTomitaPre_adjoint_extends}
    \uses{def:linearized-tomita-pre-operator, def:hall-9.1, def:hall-9.3}
    \leanok
    $C(\mathcal{M}'\Omega) \subset \text{Dom}\big( (S_0')^* \big)$ and $(S_0')^* C y = F_0 y$ for every $y \in \mathcal{M}'\Omega$; that is, $(S_0')^*$ is an extension of $F_0 \circ C$. In particular $(S_0')^*$ is densely defined.
\end{lemma}
\begin{proof}
    \leanok
    \uses{lmm:tomita-adjoint-relation, lmm:reference-conjugation-properties, lmm:characterizing-adjoint-domain-membership, lmm:tomita-pre-operators-dense-involutive}
    For $x \in \mathcal{M}\Omega$, Lemma~\ref{lmm:reference-conjugation-properties} and Lemma~\ref{lmm:tomita-adjoint-relation} give
    \begin{align}
        \left< Cy, S_0'x \right> = \left< Cy, C S_0 x \right> = \overline{\left< y, S_0x \right>} = \overline{\left< x, F_0y \right>} = \left< F_0 y, x \right>,
    \end{align}
    which is the defining relation of the adjoint (Lemma~\ref{lmm:characterizing-adjoint-domain-membership}). The set $C(\mathcal{M}'\Omega)$ is dense because $C$ is an isometric bijection and $\mathcal{M}'\Omega$ is dense.
\end{proof}

\begin{lemma}[A Densely Defined Adjoint Gives Closability]
    \label{lmm:closable-of-dense-adjoint}
    \lean{Physicslib4.Modular.isClosable_of_dense_adjoint}
    \uses{def:hall-9.1, def:hall-9.6}
    \leanok
    Let $A$ be a densely defined operator on $\mathbf{H}$ whose adjoint $A^*$ is densely defined. Then $A$ is closable.
\end{lemma}
\begin{proof}
    \leanok
    \uses{prpstn:closure-linearity-and-sequential-description, lmm:hall-dense-testing, prpstn:continuity-of-norm-and-inner-product}
    By Proposition~\ref{prpstn:closure-linearity-and-sequential-description} it suffices that $(0, \eta) \in \overline{\Gamma(A)}$ implies $\eta = 0$. Take $\psi_n \in \text{Dom}(A)$ with $\psi_n \to 0$ and $A\psi_n \to \eta$. For $\phi \in \text{Dom}(A^*)$, $\left< \phi, \eta \right> = \lim_n \left< \phi, A\psi_n \right> = \lim_n \left< A^*\phi, \psi_n \right> = 0$, so $\eta = 0$ by Lemma~\ref{lmm:hall-dense-testing}, $\text{Dom}(A^*)$ being dense.
\end{proof}

\begin{corollary}[$S_0'$ is Closable]
    \label{crllr:linearized-tomita-closable}
    \lean{Physicslib4.Modular.linearizedTomitaPre_isClosable}
    \uses{def:linearized-tomita-pre-operator, def:hall-9.6}
    \leanok
    The linearized pre-operator $S_0'$ is closable.
\end{corollary}
\begin{proof}
    \leanok
    \uses{lmm:linearized-adjoint-extends, lmm:closable-of-dense-adjoint}
    $S_0'$ is densely defined, and $(S_0')^*$ is densely defined by Lemma~\ref{lmm:linearized-adjoint-extends}; apply Lemma~\ref{lmm:closable-of-dense-adjoint}.
\end{proof}

\begin{definition}[The Tomita Operator]
    \label{def:tomita-operator}
    \lean{Physicslib4.Modular.tomitaOperator}
    \uses{crllr:linearized-tomita-closable, def:reference-conjugation, def:hall-9.6}
    \leanok
    Let $S' = (S_0')^{\text{cl}}$ be the closure of $S_0'$ (\ref{def:hall-9.6}). The \emph{Tomita operator} is the antilinear operator
    \begin{align}
        S = C \circ S', \qquad \text{Dom}(S) = \text{Dom}(S').
    \end{align}
\end{definition}

\begin{lemma}[The Graph of $S$ is the Closure of the Graph of $S_0$]
    \label{lmm:tomita-operator-independent-of-c}
    \lean{Physicslib4.Modular.tomitaOperator_graph_eq_closure}
    \uses{def:tomita-operator, def:tomita-pre-operators}
    \leanok
    $\Gamma(S)$ is the closure of $\Gamma(S_0)$ in $\mathbf{H} \times \mathbf{H}$. In particular $S$ and $\text{Dom}(S)$ do not depend on the choice of the reference conjugation $C$.
\end{lemma}
\begin{proof}
    \leanok
    \uses{lmm:reference-conjugation-properties, prpstn:closure-linearity-and-sequential-description, lmm:componentwise-convergence}
    The map $\Phi(\phi, \psi) = (\phi, C\psi)$ is an isometric real-linear bijection of $\mathbf{H} \times \mathbf{H}$ (Lemma~\ref{lmm:reference-conjugation-properties} and Lemma~\ref{lmm:componentwise-convergence}), hence a homeomorphism, so it maps closures to closures. Since $C(C\psi) = \psi$, $\Phi$ maps $\Gamma(S_0')$ onto $\Gamma(S_0)$ and $\Gamma(S') = \overline{\Gamma(S_0')}$ (Proposition~\ref{prpstn:closure-linearity-and-sequential-description}) onto $\Gamma(S)$. The closure of $\Gamma(S_0)$ involves only $\mathcal{M}$ and $\Omega$.
\end{proof}

\begin{lemma}[The Graph of $S$ is Symmetric under the Flip]
    \label{lmm:tomita-graph-flip}
    \lean{Physicslib4.Modular.tomitaOperator_graph_swap}
    \uses{def:tomita-operator}
    \leanok
    If $(\phi, \psi) \in \Gamma(S)$, then $(\psi, \phi) \in \Gamma(S)$.
\end{lemma}
\begin{proof}
    \leanok
    \uses{lmm:tomita-operator-independent-of-c, lmm:tomita-pre-operators-dense-involutive}
    The flip $(\phi, \psi) \mapsto (\psi, \phi)$ is an isometric bijection of $\mathbf{H} \times \mathbf{H}$. It maps $\Gamma(S_0)$ onto itself by Lemma~\ref{lmm:tomita-pre-operators-dense-involutive}, hence maps its closure, which is $\Gamma(S)$ by Lemma~\ref{lmm:tomita-operator-independent-of-c}, onto itself.
\end{proof}

\begin{theorem}[The Tomita Operator is an Involution]
    \label{thrm:tomita-operator-involution}
    \lean{Physicslib4.Modular.tomitaOperator_involution}
    \uses{def:tomita-operator}
    \leanok
    $S$ is injective, $S$ maps $\text{Dom}(S)$ onto $\text{Dom}(S)$, and $S S \phi = \phi$ for all $\phi \in \text{Dom}(S)$ (so $S^{-1} = S$, Lemma~\ref{lmm:tomita-inverse-is-self}). The domain and the range of $S$ are dense in $\mathbf{H}$.
\end{theorem}
\begin{proof}
    \leanok
    \uses{lmm:tomita-graph-flip, lmm:tomita-pre-operators-dense-involutive}
    For $\phi \in \text{Dom}(S)$, Lemma~\ref{lmm:tomita-graph-flip} applied to $(\phi, S\phi)$ gives $S\phi \in \text{Dom}(S)$ and $S S\phi = \phi$. This makes $S$ a bijection of $\text{Dom}(S)$ onto itself, so its range is its domain, which contains the dense set $\mathcal{M}\Omega$ (Lemma~\ref{lmm:tomita-pre-operators-dense-involutive}).
\end{proof}

\begin{corollary}[Properties of $S'$]
    \label{crllr:linearized-tomita-properties}
    \lean{Physicslib4.Modular.linearizedTomita_closure_properties}
    \uses{def:tomita-operator, def:hall-9.6, def:kernel-of-an-unbounded-operator, def:range-of-an-unbounded-operator}
    \leanok
    $S'$ is closed, densely defined and injective, and its range is dense in $\mathbf{H}$.
\end{corollary}
\begin{proof}
    \leanok
    \uses{thrm:tomita-operator-involution, lmm:reference-conjugation-properties, prpstn:closure-linearity-and-sequential-description}
    $S'$ is a closure, hence closed, and its domain contains $\mathcal{M}\Omega$. Since $S' = CS$ with $C$ an isometric bijection, $S'\phi = 0$ forces $S\phi = 0$, so $\phi = 0$ by Theorem~\ref{thrm:tomita-operator-involution}, and $\text{Range}(S') = C(\text{Range}(S))$ is dense.
\end{proof}

\subsection{The Modular Operator and the Modular Conjugation}\label{sbsctn:modular-operator-conjugation}

The modular operator is $\Delta = (S')^*S'$ and the modular conjugation comes from the polar decomposition of $S'$. The subsection first supplies the general results on unbounded operators that this needs: von Neumann's theorem that $T^*T$ is self-adjoint, a composition law and a core property for the functional calculus of Definition~\ref{def:hall-10.5}, the absolute value $|T|$, and the polar decomposition of a closed operator.

\subsubsection{Products of Operators and von Neumann's Theorem}

\begin{definition}[Product of Operators]
    \label{def:product-of-unbounded-operators}
    \lean{Physicslib4.Spectral.Unbounded.pmapMul}
    \uses{def:hall-3.1, def:antiunitary-operator, def:hall-9.6}
    \leanok
    Let $A$ and $B$ be operators on $\mathbf{H}$, each linear or antilinear. Their \emph{product} $AB$ is the operator with
    \begin{align}
        \text{Dom}(AB) = \{ \psi \in \text{Dom}(B) \mid B\psi \in \text{Dom}(A) \}, \qquad (AB)\psi = A(B\psi).
    \end{align}
    Equivalently, $(\psi, \chi) \in \Gamma(AB)$ if and only if there is $\eta$ with $(\psi, \eta) \in \Gamma(B)$ and $(\eta, \chi) \in \Gamma(A)$. An everywhere-defined map, such as a unitary or antiunitary operator, is regarded as an operator with domain $\mathbf{H}$; thus $\text{Dom}(WA) = \text{Dom}(A)$ and $\text{Dom}(AW) = W^{-1}(\text{Dom}(A))$ for a bijection $W$. Products, and the inverses of Definition~\ref{def:inverse-of-operator}, are partial maps that need not be densely defined; the density requirement of Definition~\ref{def:hall-3.1} is not imposed on them. Equality of operators includes equality of domains.
\end{definition}

\begin{definition}[Core]
    \label{def:core}
    \lean{LinearPMap.HasCore}
    \uses{def:hall-9.6}
    \leanok
    Let $A$ be a closed operator on $\mathbf{H}$. A subspace $D \subset \text{Dom}(A)$ is a \emph{core} for $A$ if $\Gamma(A\vert_D)$ is dense in $\Gamma(A)$: for every $\psi \in \text{Dom}(A)$ there are $\psi_n \in D$ with $\psi_n \to \psi$ and $A\psi_n \to A\psi$.
\end{definition}

For the results from here to Lemma~\ref{lmm:core-of-adjoint-product}, $T$ is a closed, densely defined operator on $\mathbf{H}$ and $D = T^*T$.

\begin{definition}[The Operator $T^*T$]
    \label{def:adjoint-product-operator}
    \lean{Physicslib4.Spectral.Unbounded.adjointMulSelf}
    \uses{def:product-of-unbounded-operators, def:hall-9.1}
    \leanok
    Let $T$ be a densely defined operator on $\mathbf{H}$. The operator $T^*T$ is the product (\ref{def:product-of-unbounded-operators}) of $T^*$ with $T$:
    \begin{align}
        \text{Dom}(T^*T) = \{ x \in \text{Dom}(T) \mid Tx \in \text{Dom}(T^*) \}, \qquad (T^*T)x = T^*(Tx).
    \end{align}
    In the Lean this is a definition, the composite of \texttt{T\ensuremath{^\dagger}} with the restriction of $T$ to $\text{Dom}(T^*T)$ (\texttt{LinearPMap.comp} with \texttt{LinearPMap.domRestrict}), not a hypothesis characterising its graph; its graph is $\{(x, z) \mid \exists w,\ (x, w) \in \Gamma(T),\ (w, z) \in \Gamma(T^*)\}$.
\end{definition}

\begin{lemma}[Solvability of $x + T^*Tx = y$]
    \label{lmm:von-neumann-solvability}
    \lean{Physicslib4.Spectral.Unbounded.exists_add_adjointMulSelf_eq}
    \uses{def:adjoint-product-operator, def:hall-9.1, def:hall-9.6}
    \leanok
    For every $y \in \mathbf{H}$ there is $x \in \text{Dom}(T^*T)$ with $x + T^*Tx = y$.
\end{lemma}
\begin{proof}
    \leanok
    \uses{def:product-hilbert-space, prpstn:hall-a.49, lmm:characterizing-adjoint-domain-membership}
    $\Gamma(T)$ is a closed subspace of $\mathbf{H} \times \mathbf{H}$ (\ref{def:product-hilbert-space}), so by Proposition~\ref{prpstn:hall-a.49}, $(y, 0) = (x, Tx) + (u, v)$ with $x \in \text{Dom}(T)$ and $(u, v) \perp \Gamma(T)$. Orthogonality says $\left< v, Tx' \right> = \left< -u, x' \right>$ for all $x' \in \text{Dom}(T)$, so $v \in \text{Dom}(T^*)$ with $T^*v = -u$ (Lemma~\ref{lmm:characterizing-adjoint-domain-membership}). The components give $Tx = -v \in \text{Dom}(T^*)$, hence $T^*Tx = u$, and $y = x + u = x + T^*Tx$.
\end{proof}

\begin{lemma}[The Form of $T^*T$]
    \label{lmm:von-neumann-form}
    \lean{Physicslib4.Spectral.Unbounded.inner_adjointMulSelf}
    \uses{def:adjoint-product-operator, def:hall-9.1}
    \leanok
    For $x \in \text{Dom}(T^*T)$ and $x' \in \text{Dom}(T)$, $\left< x', T^*Tx \right> = \left< Tx', Tx \right>$. In particular $\left< x, T^*Tx \right> = \|Tx\|^2 \ge 0$.
\end{lemma}
\begin{proof}
    \leanok
    \uses{lmm:characterizing-adjoint-domain-membership}
    The defining relation of the adjoint for $\phi = Tx \in \text{Dom}(T^*)$ and $\psi = x'$ reads $\left< Tx, Tx' \right> = \left< T^*Tx, x' \right>$; take complex conjugates.
\end{proof}

\begin{lemma}[$T^*T$ is Densely Defined]
    \label{lmm:von-neumann-dense-domain}
    \lean{Physicslib4.Spectral.Unbounded.hasDenseDomain_adjointMulSelf}
    \uses{def:adjoint-product-operator}
    \leanok
    $\text{Dom}(T^*T)$ is dense in $\mathbf{H}$.
\end{lemma}
\begin{proof}
    \leanok
    \uses{lmm:von-neumann-solvability, lmm:von-neumann-form, crllr:trivial-complement-characterizes-density}
    Let $y \perp \text{Dom}(T^*T)$ and write $y = x + T^*Tx$ (Lemma~\ref{lmm:von-neumann-solvability}). Then $0 = \left< x, y \right> = \|x\|^2 + \|Tx\|^2$ by Lemma~\ref{lmm:von-neumann-form}, so $x = 0$ and $y = 0$; conclude by Corollary~\ref{crllr:trivial-complement-characterizes-density}.
\end{proof}

\begin{lemma}[Self-Adjointness from a Surjective Shift]
    \label{lmm:self-adjoint-of-range-eq-top}
    \lean{Physicslib4.Spectral.Unbounded.adjoint_eq_of_range_eq_top}
    \uses{def:hall-9.1, def:hall-9.3, def:hall-9.5, def:range-of-an-unbounded-operator, def:kernel-of-an-unbounded-operator}
    \leanok
    Let $A$ be a densely defined operator such that $A^*$ is an extension of $A$, and let $\mu \in \mathbb{C}$ satisfy $\text{Range}(A - \mu\mathbf{1}) = \mathbf{H}$ and $\text{Ker}(A^* - \mu\mathbf{1}) = \{0\}$. Then $A^* = A$; that is, $A$ is self-adjoint. In Lean this is \texttt{adjoint\_eq\_of\_range\_eq\_top}.
\end{lemma}
\begin{proof}
    \leanok
    \uses{def:hall-9.3}
    Let $\phi \in \text{Dom}(A^*)$ and pick $\chi \in \text{Dom}(A)$ with $(A - \mu\mathbf{1})\chi = (A^* - \mu\mathbf{1})\phi$. Since $A^*$ extends $A$, $(A^* - \mu\mathbf{1})(\phi - \chi) = 0$, so $\phi = \chi \in \text{Dom}(A)$. Hence $\text{Dom}(A^*) = \text{Dom}(A)$, and $A^* = A$ there.
\end{proof}

\begin{lemma}[$T^*T$ is Symmetric and Positive]
    \label{lmm:von-neumann-symmetric-positive}
    \lean{Physicslib4.Spectral.Unbounded.adjointMulSelf_symmetric_positive}
    \uses{def:adjoint-product-operator, def:hall-9.2, def:hall-9.3, def:positive-unbounded-operator}
    \leanok
    $D$ is symmetric, $D^*$ is an extension of $D$, and $D$ is positive.
\end{lemma}
\begin{proof}
    \leanok
    \uses{lmm:von-neumann-form, lmm:von-neumann-dense-domain, prpstn:hall-9.4}
    For $x, x' \in \text{Dom}(D)$, Lemma~\ref{lmm:von-neumann-form} gives $\left< x', Dx \right> = \left< Tx', Tx \right> = \overline{\left< x, Dx' \right>} = \left< Dx', x \right>$, and $\left< x, Dx \right> = \|Tx\|^2 \ge 0$. $D$ is densely defined (Lemma~\ref{lmm:von-neumann-dense-domain}), so Proposition~\ref{prpstn:hall-9.4} gives $D \subset D^*$.
\end{proof}

\begin{lemma}[$D^* + \mathbf{1}$ is Injective]
    \label{lmm:von-neumann-kernel}
    \lean{Physicslib4.Spectral.Unbounded.ker_subSmul_adjoint_adjointMulSelf}
    \uses{def:adjoint-product-operator, def:kernel-of-an-unbounded-operator}
    \leanok
    $\text{Ker}(D^* + \mathbf{1}) = \{0\}$.
\end{lemma}
\begin{proof}
    \leanok
    \uses{lmm:von-neumann-solvability, lmm:von-neumann-dense-domain, prpstn:hall-9.12, prpstn:hall-9.13, lmm:adjoint-of-scalar-multiple-of-identity}
    By Propositions~\ref{prpstn:hall-9.12} and~\ref{prpstn:hall-9.13} with Lemma~\ref{lmm:adjoint-of-scalar-multiple-of-identity}, $\text{Ker}(D^* + \mathbf{1}) = \text{Range}(D + \mathbf{1})^\perp$, and $\text{Range}(D + \mathbf{1}) = \mathbf{H}$ by Lemma~\ref{lmm:von-neumann-solvability}.
\end{proof}

\begin{theorem}[von Neumann's Theorem]
    \label{thrm:von-neumann-adjoint-product}
    \lean{Physicslib4.Spectral.Unbounded.isSelfAdjoint_adjointMulSelf}
    \uses{def:adjoint-product-operator, def:hall-9.5, def:positive-unbounded-operator}
    \leanok
    $T^*T$ is a positive, self-adjoint, densely defined operator.
\end{theorem}
\begin{proof}
    \leanok
    \uses{lmm:self-adjoint-of-range-eq-top, lmm:von-neumann-symmetric-positive, lmm:von-neumann-kernel, lmm:von-neumann-solvability, lmm:von-neumann-dense-domain}
    Apply Lemma~\ref{lmm:self-adjoint-of-range-eq-top} to $D$ with $\mu = -1$: $D$ is densely defined (Lemma~\ref{lmm:von-neumann-dense-domain}), $D \subset D^*$ (Lemma~\ref{lmm:von-neumann-symmetric-positive}), $\text{Range}(D + \mathbf{1}) = \mathbf{H}$ (Lemma~\ref{lmm:von-neumann-solvability}) and $\text{Ker}(D^* + \mathbf{1}) = \{0\}$ (Lemma~\ref{lmm:von-neumann-kernel}). Positivity is Lemma~\ref{lmm:von-neumann-symmetric-positive}.
\end{proof}

\begin{lemma}[$\text{Dom}(T^*T)$ is a Core for $T$]
    \label{lmm:core-of-adjoint-product}
    \lean{Physicslib4.Spectral.Unbounded.hasCore_domain_adjointMulSelf}
    \uses{def:core, def:adjoint-product-operator}
    \leanok
    $\text{Dom}(T^*T)$ is a core for $T$.
\end{lemma}
\begin{proof}
    \leanok
    \uses{lmm:von-neumann-solvability, lmm:von-neumann-form, crllr:trivial-complement-characterizes-density}
    $\Gamma(T)$ is closed in $\mathbf{H} \times \mathbf{H}$, hence complete (Mathlib's \texttt{IsClosed.completeSpace\_coe}), so it is a Hilbert space. If $(y, Ty) \in \Gamma(T)$ is orthogonal to $(x, Tx)$ for every $x \in \text{Dom}(T^*T)$, then by Lemma~\ref{lmm:von-neumann-form}, $0 = \left< y, x \right> + \left< Ty, Tx \right> = \left< y, x + T^*Tx \right>$, and $x + T^*Tx$ ranges over all of $\mathbf{H}$ (Lemma~\ref{lmm:von-neumann-solvability}); so $y = 0$. Corollary~\ref{crllr:trivial-complement-characterizes-density} in $\Gamma(T)$ gives density.
\end{proof}

\subsubsection{The Composition Law of the Functional Calculus}

The results of this paragraph concern a projection-valued measure $\mu$ on a measurable space $(X, \Sigma)$ (the $\sigma$-algebra is written $\Sigma$ rather than $\Omega(X)$ as in Section~\ref{sctn:unbounded-spectral-theorems}, because $\Omega$ is the cyclic vector) and the operators $\int_X f \, d\mu$ of Proposition~\ref{prpstn:hall-10.1}, with domains $W_f$ (Proposition~\ref{prpstn:hall-10.2}). For a self-adjoint $A$ and $\mu = \mu^A$ they are statements about the functional calculus $f(A) = \int f \, d\mu^A$ of Definition~\ref{def:hall-10.5}. Write $h(\mu) = \int_X h \, d\mu$ for short.

\begin{lemma}[Bounded Approximants are Close to the Truncations]
    \label{lmm:fc-difference-to-truncation}
    \lean{Physicslib4.Spectral.Unbounded.tendsto_norm_integral_sub_truncation}
    \uses{prpstn:hall-10.1, prpstn:hall-10.2}
    \leanok
    Let $h$ be measurable, $\psi \in W_h$, and let $h_n$ be bounded measurable functions with $h_n \to h$ pointwise and $|h_n| \le |h|$. With the truncations $t_n = h \cdot 1_{\{|h| < n\}}$, $\|h_n(\mu)\psi - t_n(\mu)\psi\| \to 0$.
\end{lemma}
\begin{proof}
    \leanok
    \uses{lmm:norm-identity-bounded-integral, thrm:dominated-convergence-theorem, thrm:operator-valued-integration}
    By linearity of the bounded integral and Lemma~\ref{lmm:norm-identity-bounded-integral}, $\|h_n(\mu)\psi - t_n(\mu)\psi\|^2 = \int |h_n - t_n|^2 \, d\mu_\psi$. The integrand tends to $0$ pointwise and is dominated by $4|h|^2$, which is $\mu_\psi$-integrable because $\psi \in W_h$; apply Theorem~\ref{thrm:dominated-convergence-theorem}.
\end{proof}

\begin{lemma}[Dominated Convergence for the Spectral Integral]
    \label{lmm:fc-dominated-convergence}
    \lean{Physicslib4.Spectral.Unbounded.tendsto_integral_apply_of_dominated}
    \uses{prpstn:hall-10.1, prpstn:hall-10.2}
    \leanok
    Let $h$ be measurable, $\psi \in W_h$, and let $h_n$ be bounded measurable functions with $h_n \to h$ pointwise and $|h_n| \le |h|$. Then $h_n(\mu)\psi \to h(\mu)\psi$.
\end{lemma}
\begin{proof}
    \leanok
    \uses{lmm:fc-difference-to-truncation, lmm:truncations-converge}
    $t_n(\mu)\psi \to h(\mu)\psi$ by Lemma~\ref{lmm:truncations-converge}, and $h_n(\mu)\psi - t_n(\mu)\psi \to 0$ by Lemma~\ref{lmm:fc-difference-to-truncation}.
\end{proof}

\begin{lemma}[Composition with a Bounded Function on the Left]
    \label{lmm:fc-bounded-left-composition}
    \lean{Physicslib4.Spectral.Unbounded.integral_apply_pmapIntegral}
    \uses{prpstn:hall-10.1, prpstn:hall-10.2}
    \leanok
    Let $u$ be bounded measurable, $h$ measurable and $\psi \in W_h$. Then $\psi \in W_{uh}$ and $u(\mu)\,h(\mu)\psi = (uh)(\mu)\psi$.
\end{lemma}
\begin{proof}
    \leanok
    \uses{lmm:fc-dominated-convergence, lmm:truncations-converge, prpstn:integral-multiplicative}
    $|uh|^2 \le \sup|u|^2 |h|^2$ gives $\psi \in W_{uh}$. With $t_n$ the truncations of $h$, $u(\mu) t_n(\mu)\psi = (u t_n)(\mu)\psi$ by multiplicativity of the bounded integral (Proposition~\ref{prpstn:integral-multiplicative}). The left side tends to $u(\mu)h(\mu)\psi$ by Lemma~\ref{lmm:truncations-converge} and continuity of $u(\mu)$, the right side to $(uh)(\mu)\psi$ by Lemma~\ref{lmm:fc-dominated-convergence}.
\end{proof}

\begin{lemma}[The Associated Measure of $h(\mu)\psi$]
    \label{lmm:fc-associated-measure}
    \lean{Physicslib4.Spectral.Unbounded.assoc_pmapIntegral}
    \uses{prpstn:hall-10.1, prpstn:hall-10.2}
    \leanok
    Let $h$ be measurable, not necessarily bounded, and $\psi \in W_h$. Then for every $E \in \Sigma$,
    \begin{align}
        \mu_{h(\mu)\psi}(E) = \int_E |h|^2 \, d\mu_\psi;
    \end{align}
    that is, $\mu_{h(\mu)\psi}$ is the measure with density $|h|^2$ with respect to $\mu_\psi$ (in Mathlib, $\mu_\psi$\texttt{.withDensity} $|h|^2$). This extends Lemma~\ref{lmm:associated-measure-of-image} from bounded to unbounded $h$.
\end{lemma}
\begin{proof}
    \leanok
    \uses{lmm:fc-bounded-left-composition, lmm:associated-measure-of-image, prpstn:hall-10.2}
    $\mu_{h(\mu)\psi}(E) = \|\mu(E)h(\mu)\psi\|^2 = \|(1_E h)(\mu)\psi\|^2$ by Lemma~\ref{lmm:fc-bounded-left-composition} with $u = 1_E$, and the norm identity of Proposition~\ref{prpstn:hall-10.2} gives $\int |1_E h|^2 \, d\mu_\psi$.
\end{proof}

\begin{lemma}[The Domain of a Composite]
    \label{lmm:fc-composition-domain}
    \lean{Physicslib4.Spectral.Unbounded.pmapIntegral_mem_integralDomain_iff}
    \uses{prpstn:hall-10.2}
    \leanok
    Let $f, g$ be measurable and $\psi \in W_g$. Then $g(\mu)\psi \in W_f$ if and only if $\psi \in W_{fg}$.
\end{lemma}
\begin{proof}
    \leanok
    \uses{lmm:fc-associated-measure}
    By Lemma~\ref{lmm:fc-associated-measure}, $\mu_{g(\mu)\psi}$ has density $|g|^2$ with respect to $\mu_\psi$, so integration against a measure with a density (Mathlib's \texttt{lintegral\_withDensity\_eq\_lintegral\_mul}) gives $\int |f|^2 \, d\mu_{g(\mu)\psi} = \int |f|^2 |g|^2 \, d\mu_\psi = \int |fg|^2 \, d\mu_\psi$.
\end{proof}

\begin{theorem}[Composition Law of the Functional Calculus]
    \label{thrm:fc-composition}
    \lean{Physicslib4.Spectral.Unbounded.pmapIntegral_comp}
    \uses{def:product-of-unbounded-operators, prpstn:hall-10.1, def:hall-9.3}
    \leanok
    For measurable $f, g$, the operator $(fg)(\mu)$ is an extension of $f(\mu)\,g(\mu)$, and
    \begin{align}
        \text{Dom}\big( f(\mu)\,g(\mu) \big) = W_g \cap W_{fg}.
    \end{align}
    In particular $f(A)\,g(A) \subset (fg)(A)$ for a self-adjoint $A$.
\end{theorem}
\begin{proof}
    \leanok
    \uses{lmm:fc-composition-domain, lmm:fc-bounded-left-composition, lmm:fc-dominated-convergence, lmm:truncations-converge}
    The domain identity is Lemma~\ref{lmm:fc-composition-domain}. Let $\psi$ be in it and $E_n = \{|f| < n\}$. Applying Lemma~\ref{lmm:fc-bounded-left-composition} twice, $\mu(E_n)\,(fg)(\mu)\psi = (1_{E_n} f g)(\mu)\psi = (1_{E_n} f)(\mu)\,g(\mu)\psi$. As $n \to \infty$ the left side tends to $(fg)(\mu)\psi$ (Lemma~\ref{lmm:fc-dominated-convergence} with $h = 1$) and the right side to $f(\mu)g(\mu)\psi$ (Lemma~\ref{lmm:truncations-converge}).
\end{proof}

\begin{lemma}[Spectral Projections Commute with $h(\mu)$]
    \label{lmm:fc-commutes-with-projections}
    \lean{Physicslib4.Spectral.Unbounded.pmapIntegral_apply_projection}
    \uses{prpstn:hall-10.1, prpstn:hall-10.2}
    \leanok
    Let $h$ be measurable, $\psi \in W_h$ and $E \in \Sigma$. Then $\mu(E)\psi \in W_h$ and $h(\mu)\mu(E)\psi = \mu(E)h(\mu)\psi$.
\end{lemma}
\begin{proof}
    \leanok
    \uses{thrm:fc-composition, lmm:fc-bounded-left-composition}
    $\mu(E) = 1_E(\mu)$ is defined on all of $\mathbf{H}$, and $\psi \in W_{h 1_E}$ by Lemma~\ref{lmm:fc-bounded-left-composition}. Theorem~\ref{thrm:fc-composition} with $f = h$, $g = 1_E$ gives $\mu(E)\psi \in W_h$ and $h(\mu)\mu(E)\psi = (h1_E)(\mu)\psi = (1_E h)(\mu)\psi$, which is $\mu(E)h(\mu)\psi$ by Lemma~\ref{lmm:fc-bounded-left-composition}.
\end{proof}

\begin{lemma}[A Real Function of $A$ is Closed]
    \label{lmm:fc-real-closed}
    \lean{Physicslib4.Spectral.Unbounded.isSelfAdjoint_functionalCalculus_of_real}
    \uses{def:hall-10.5, def:hall-9.6}
    \leanok
    Let $A$ be self-adjoint and $f : \mathbb{R} \to \mathbb{R}$ measurable. Then $f(A)$ is self-adjoint, hence closed.
\end{lemma}
\begin{proof}
    \leanok
    \uses{prpstn:hall-10.3, prpstn:hall-9.8}
    $f(A)$ is self-adjoint by Proposition~\ref{prpstn:hall-10.3}, so it equals its adjoint, whose graph is closed (Proposition~\ref{prpstn:hall-9.8}).
\end{proof}

\begin{lemma}[$\text{Dom}(A)$ Gives a Core for Every $f(A)$]
    \label{lmm:fc-core}
    \lean{Physicslib4.Spectral.Unbounded.hasCore_functionalCalculus}
    \uses{def:core, def:hall-10.5, def:spectral-measure, lmm:fc-real-closed}
    \leanok
    Let $A$ be self-adjoint and $f : \mathbb{R} \to \mathbb{R}$ measurable, so that $f(A)$ is closed (Lemma~\ref{lmm:fc-real-closed}). Then $\text{Dom}(A) \cap \text{Dom}(f(A))$ is a core for $f(A)$. More precisely, for $\psi \in \text{Dom}(f(A))$ and $E_n = \{ \lambda \mid |\lambda| < n \}$, the vectors $\mu^A(E_n)\psi$ lie in $\text{Dom}(A) \cap \text{Dom}(f(A))$, and $\mu^A(E_n)\psi \to \psi$, $f(A)\mu^A(E_n)\psi \to f(A)\psi$.
\end{lemma}
\begin{proof}
    \leanok
    \uses{lmm:fc-commutes-with-projections, lmm:bounded-on-set-range-in-domain, lmm:fc-dominated-convergence}
    $\mu^A(E_n)\psi \in W_\iota = \text{Dom}(A)$ by Lemma~\ref{lmm:bounded-on-set-range-in-domain}, and $\mu^A(E_n)\psi \in W_f$ with $f(A)\mu^A(E_n)\psi = \mu^A(E_n)f(A)\psi$ by Lemma~\ref{lmm:fc-commutes-with-projections}. Lemma~\ref{lmm:fc-dominated-convergence} with $h = 1$ and $h_n = 1_{E_n}$ gives $\mu^A(E_n)\chi \to \chi$ for every $\chi$, applied to $\chi = \psi$ and $\chi = f(A)\psi$.
\end{proof}

\begin{lemma}[The Pushforward is a Projection-Valued Measure]
    \label{lmm:pushforward-pvm-is-pvm}
    \lean{Physicslib4.Spectral.Unbounded.assoc_pushforwardPVM}
    \uses{def:projection-valued-measure}
    \leanok
    Let $\varphi : X \to Y$ be measurable and set $\varphi_*\mu(F) = \mu(\varphi^{-1}(F))$. Then $\varphi_*\mu$ is a projection-valued measure on $Y$, and $(\varphi_*\mu)_\psi = \varphi_*(\mu_\psi)$ for every $\psi \in \mathbf{H}$.
\end{lemma}
\begin{proof}
    \leanok
    Preimages preserve complements, intersections and countable disjoint unions, so the axioms of Definition~\ref{def:projection-valued-measure} transfer; $(\varphi_*\mu)_\psi(F) = \left< \psi, \mu(\varphi^{-1}F)\psi \right> = \mu_\psi(\varphi^{-1}F)$.
\end{proof}

\begin{lemma}[Domains and Forms under Pushforward]
    \label{lmm:pushforward-pvm-domain}
    \lean{Physicslib4.Spectral.Unbounded.integralDomain_pushforwardPVM}
    \uses{lmm:pushforward-pvm-is-pvm, prpstn:hall-10.2}
    \leanok
    For measurable $g : Y \to \mathbb{C}$, $W^{\varphi_*\mu}_g = W^\mu_{g \circ \varphi}$, and $\int_Y g \, d(\varphi_*\mu)_\psi = \int_X (g \circ \varphi) \, d\mu_\psi$ for $\psi$ in this domain.
\end{lemma}
\begin{proof}
    \leanok
    \uses{lmm:pushforward-pvm-is-pvm}
    By Lemma~\ref{lmm:pushforward-pvm-is-pvm} and the pushforward integration formula (Mathlib's \texttt{lintegral\_map} for $\int |g|^2$, which gives the domains, and \texttt{integral\_map} for the complex integral $\int g$, which is defined because $g \in L^2$ of the finite measure $\mu_\psi$).
\end{proof}

\begin{lemma}[Spectral Integrals under Pushforward]
    \label{lmm:pushforward-pvm}
    \lean{Physicslib4.Spectral.Unbounded.pmapIntegral_pushforwardPVM}
    \uses{lmm:pushforward-pvm-is-pvm, prpstn:hall-10.1}
    \leanok
    For measurable $g : Y \to \mathbb{C}$, $\int_Y g \, d(\varphi_*\mu) = \int_X (g \circ \varphi) \, d\mu$, including equality of domains.
\end{lemma}
\begin{proof}
    \leanok
    \uses{lmm:pushforward-pvm-domain, prpstn:hall-10.1, prpstn:hall-10.2}
    The domains agree and the quadratic forms $\int g \, d(\varphi_*\mu)_\psi$ and $\int (g \circ \varphi) \, d\mu_\psi$ agree (Lemma~\ref{lmm:pushforward-pvm-domain}); polarization gives equal sesquilinear forms, and the uniqueness clause of Proposition~\ref{prpstn:hall-10.1} identifies the operators.
\end{proof}

\begin{corollary}[A Function of a Function of $A$]
    \label{crllr:fc-of-fc}
    \lean{Physicslib4.Spectral.Unbounded.functionalCalculus_functionalCalculus}
    \uses{def:spectral-measure, def:hall-10.5}
    \leanok
    Let $A$ be self-adjoint and $f : \mathbb{R} \to \mathbb{R}$ measurable. Then $f(A)$ is self-adjoint with spectral measure $\mu^{f(A)} = f_*\mu^A$, and $g(f(A)) = (g \circ f)(A)$ for every measurable $g : \mathbb{R} \to \mathbb{C}$.
\end{corollary}
\begin{proof}
    \leanok
    \uses{lmm:pushforward-pvm, prpstn:hall-10.3, thrm:hall-10.4}
    $f(A)$ is self-adjoint by Proposition~\ref{prpstn:hall-10.3}. Lemma~\ref{lmm:pushforward-pvm} with $g(\lambda) = \lambda$ gives $\int \lambda \, d(f_*\mu^A) = f(A)$, so $\mu^{f(A)} = f_*\mu^A$ by the uniqueness in Theorem~\ref{thrm:hall-10.4}; Lemma~\ref{lmm:pushforward-pvm} for general $g$ gives the last claim.
\end{proof}

\subsubsection{The Absolute Value of a Closed Operator}

Write $s(\lambda) = \sqrt{\max(\lambda, 0)}$, a real, non-negative, measurable function on $\mathbb{R}$. For a positive self-adjoint $A$, $A^{1/2} \equiv s(A)$.

\begin{definition}[Absolute Value]
    \label{def:absolute-value}
    \lean{Physicslib4.Spectral.Unbounded.absoluteValue}
    \uses{def:adjoint-product-operator, thrm:von-neumann-adjoint-product, def:hall-10.5}
    \leanok
    For a closed, densely defined operator $T$ on $\mathbf{H}$, the \emph{absolute value} of $T$ is $|T| \equiv (T^*T)^{1/2} = s(T^*T)$, the functional calculus of the self-adjoint operator $T^*T$ (Theorem~\ref{thrm:von-neumann-adjoint-product}).
\end{definition}

For the next four results, $T$ is a closed, densely defined operator on $\mathbf{H}$ and $D = T^*T$.

\begin{lemma}[$\text{Dom}(T^*T) \subset \text{Dom}(|T|)$]
    \label{lmm:adjoint-product-domain-in-sqrt-domain}
    \lean{Physicslib4.Spectral.Unbounded.adjointMulSelf_domain_le_absoluteValue}
    \uses{def:absolute-value}
    \leanok
    Every $x \in \text{Dom}(D)$ lies in $\text{Dom}(|T|)$.
\end{lemma}
\begin{proof}
    \leanok
    \uses{lmm:positive-spectral-support, thrm:von-neumann-adjoint-product, lmm:l2-implies-l1, thrm:hall-10.4, def:spectral-measure}
    By Lemma~\ref{lmm:positive-spectral-support}, $\mu^D_x$ is concentrated on $[0, \infty)$, where $s^2(\lambda) = \lambda$. Since $x \in \text{Dom}(D) = W_\iota$ (Theorem~\ref{thrm:hall-10.4}), $\int \lambda^2 \, d\mu^D_x < \infty$, and Lemma~\ref{lmm:l2-implies-l1} gives $\int s^2 \, d\mu^D_x = \int \lambda \, d\mu^D_x < \infty$; so $x \in W_s$.
\end{proof}

\begin{lemma}[The Norm Identity on $\text{Dom}(T^*T)$]
    \label{lmm:absolute-value-norm-on-core}
    \lean{Physicslib4.Spectral.Unbounded.norm_absoluteValue_of_mem_adjointMulSelf}
    \uses{def:absolute-value, lmm:adjoint-product-domain-in-sqrt-domain}
    \leanok
    For $x \in \text{Dom}(D)$, $\||T|x\| = \|Tx\|$.
\end{lemma}
\begin{proof}
    \leanok
    \uses{lmm:adjoint-product-domain-in-sqrt-domain, lmm:positive-spectral-support, prpstn:hall-10.2, thrm:hall-10.4, lmm:von-neumann-form}
    By Proposition~\ref{prpstn:hall-10.2}, Lemma~\ref{lmm:positive-spectral-support} and $D = \int \lambda \, d\mu^D$ (Theorem~\ref{thrm:hall-10.4}), $\||T|x\|^2 = \int s^2 \, d\mu^D_x = \int \lambda \, d\mu^D_x = \left< x, Dx \right>$, which is $\|Tx\|^2$ by Lemma~\ref{lmm:von-neumann-form}.
\end{proof}

\begin{lemma}[$|T|$ is Closed and Positive, with Core $\text{Dom}(T^*T)$]
    \label{lmm:absolute-value-closed}
    \lean{Physicslib4.Spectral.Unbounded.absoluteValue_closed}
    \uses{def:absolute-value, def:core, def:positive-unbounded-operator}
    \leanok
    $|T|$ is a positive self-adjoint operator, hence closed, and $\text{Dom}(D)$ is a core for $|T|$.
\end{lemma}
\begin{proof}
    \leanok
    \uses{prpstn:hall-10.3, prpstn:hall-10.2, prpstn:hall-9.8, lmm:fc-core, lmm:adjoint-product-domain-in-sqrt-domain}
    $|T|$ is self-adjoint by Proposition~\ref{prpstn:hall-10.3}, so it equals its adjoint, whose graph is closed (Proposition~\ref{prpstn:hall-9.8}); it is positive since $\left< \psi, |T|\psi \right> = \int s \, d\mu^D_\psi \ge 0$ (Proposition~\ref{prpstn:hall-10.2}). By Lemma~\ref{lmm:fc-core}, $\text{Dom}(D) \cap \text{Dom}(|T|)$ is a core, and it equals $\text{Dom}(D)$ by Lemma~\ref{lmm:adjoint-product-domain-in-sqrt-domain}.
\end{proof}

\begin{lemma}[Transfer of Domains through a Core, One Direction]
    \label{lmm:closure-transfer-one-direction}
    \lean{Physicslib4.Spectral.Unbounded.domain_le_of_hasCore_norm_eq}
    \uses{def:core, def:hall-9.6}
    \leanok
    Let $A$ and $B$ be closed operators on $\mathbf{H}$ and $D$ a core for $A$ with $D \subset \text{Dom}(B)$ and $\|A\psi\| = \|B\psi\|$ for $\psi \in D$. Then $\text{Dom}(A) \subset \text{Dom}(B)$ and $\|A\psi\| = \|B\psi\|$ for $\psi \in \text{Dom}(A)$.
\end{lemma}
\begin{proof}
    \leanok
    \uses{prpstn:closure-linearity-and-sequential-description, prpstn:continuity-of-norm-and-inner-product}
    Let $\psi \in \text{Dom}(A)$ and take $\psi_n \in D$ with $\psi_n \to \psi$, $A\psi_n \to A\psi$. Since $\|B\psi_n - B\psi_m\| = \|A\psi_n - A\psi_m\|$, $(B\psi_n)$ is Cauchy and converges; $B$ is closed, so $\psi \in \text{Dom}(B)$ and $B\psi_n \to B\psi$, and the norms pass to the limit.
\end{proof}

\begin{lemma}[Transfer of Domains through a Common Core]
    \label{lmm:closure-transfer-common-core}
    \lean{Physicslib4.Spectral.Unbounded.domain_eq_of_hasCore_norm_eq}
    \uses{def:core, def:hall-9.6}
    \leanok
    Let $A$ and $B$ be closed operators on $\mathbf{H}$ and $D$ a subspace that is a core for both, with $\|A\psi\| = \|B\psi\|$ for $\psi \in D$. Then $\text{Dom}(A) = \text{Dom}(B)$ and $\|A\psi\| = \|B\psi\|$ for all $\psi$ in it.
\end{lemma}
\begin{proof}
    \leanok
    \uses{lmm:closure-transfer-one-direction}
    Apply Lemma~\ref{lmm:closure-transfer-one-direction} to $(A, B)$ and to $(B, A)$.
\end{proof}

\begin{theorem}[Domain and Norm of $|T|$]
    \label{thrm:absolute-value-domain}
    \lean{Physicslib4.Spectral.Unbounded.absoluteValue_domain}
    \uses{def:absolute-value}
    \leanok
    $|T|$ is positive and self-adjoint, $\text{Dom}(|T|) = \text{Dom}(T)$, and $\||T|\psi\| = \|T\psi\|$ for every $\psi \in \text{Dom}(T)$.
\end{theorem}
\begin{proof}
    \leanok
    \uses{lmm:closure-transfer-common-core, lmm:core-of-adjoint-product, lmm:absolute-value-closed, lmm:absolute-value-norm-on-core}
    Apply Lemma~\ref{lmm:closure-transfer-common-core} to $T$ and $|T|$ with the common core $\text{Dom}(D)$ (Lemmas~\ref{lmm:core-of-adjoint-product} and~\ref{lmm:absolute-value-closed}), on which the norms agree by Lemma~\ref{lmm:absolute-value-norm-on-core}.
\end{proof}

\begin{lemma}[Uniqueness of the Positive Square Root]
    \label{lmm:positive-square-root-unique}
    \lean{Physicslib4.Spectral.Unbounded.sqrt_pmapMul_self}
    \uses{def:product-of-unbounded-operators, def:positive-unbounded-operator, def:hall-10.5}
    \leanok
    Let $P$ be a positive self-adjoint operator. Then $PP$ is self-adjoint and $(PP)^{1/2} = P$.
\end{lemma}
\begin{proof}
    \leanok
    \uses{thrm:fc-composition, crllr:fc-of-fc, lmm:positive-spectral-support, lmm:integral-ignores-null-sets, thrm:hall-10.4}
    With $q(\lambda) = \lambda^2$, Theorem~\ref{thrm:fc-composition} gives $PP \subset q(P)$ with $\text{Dom}(PP) = W_\iota \cap W_q = W_q$ (as $\lambda^2 \le 1 + \lambda^4$ and $\mu^P_\psi$ is finite), so $PP = q(P)$, which is self-adjoint. By Corollary~\ref{crllr:fc-of-fc}, $s(q(P)) = (s \circ q)(P)$, and $s \circ q = |\lambda|$ agrees with $\lambda$ off the $\mu^P$-null set $(-\infty, 0)$ (Lemma~\ref{lmm:positive-spectral-support}); so $(s \circ q)(P) = P$ by Lemma~\ref{lmm:integral-ignores-null-sets} and Theorem~\ref{thrm:hall-10.4}.
\end{proof}

\subsubsection{Polar Decomposition}

For the next three results, $T$ is a closed, densely defined, injective operator on $\mathbf{H}$ with dense range.

\begin{lemma}[$|T|$ is Injective with Dense Range]
    \label{lmm:absolute-value-injective-dense-range}
    \lean{Physicslib4.Spectral.Unbounded.absoluteValue_injective_denseRange}
    \uses{def:absolute-value, def:kernel-of-an-unbounded-operator, def:range-of-an-unbounded-operator}
    \leanok
    $|T|$ is injective and $\text{Range}(|T|)$ is dense in $\mathbf{H}$.
\end{lemma}
\begin{proof}
    \leanok
    \uses{thrm:absolute-value-domain, prpstn:hall-9.12, crllr:trivial-complement-characterizes-density}
    If $|T|\psi = 0$, then $\|T\psi\| = 0$ by Theorem~\ref{thrm:absolute-value-domain}, so $\psi = 0$. Since $|T|$ is self-adjoint, $\text{Range}(|T|)^\perp = \text{Ker}(|T|^*) = \text{Ker}(|T|) = \{0\}$ (Proposition~\ref{prpstn:hall-9.12}).
\end{proof}

\begin{lemma}[The Isometry on the Range of $|T|$]
    \label{lmm:polar-isometry-on-range}
    \lean{Physicslib4.Spectral.Unbounded.exists_polarIsometry_on_range}
    \uses{def:absolute-value, def:range-of-an-unbounded-operator}
    \leanok
    The assignment $V_0(|T|\psi) = T\psi$, for $\psi \in \text{Dom}(T)$, is a well-defined, linear and isometric map from $\text{Range}(|T|)$ to $\mathbf{H}$.
\end{lemma}
\begin{proof}
    \leanok
    \uses{thrm:absolute-value-domain}
    $\text{Dom}(|T|) = \text{Dom}(T)$ and $\||T|\psi\| = \|T\psi\|$ (Theorem~\ref{thrm:absolute-value-domain}). If $|T|\psi = |T|\psi'$, then $\|T\psi - T\psi'\| = \||T|(\psi - \psi')\| = 0$, so $V_0$ is well defined; linearity and isometry follow from those of $T$ and $|T|$ and the norm identity.
\end{proof}

\begin{lemma}[The Isometry of the Polar Decomposition]
    \label{lmm:polar-isometry}
    \lean{Physicslib4.Spectral.Unbounded.existsUnique_polarIsometry}
    \uses{def:absolute-value, def:product-of-unbounded-operators}
    \leanok
    There is a unique $V \in \mathcal{B}(\mathbf{H})$ with $V|T|\psi = T\psi$ for all $\psi \in \text{Dom}(T)$. It is isometric, and $T = V|T|$.
\end{lemma}
\begin{proof}
    \leanok
    \uses{lmm:polar-isometry-on-range, lmm:absolute-value-injective-dense-range, thrm:bounded-linear-transformation-theorem, thrm:absolute-value-domain}
    $V_0$ of Lemma~\ref{lmm:polar-isometry-on-range} has dense domain (Lemma~\ref{lmm:absolute-value-injective-dense-range}), so it has a unique bounded extension $V$ (Theorem~\ref{thrm:bounded-linear-transformation-theorem}), isometric by continuity of the norm. The equality $\text{Dom}(|T|) = \text{Dom}(T)$ (Theorem~\ref{thrm:absolute-value-domain}) gives $T = V|T|$.
\end{proof}

\begin{theorem}[Polar Decomposition]
    \label{thrm:polar-decomposition}
    \lean{Physicslib4.Spectral.Unbounded.polar_decomposition}
    \uses{def:absolute-value, def:unitary-operator, def:product-of-unbounded-operators}
    \leanok
    The isometry $V$ of Lemma~\ref{lmm:polar-isometry} is unitary, and $T = V|T|$ with $|T|$ positive, self-adjoint and injective.
\end{theorem}
\begin{proof}
    \leanok
    \uses{lmm:polar-isometry, lmm:absolute-value-injective-dense-range, thrm:absolute-value-domain, prpstn:polarization-identity}
    The range of an isometry on a complete space is complete, hence closed; it contains the dense set $\text{Range}(T)$, so $V$ is onto. An isometry preserves inner products by Proposition~\ref{prpstn:polarization-identity}, so $V$ is unitary. The other claims are Theorem~\ref{thrm:absolute-value-domain} and Lemma~\ref{lmm:absolute-value-injective-dense-range}.
\end{proof}

\begin{lemma}[Adjoint of a Unitary Times an Operator]
    \label{lmm:adjoint-unitary-product}
    \lean{Physicslib4.Spectral.Unbounded.adjoint_unitary_compPMap}
    \uses{def:unitary-operator, def:product-of-unbounded-operators, def:hall-9.1}
    \leanok
    Let $U$ be unitary and $P$ densely defined. Then $(UP)^* = P^*U^{-1}$.
\end{lemma}
\begin{proof}
    \leanok
    \uses{lmm:characterizing-adjoint-domain-membership}
    $\left< \phi, UP\psi \right> = \left< U^{-1}\phi, P\psi \right>$ for $\psi \in \text{Dom}(P) = \text{Dom}(UP)$, so by Lemma~\ref{lmm:characterizing-adjoint-domain-membership}, $\phi \in \text{Dom}((UP)^*)$ if and only if $U^{-1}\phi \in \text{Dom}(P^*)$, and then $(UP)^*\phi = P^*U^{-1}\phi$.
\end{proof}

\begin{theorem}[Uniqueness of the Polar Decomposition]
    \label{thrm:polar-decomposition-unique}
    \lean{Physicslib4.Spectral.Unbounded.polar_decomposition_unique}
    \uses{thrm:polar-decomposition, def:positive-unbounded-operator}
    \leanok
    Let $T$ be closed, densely defined and injective with dense range, and suppose $T = UP$ with $U$ unitary and $P$ positive and self-adjoint. Then $P = |T|$ and $U = V$.
\end{theorem}
\begin{proof}
    \leanok
    \uses{lmm:adjoint-unitary-product, lmm:positive-square-root-unique, lmm:absolute-value-injective-dense-range, lmm:polar-isometry}
    By Lemma~\ref{lmm:adjoint-unitary-product}, $T^*T = PU^{-1}UP = PP$, so $|T| = (PP)^{1/2} = P$ by Lemma~\ref{lmm:positive-square-root-unique}. Then $U$ and $V$ agree on the dense set $\text{Range}(|T|)$ (Lemmas~\ref{lmm:polar-isometry} and~\ref{lmm:absolute-value-injective-dense-range}), and both are continuous.
\end{proof}

\subsubsection{Inverses and Antiunitary Conjugation}

\begin{definition}[Inverse of an Injective Operator]
    \label{def:inverse-of-operator}
    \lean{LinearPMap.inverse}
    \uses{def:hall-9.6, def:antiunitary-operator, def:kernel-of-an-unbounded-operator, def:range-of-an-unbounded-operator}
    \leanok
    Let $A$ be an injective operator on $\mathbf{H}$, linear or antilinear. Its \emph{inverse} $A^{-1}$ is the operator whose graph is the flip of $\Gamma(A)$,
    \begin{align}
        \Gamma(A^{-1}) = \{ (A\psi, \psi) \mid \psi \in \text{Dom}(A) \}, \qquad \text{Dom}(A^{-1}) = \text{Range}(A).
    \end{align}
    It is a partial map of the same kind as $A$ and need not be densely defined. For linear $A$ this is Mathlib's \texttt{LinearPMap.inverse}.
\end{definition}

\begin{lemma}[Inverse of a Product with a Bijection]
    \label{lmm:inverse-of-product-with-bijection}
    \lean{Physicslib4.Spectral.Unbounded.inverse_unitary_compPMap}
    \uses{def:inverse-of-operator, def:product-of-unbounded-operators, def:unitary-operator, def:antiunitary-operator}
    \leanok
    Let $A$ be injective and $W$ unitary or antiunitary. Then $WA$ is injective and $(WA)^{-1} = A^{-1}W^{-1}$.
\end{lemma}
\begin{proof}
    \leanok
    \uses{def:inverse-of-operator, def:product-of-unbounded-operators}
    $WA$ is injective because $W$ is. Both sides have domain $\text{Range}(WA) = W(\text{Range}(A)) = \{ \phi \mid W^{-1}\phi \in \text{Range}(A) \}$, and for $\phi = WA\psi$ both send $\phi$ to $\psi$.
\end{proof}

\begin{lemma}[An Injective Self-Adjoint Operator has no Atom at $0$]
    \label{lmm:injective-no-atom-at-zero}
    \lean{Physicslib4.Spectral.Unbounded.spectralMeasure_singleton_zero}
    \uses{def:spectral-measure, def:kernel-of-an-unbounded-operator}
    \leanok
    If $A$ is self-adjoint and injective, then $\mu^A(\{0\}) = 0$.
\end{lemma}
\begin{proof}
    \leanok
    \uses{lmm:bounded-on-set-range-in-domain, lmm:range-membership-concentrates-measure, prpstn:hall-10.2}
    Let $\psi \in \text{Range}(\mu^A(\{0\}))$. Then $\psi \in \text{Dom}(A)$ (Lemma~\ref{lmm:bounded-on-set-range-in-domain}), $\mu^A_\psi$ is concentrated on $\{0\}$ (Lemma~\ref{lmm:range-membership-concentrates-measure}), and $\|A\psi\|^2 = \int \lambda^2 \, d\mu^A_\psi = 0$ (Proposition~\ref{prpstn:hall-10.2}); so $\psi = 0$.
\end{proof}

\begin{lemma}[$\iota^-(A)$ Inverts $A$ on $\text{Dom}(A)$]
    \label{lmm:inverse-fc-extends}
    \lean{Physicslib4.Spectral.Unbounded.functionalCalculus_inv_apply_self}
    \uses{def:hall-10.5, def:spectral-measure}
    \leanok
    Let $A$ be self-adjoint and injective, and let $\iota^{-}(\lambda) = 1/\lambda$ for $\lambda \ne 0$ and $\iota^{-}(0) = 0$. For $\psi \in \text{Dom}(A)$, $A\psi \in \text{Dom}(\iota^{-}(A))$ and $\iota^{-}(A)A\psi = \psi$.
\end{lemma}
\begin{proof}
    \leanok
    \uses{lmm:injective-no-atom-at-zero, thrm:fc-composition, lmm:integral-ignores-null-sets, lmm:functional-calculus-bounded, thrm:hall-10.4}
    $\iota^{-}\iota = 1$ off the $\mu^A$-null set $\{0\}$ (Lemma~\ref{lmm:injective-no-atom-at-zero}), so $(\iota^{-}\iota)(A) = \mathbf{1}$ (Lemmas~\ref{lmm:integral-ignores-null-sets} and~\ref{lmm:functional-calculus-bounded}). Theorem~\ref{thrm:fc-composition} with $f = \iota^{-}$, $g = \iota$ gives $\text{Dom}(\iota^{-}(A)A) = W_\iota \cap W_{\iota^-\iota} = \text{Dom}(A)$ and $\iota^{-}(A)A\psi = \psi$.
\end{proof}

\begin{lemma}[$\text{Dom}(\iota^-(A)) \subset \text{Range}(A)$]
    \label{lmm:inverse-fc-domain}
    \lean{Physicslib4.Spectral.Unbounded.self_apply_functionalCalculus_inv}
    \uses{lmm:inverse-fc-extends, def:range-of-an-unbounded-operator}
    \leanok
    Under the hypotheses of Lemma~\ref{lmm:inverse-fc-extends}, every $\phi \in \text{Dom}(\iota^{-}(A))$ satisfies $\iota^-(A)\phi \in \text{Dom}(A)$ and $A\iota^{-}(A)\phi = \phi$; in particular $\text{Dom}(\iota^{-}(A)) \subset \text{Range}(A)$.
\end{lemma}
\begin{proof}
    \leanok
    \uses{lmm:injective-no-atom-at-zero, thrm:fc-composition, lmm:integral-ignores-null-sets, lmm:functional-calculus-bounded, thrm:hall-10.4}
    As in Lemma~\ref{lmm:inverse-fc-extends}, $(\iota\iota^{-})(A) = \mathbf{1}$, and Theorem~\ref{thrm:fc-composition} with $f = \iota$, $g = \iota^{-}$ gives $\text{Dom}(A\,\iota^{-}(A)) = W_{\iota^-} \cap W_{\iota\iota^-} = W_{\iota^-}$ and $A\iota^{-}(A)\phi = \phi$.
\end{proof}

\begin{lemma}[The Inverse by the Functional Calculus]
    \label{lmm:inverse-functional-calculus}
    \lean{Physicslib4.Spectral.Unbounded.inverse_eq_functionalCalculus_inv}
    \uses{def:hall-10.5, def:inverse-of-operator, lmm:inverse-fc-extends}
    \leanok
    Let $A$ be self-adjoint and injective. Then $A^{-1} = \iota^{-}(A)$.
\end{lemma}
\begin{proof}
    \leanok
    \uses{lmm:inverse-fc-extends, lmm:inverse-fc-domain, def:inverse-of-operator}
    By Lemma~\ref{lmm:inverse-fc-extends}, $\iota^{-}(A)$ is an extension of $A^{-1}$ (Definition~\ref{def:inverse-of-operator}), and by Lemma~\ref{lmm:inverse-fc-domain} its domain is contained in $\text{Range}(A) = \text{Dom}(A^{-1})$.
\end{proof}

\begin{lemma}[The Square Root is Positive and Self-Adjoint]
    \label{lmm:square-root-positive}
    \lean{Physicslib4.Spectral.Unbounded.sqrt_functionalCalculus_positive}
    \uses{def:hall-10.5, def:positive-unbounded-operator}
    \leanok
    Let $A$ be self-adjoint. Then $A^{1/2} = s(A)$ is positive and self-adjoint.
\end{lemma}
\begin{proof}
    \leanok
    \uses{prpstn:hall-10.2, prpstn:hall-10.3}
    $s \ge 0$ is real and measurable: self-adjointness is Proposition~\ref{prpstn:hall-10.3}, and $\left< \psi, A^{1/2}\psi \right> = \int s \, d\mu^A_\psi \ge 0$ by Proposition~\ref{prpstn:hall-10.2}.
\end{proof}

\begin{lemma}[The Square Root of an Injective Positive Operator is Injective]
    \label{lmm:square-root-injective}
    \lean{Physicslib4.Spectral.Unbounded.ker_sqrt_functionalCalculus}
    \uses{def:hall-10.5, def:positive-unbounded-operator, def:kernel-of-an-unbounded-operator}
    \leanok
    Let $A$ be positive, self-adjoint and injective. Then $A^{1/2}$ is injective.
\end{lemma}
\begin{proof}
    \leanok
    \uses{prpstn:hall-10.2, lmm:positive-spectral-support, lmm:injective-no-atom-at-zero, lmm:associated-measure-total-mass}
    Let $A^{1/2}\psi = 0$. By Proposition~\ref{prpstn:hall-10.2}, $0 = \|A^{1/2}\psi\|^2 = \int \max(\lambda, 0) \, d\mu^A_\psi$, and a nonnegative integral vanishes only if the integrand vanishes almost everywhere (Mathlib's \texttt{lintegral\_eq\_zero\_iff}); so $\mu^A_\psi((0, \infty)) = 0$. Also $\mu^A((-\infty, 0]) = 0$ (Lemmas~\ref{lmm:positive-spectral-support} and~\ref{lmm:injective-no-atom-at-zero}). Hence $\|\psi\|^2 = \mu^A_\psi(\mathbb{R}) = 0$ (Lemma~\ref{lmm:associated-measure-total-mass}).
\end{proof}

\begin{lemma}[The Inverse Square Root is Positive]
    \label{lmm:inverse-square-root-positive}
    \lean{Physicslib4.Spectral.Unbounded.inv_sqrt_functionalCalculus_positive}
    \uses{def:hall-10.5, def:positive-unbounded-operator}
    \leanok
    Let $A$ be self-adjoint and $r(\lambda) = \lambda^{-1/2}$ for $\lambda > 0$, $r(\lambda) = 0$ otherwise. Then $A^{-1/2} \equiv r(A)$ is positive and self-adjoint.
\end{lemma}
\begin{proof}
    \leanok
    \uses{prpstn:hall-10.2, prpstn:hall-10.3}
    $r \ge 0$ is real and measurable: self-adjointness is Proposition~\ref{prpstn:hall-10.3}, and $\left< \psi, r(A)\psi \right> = \int r \, d\mu^A_\psi \ge 0$ by Proposition~\ref{prpstn:hall-10.2}.
\end{proof}

\begin{corollary}[The Inverse Square Root]
    \label{crllr:inverse-square-root}
    \lean{Physicslib4.Spectral.Unbounded.inverse_sqrt_functionalCalculus}
    \uses{lmm:square-root-positive, lmm:square-root-injective, lmm:inverse-square-root-positive, def:inverse-of-operator}
    \leanok
    Let $A$ be positive, self-adjoint and injective. Then $(A^{1/2})^{-1} = A^{-1/2}$.
\end{corollary}
\begin{proof}
    \leanok
    \uses{lmm:square-root-injective, lmm:inverse-functional-calculus, crllr:fc-of-fc}
    $A^{1/2}$ is injective (Lemma~\ref{lmm:square-root-injective}), so $(A^{1/2})^{-1} = \iota^{-}(A^{1/2}) = (\iota^{-} \circ s)(A)$ by Lemma~\ref{lmm:inverse-functional-calculus} and Corollary~\ref{crllr:fc-of-fc}, and $\iota^{-} \circ s = r$ everywhere.
\end{proof}

\begin{lemma}[The Square of the Square Root]
    \label{lmm:square-of-square-root}
    \lean{Physicslib4.Spectral.Unbounded.sq_sqrt_functionalCalculus}
    \uses{def:hall-10.5, def:positive-unbounded-operator}
    \leanok
    Let $A$ be positive and self-adjoint and $q(\lambda) = \lambda^2$. Then $q(A^{1/2}) = A$.
\end{lemma}
\begin{proof}
    \leanok
    \uses{crllr:fc-of-fc, lmm:positive-spectral-support, lmm:integral-ignores-null-sets, thrm:hall-10.4}
    By Corollary~\ref{crllr:fc-of-fc}, $q(A^{1/2}) = (q \circ s)(A)$, and $q \circ s = \lambda$ off the $\mu^A$-null set $(-\infty, 0)$ (Lemma~\ref{lmm:positive-spectral-support}); so $(q \circ s)(A) = \int \lambda \, d\mu^A = A$ by Lemma~\ref{lmm:integral-ignores-null-sets} and Theorem~\ref{thrm:hall-10.4}.
\end{proof}

\begin{lemma}[The Square of the Inverse Square Root]
    \label{lmm:square-of-inverse-square-root}
    \lean{Physicslib4.Spectral.Unbounded.sq_inv_sqrt_functionalCalculus}
    \uses{lmm:inverse-square-root-positive, def:inverse-of-operator}
    \leanok
    Let $A$ be positive, self-adjoint and injective, and $q(\lambda) = \lambda^2$. Then $q(A^{-1/2}) = A^{-1}$.
\end{lemma}
\begin{proof}
    \leanok
    \uses{crllr:fc-of-fc, lmm:positive-spectral-support, lmm:injective-no-atom-at-zero, lmm:integral-ignores-null-sets, lmm:inverse-functional-calculus}
    By Corollary~\ref{crllr:fc-of-fc}, $q(A^{-1/2}) = (q \circ r)(A)$, and $q \circ r = \iota^{-}$ off $(-\infty, 0]$, which is $\mu^A$-null (Lemmas~\ref{lmm:positive-spectral-support} and~\ref{lmm:injective-no-atom-at-zero}); so $(q \circ r)(A) = \iota^-(A) = A^{-1}$ by Lemmas~\ref{lmm:integral-ignores-null-sets} and~\ref{lmm:inverse-functional-calculus}.
\end{proof}

\begin{lemma}[Antiunitary Conjugation of a Projection-Valued Measure]
    \label{lmm:antiunitary-conjugate-pvm-is-pvm}
    \lean{Physicslib4.Spectral.Unbounded.exists_antiunitaryConj_pvm}
    \uses{def:antiunitary-operator, def:projection-valued-measure}
    \leanok
    Let $W$ be antiunitary and $\mu$ a projection-valued measure on $X$. Then $\nu(E) = W\mu(E)W^{-1}$ is a projection-valued measure with $\nu_\psi = \mu_{W^{-1}\psi}$.
\end{lemma}
\begin{proof}
    \leanok
    \uses{lmm:antiunitary-isometric}
    Each $\nu(E)$ is linear and bounded (Lemma~\ref{lmm:antiunitary-isometric}), and $\left< \phi, \nu(E)\psi \right> = \left< \mu(E)W^{-1}\psi, W^{-1}\phi \right>$ shows that it is a self-adjoint idempotent and that $\nu_\psi(E) = \mu_{W^{-1}\psi}(E)$; continuity of $W$ transfers countable additivity.
\end{proof}

\begin{lemma}[Domains under Antiunitary Conjugation]
    \label{lmm:antiunitary-conjugate-pvm-domain}
    \lean{Physicslib4.Spectral.Unbounded.mem_integralDomain_antiunitaryConj}
    \uses{lmm:antiunitary-conjugate-pvm-is-pvm, prpstn:hall-10.2}
    \leanok
    With $\nu$ as in Lemma~\ref{lmm:antiunitary-conjugate-pvm-is-pvm}, $W^\nu_f = W(W^\mu_f)$ for every measurable $f$.
\end{lemma}
\begin{proof}
    \leanok
    \uses{lmm:antiunitary-conjugate-pvm-is-pvm}
    $\int |f|^2 \, d\nu_\psi = \int |f|^2 \, d\mu_{W^{-1}\psi}$ by Lemma~\ref{lmm:antiunitary-conjugate-pvm-is-pvm}, so $\psi \in W^\nu_f$ if and only if $W^{-1}\psi \in W^\mu_f$.
\end{proof}

\begin{lemma}[Spectral Integrals under Antiunitary Conjugation]
    \label{lmm:antiunitary-conjugate-pvm}
    \lean{Physicslib4.Spectral.Unbounded.pmapIntegral_antiunitaryConj}
    \uses{lmm:antiunitary-conjugate-pvm-is-pvm, def:product-of-unbounded-operators, prpstn:hall-10.1}
    \leanok
    With $\nu$ as in Lemma~\ref{lmm:antiunitary-conjugate-pvm-is-pvm} and $f$ real and measurable, $\int f \, d\nu = W\big( \int f \, d\mu \big) W^{-1}$.
\end{lemma}
\begin{proof}
    \leanok
    \uses{lmm:antiunitary-conjugate-pvm-domain, prpstn:hall-10.1, prpstn:hall-10.2}
    Both sides are linear with domain $W(W^\mu_f)$ (Lemma~\ref{lmm:antiunitary-conjugate-pvm-domain}). The quadratic form of the right side is $\overline{\int f \, d\mu_{W^{-1}\psi}} = \int f \, d\nu_\psi$, since $f$ is real; polarization and the uniqueness in Proposition~\ref{prpstn:hall-10.1} give equality.
\end{proof}

\begin{lemma}[Spectral Measure of an Antiunitary Conjugate]
    \label{lmm:antiunitary-conjugate-spectral-measure}
    \lean{Physicslib4.Spectral.Unbounded.spectralMeasure_antiunitaryConj}
    \uses{def:antiunitary-operator, def:product-of-unbounded-operators, def:spectral-measure}
    \leanok
    Let $W$ be antiunitary and $A$ self-adjoint. Then $WAW^{-1}$ is self-adjoint with spectral measure $W\mu^A(\cdot)W^{-1}$.
\end{lemma}
\begin{proof}
    \leanok
    \uses{lmm:antiunitary-conjugate-pvm-is-pvm, lmm:antiunitary-conjugate-pvm, prpstn:hall-10.3, thrm:hall-10.4}
    With $\nu$ as in Lemma~\ref{lmm:antiunitary-conjugate-pvm-is-pvm}, Lemma~\ref{lmm:antiunitary-conjugate-pvm} with $f(\lambda) = \lambda$ shows $WAW^{-1} = \int \lambda \, d\nu$, self-adjoint by Proposition~\ref{prpstn:hall-10.3}; the uniqueness in Theorem~\ref{thrm:hall-10.4} gives $\mu^{WAW^{-1}} = \nu$.
\end{proof}

\begin{corollary}[Antiunitary Conjugation and the Functional Calculus]
    \label{crllr:antiunitary-conjugate-functional-calculus}
    \lean{Physicslib4.Spectral.Unbounded.functionalCalculus_antiunitaryConj}
    \uses{lmm:antiunitary-conjugate-spectral-measure, def:hall-10.5}
    \leanok
    Let $W$ be antiunitary and $A$ self-adjoint. Then $Wf(A)W^{-1} = f(WAW^{-1})$ for every real measurable $f$.
\end{corollary}
\begin{proof}
    \leanok
    \uses{lmm:antiunitary-conjugate-spectral-measure, lmm:antiunitary-conjugate-pvm}
    $f(WAW^{-1}) = \int f \, d\nu$ by Lemma~\ref{lmm:antiunitary-conjugate-spectral-measure}, which is $Wf(A)W^{-1}$ by Lemma~\ref{lmm:antiunitary-conjugate-pvm}.
\end{proof}

\begin{lemma}[Antiunitary Conjugation Preserves Positivity and Injectivity]
    \label{lmm:antiunitary-conjugate-positive}
    \lean{Physicslib4.Spectral.Unbounded.antiunitaryConj_positive_injective}
    \uses{def:antiunitary-operator, def:product-of-unbounded-operators, def:positive-unbounded-operator}
    \leanok
    Let $W$ be antiunitary and $A$ an operator. If $A$ is positive, or injective, so is $WAW^{-1}$.
\end{lemma}
\begin{proof}
    \leanok
    \uses{lmm:conjugate-linear-isometry-inner}
    By Lemma~\ref{lmm:conjugate-linear-isometry-inner}, $\left< \psi, WAW^{-1}\psi \right>$ is the complex conjugate of $\left< W^{-1}\psi, AW^{-1}\psi \right> \ge 0$. Injectivity holds because $W$ and $W^{-1}$ are injective.
\end{proof}

\subsubsection{The Modular Operator and the Modular Conjugation}

\begin{definition}[The Modular Operator]
    \label{def:modular-operator}
    \lean{Physicslib4.Modular.modularOperator}
    \uses{def:tomita-operator, def:adjoint-product-operator, def:absolute-value}
    \leanok
    The \emph{modular operator} of $(\mathcal{M}, \Omega)$ is $\Delta = (S')^*S'$ (Definition~\ref{def:adjoint-product-operator}), and $\Delta^{1/2} = s(\Delta) = |S'|$.
\end{definition}

\begin{theorem}[Properties of the Modular Operator]
    \label{thrm:modular-operator-properties}
    \lean{Physicslib4.Modular.modularOperator_properties}
    \uses{def:modular-operator, def:hall-9.5, def:positive-unbounded-operator}
    \leanok
    $\Delta$ is positive, self-adjoint, densely defined and injective. Moreover $\text{Dom}(\Delta^{1/2}) = \text{Dom}(S)$ and $\|\Delta^{1/2}\psi\| = \|S\psi\|$ for $\psi \in \text{Dom}(S)$.
\end{theorem}
\begin{proof}
    \leanok
    \uses{crllr:linearized-tomita-properties, thrm:von-neumann-adjoint-product, thrm:absolute-value-domain, lmm:von-neumann-form, lmm:reference-conjugation-properties}
    $S'$ is closed and densely defined (Corollary~\ref{crllr:linearized-tomita-properties}), so Theorems~\ref{thrm:von-neumann-adjoint-product} and~\ref{thrm:absolute-value-domain} apply, with $\|S'\psi\| = \|CS\psi\| = \|S\psi\|$ and $\text{Dom}(S') = \text{Dom}(S)$. If $\Delta\psi = 0$, then $\|S'\psi\|^2 = \left< \psi, \Delta\psi \right> = 0$ (Lemma~\ref{lmm:von-neumann-form}), so $\psi = 0$ by injectivity of $S'$.
\end{proof}

\begin{lemma}[An Intrinsic Description of $\Delta$]
    \label{lmm:modular-operator-intrinsic}
    \lean{Physicslib4.Modular.mem_modularOperator_iff}
    \uses{def:modular-operator}
    \leanok
    $\psi \in \text{Dom}(\Delta)$ with $\Delta\psi = \chi$ if and only if $\psi \in \text{Dom}(S)$ and $\left< S\phi, S\psi \right> = \left< \chi, \phi \right>$ for every $\phi \in \text{Dom}(S)$.
\end{lemma}
\begin{proof}
    \leanok
    \uses{lmm:characterizing-adjoint-domain-membership, lmm:reference-conjugation-properties}
    By Lemma~\ref{lmm:characterizing-adjoint-domain-membership}, $S'\psi \in \text{Dom}((S')^*)$ with $(S')^*S'\psi = \chi$ if and only if $\left< S'\psi, S'\phi \right> = \left< \chi, \phi \right>$ for all $\phi \in \text{Dom}(S')$, and $\left< S'\psi, S'\phi \right> = \left< CS\psi, CS\phi \right> = \left< S\phi, S\psi \right>$ by Lemma~\ref{lmm:reference-conjugation-properties}.
\end{proof}

\begin{definition}[The Modular Conjugation]
    \label{def:modular-conjugation}
    \lean{Physicslib4.Modular.modularConjugation}
    \uses{def:modular-operator, thrm:polar-decomposition, crllr:linearized-tomita-properties}
    \leanok
    Let $V$ be the unitary of the polar decomposition $S' = V|S'| = V\Delta^{1/2}$ (Theorem~\ref{thrm:polar-decomposition}, applicable by Corollary~\ref{crllr:linearized-tomita-properties}). The \emph{modular conjugation} is $J = CV$.
\end{definition}

\begin{theorem}[$S = J\Delta^{1/2}$]
    \label{thrm:modular-conjugation-polar}
    \lean{Physicslib4.Modular.tomitaOperator_eq_modularConjugation_sqrt}
    \uses{def:modular-conjugation, def:antiunitary-operator, def:product-of-unbounded-operators}
    \leanok
    $J$ is antiunitary and $S = J\Delta^{1/2}$, including $\text{Dom}(S) = \text{Dom}(\Delta^{1/2})$.
\end{theorem}
\begin{proof}
    \leanok
    \uses{lmm:antiunitary-composition, thrm:polar-decomposition, lmm:reference-conjugation-properties}
    $J$ is antiunitary by Lemma~\ref{lmm:antiunitary-composition}. Since $C(C\phi) = \phi$ for all $\phi$, $S = CS' = CV\Delta^{1/2} = J\Delta^{1/2}$ by Theorem~\ref{thrm:polar-decomposition}.
\end{proof}

\begin{lemma}[An Intrinsic Description of $J$]
    \label{lmm:modular-conjugation-intrinsic}
    \lean{Physicslib4.Modular.eq_modularConjugation_of_continuous}
    \uses{def:modular-conjugation}
    \leanok
    $J$ is the only continuous map $\mathbf{H} \to \mathbf{H}$ with $J\Delta^{1/2}\psi = S\psi$ for all $\psi \in \text{Dom}(\Delta^{1/2})$.
\end{lemma}
\begin{proof}
    \leanok
    \uses{thrm:modular-conjugation-polar, lmm:absolute-value-injective-dense-range, crllr:linearized-tomita-properties}
    Theorem~\ref{thrm:modular-conjugation-polar} gives the identity, and two continuous maps with this property agree on the dense set $\text{Range}(\Delta^{1/2}) = \text{Range}(|S'|)$ (Lemma~\ref{lmm:absolute-value-injective-dense-range}).
\end{proof}

\begin{theorem}[$\Delta$ and $J$ do not Depend on the Reference Conjugation]
    \label{thrm:modular-objects-independent-of-c}
    \lean{Physicslib4.Modular.modular_objects_independent_of_conjugation}
    \uses{def:modular-operator, def:modular-conjugation, def:reference-conjugation}
    \leanok
    The operators $\Delta$, $\Delta^{1/2}$ and $J$ are the same for every choice of the reference conjugation $C$.
\end{theorem}
\begin{proof}
    \leanok
    \uses{lmm:tomita-operator-independent-of-c, lmm:modular-operator-intrinsic, lmm:modular-conjugation-intrinsic}
    $S$ and $\text{Dom}(S)$ do not depend on $C$ (Lemma~\ref{lmm:tomita-operator-independent-of-c}), so neither does $\Delta$ by Lemma~\ref{lmm:modular-operator-intrinsic}, nor $\Delta^{1/2} = s(\Delta)$, nor $J$ by Lemma~\ref{lmm:modular-conjugation-intrinsic}.
\end{proof}

\begin{lemma}[$S^{-1} = S$]
    \label{lmm:tomita-inverse-is-self}
    \lean{Physicslib4.Modular.tomitaOperator_inverse_eq_self}
    \uses{def:tomita-operator, def:inverse-of-operator}
    \leanok
    $S$ is injective and $S^{-1} = S$.
\end{lemma}
\begin{proof}
    \leanok
    \uses{thrm:tomita-operator-involution, lmm:tomita-graph-flip}
    $S$ is injective by Theorem~\ref{thrm:tomita-operator-involution}, and the flip of $\Gamma(S)$ is $\Gamma(S)$ by Lemma~\ref{lmm:tomita-graph-flip} (applied in both directions).
\end{proof}

\begin{lemma}[$S$ as $J^{-1}$ Times a Conjugated Operator]
    \label{lmm:tomita-inverse-as-product}
    \lean{Physicslib4.Modular.tomitaOperator_eq_symm_conj_invSqrt}
    \uses{def:modular-conjugation, def:modular-operator, lmm:inverse-square-root-positive, def:product-of-unbounded-operators}
    \leanok
    Let $P = J\Delta^{-1/2}J^{-1}$. Then $S = J^{-1}P$.
\end{lemma}
\begin{proof}
    \leanok
    \uses{lmm:tomita-inverse-is-self, thrm:modular-conjugation-polar, lmm:inverse-of-product-with-bijection, crllr:inverse-square-root, thrm:modular-operator-properties}
    By Lemma~\ref{lmm:tomita-inverse-is-self}, Theorem~\ref{thrm:modular-conjugation-polar} and Lemma~\ref{lmm:inverse-of-product-with-bijection}, $S = S^{-1} = (J\Delta^{1/2})^{-1} = (\Delta^{1/2})^{-1}J^{-1}$, which is $\Delta^{-1/2}J^{-1} = J^{-1}P$ by Corollary~\ref{crllr:inverse-square-root} for the positive, self-adjoint, injective $\Delta$ (Theorem~\ref{thrm:modular-operator-properties}).
\end{proof}

\begin{lemma}[A Second Polar Decomposition of $S'$]
    \label{lmm:tomita-inverse-polar}
    \lean{Physicslib4.Modular.linearizedTomita_closure_eq_unitary_conj_invSqrt}
    \uses{lmm:tomita-inverse-as-product}
    \leanok
    With $P = J\Delta^{-1/2}J^{-1}$, $P$ is positive and self-adjoint, $CJ^{-1}$ is unitary, and $S' = (CJ^{-1})P$.
\end{lemma}
\begin{proof}
    \leanok
    \uses{lmm:tomita-inverse-as-product, lmm:inverse-square-root-positive, lmm:antiunitary-conjugate-spectral-measure, lmm:antiunitary-conjugate-positive, lmm:antiunitary-isometric, lmm:antiunitary-composition, lmm:reference-conjugation-properties}
    $S' = CS = (CJ^{-1})P$ by Lemma~\ref{lmm:tomita-inverse-as-product} and $C(C\phi) = \phi$. $P$ is positive and self-adjoint by Lemmas~\ref{lmm:inverse-square-root-positive}, \ref{lmm:antiunitary-conjugate-spectral-measure} and~\ref{lmm:antiunitary-conjugate-positive}, and $CJ^{-1}$ is unitary by Lemmas~\ref{lmm:antiunitary-isometric} and~\ref{lmm:antiunitary-composition}.
\end{proof}

\begin{theorem}[$J$ is an Involution and $J\Delta^{1/2}J = \Delta^{-1/2}$]
    \label{thrm:modular-conjugation-involution}
    \lean{Physicslib4.Modular.modularConjugation_involution}
    \uses{def:modular-conjugation, def:modular-operator, crllr:inverse-square-root}
    \leanok
    $J(J\phi) = \phi$ for every $\phi \in \mathbf{H}$, and $J\Delta^{1/2}J = \Delta^{-1/2}$.
\end{theorem}
\begin{proof}
    \leanok
    \uses{lmm:tomita-inverse-polar, thrm:polar-decomposition-unique, crllr:linearized-tomita-properties, lmm:reference-conjugation-properties}
    By Lemma~\ref{lmm:tomita-inverse-polar} and Theorem~\ref{thrm:polar-decomposition-unique}, applied to $S'$ (Corollary~\ref{crllr:linearized-tomita-properties}), $CJ^{-1} = V = C^{-1}J = CJ$ and $J\Delta^{-1/2}J^{-1} = |S'| = \Delta^{1/2}$. The first gives $J^{-1} = J$, that is $J(J\phi) = \phi$ for every $\phi$; conjugating the second by $J$ gives $J\Delta^{1/2}J = \Delta^{-1/2}$.
\end{proof}

\begin{corollary}[$J\Delta J = \Delta^{-1}$]
    \label{crllr:modular-conjugation-inverts}
    \lean{Physicslib4.Modular.modularConjugation_conj_modularOperator}
    \uses{def:modular-conjugation, def:modular-operator, def:inverse-of-operator}
    \leanok
    $J\Delta J = \Delta^{-1}$.
\end{corollary}
\begin{proof}
    \leanok
    \uses{thrm:modular-conjugation-involution, crllr:antiunitary-conjugate-functional-calculus, lmm:square-of-square-root, lmm:square-of-inverse-square-root, thrm:modular-operator-properties, thrm:hall-10.4}
    $\Delta = q(\Delta^{1/2})$ and $\Delta^{-1} = q(\Delta^{-1/2})$ by Lemmas~\ref{lmm:square-of-square-root} and~\ref{lmm:square-of-inverse-square-root} (Theorem~\ref{thrm:modular-operator-properties}). With $J = J^{-1}$, Corollary~\ref{crllr:antiunitary-conjugate-functional-calculus} and Theorem~\ref{thrm:modular-conjugation-involution} give $J\Delta J = q(J\Delta^{1/2}J) = q(\Delta^{-1/2}) = \Delta^{-1}$.
\end{proof}

\subsection{Takesaki's Lemma}\label{sbsctn:takesaki-lemma}

Takesaki's lemma states that, for $z \in \mathbb{C} \setminus [0, \infty)$ and $a' \in \mathcal{M}'$, the vector $(z - \Delta)^{-1}a'\Omega$ is of the form $a\Omega$ for a unique $a \in \mathcal{M}$, with $\|a\| \le \|a'\| / \sqrt{2(|z| - \operatorname{Re} z)}$. This is Lemma~2.6 of Sorce's proof of Tomita's theorem (J.~Sorce, \emph{A short proof of Tomita's theorem}, \emph{J. Funct. Anal.} (2024), arXiv:2309.16762), which cites it without proof. The statement and the constant are those of a second paper, J.~Sorce, \emph{An intuitive construction of modular flow}, \emph{JHEP} 12 (2023) 079 (arXiv:2309.16766), Lemma~4.1, whose proof uses an operator affiliated with $\mathcal{M}$ and its spectral projections. The proof below uses only bounded operators. It follows the Hahn--Banach argument of M.~A.~Rieffel and A.~van Daele, \emph{A bounded operator approach to Tomita--Takesaki theory}, \emph{Pacific J. Math.} 69 (1977) 187--221, Section~4, Lemma~4.3, which adapts the proof of Sakai's linear Radon--Nikodym theorem. Two changes keep the argument inside the Hilbert space: the separation takes place in $\mathbf{H}$ rather than in the predual of $\mathcal{M}$, and an approximate polar decomposition $h(\varepsilon + h^*h)^{-1/2}$ from the continuous functional calculus replaces the polar decomposition of $h$. The last step is the commutant extension lemma (Lemma~\ref{lmm:commutant-extension}). The same lemma takes the place of Lemma~2.2 of Sorce's proof of Tomita's theorem, a description of $\text{Dom}(S)$ by affiliated operators, in the proof of his Theorem~3.3.

\subsubsection{The Commutant Extension Lemma}

\begin{lemma}[The Commutant Extension Map is Well Defined]
    \label{lmm:commutant-extension-well-defined}
    \lean{Physicslib4.Modular.commutantExtension_linearMap}
    \uses{def:cyclic-separating-vector, lmm:commutant-von-neumann-algebra}
    \leanok
    Let $\psi \in \mathbf{H}$ and $c \ge 0$ satisfy $\|b'\psi\| \le c\|b'\Omega\|$ for all $b' \in \mathcal{M}'$. Then there is a unique linear map $T_0 : \mathcal{M}'\Omega \to \mathbf{H}$ with $T_0(b'\Omega) = b'\psi$ for every $b' \in \mathcal{M}'$, and $\|T_0\eta\| \le c\|\eta\|$ for every $\eta \in \mathcal{M}'\Omega$.
\end{lemma}
\begin{proof}
    \leanok
    \uses{lmm:commutant-von-neumann-algebra}
    $\mathcal{M}'$ is a linear subspace (Lemma~\ref{lmm:commutant-von-neumann-algebra}), so $\mathcal{M}'\Omega$ is the range of the linear map $b' \mapsto b'\Omega$, and $b' \mapsto b'\psi$ is linear. If $b'\Omega = d'\Omega$, the hypothesis applied to $b' - d' \in \mathcal{M}'$ gives $\|(b' - d')\psi\| \le c\|(b' - d')\Omega\| = 0$, so $b'\psi = d'\psi$: the kernel of the first map lies in the kernel of the second, and $T_0$ is the induced map on the range. The norm bound is the hypothesis.
\end{proof}

\begin{lemma}[The Commutant Extension Map Extends to a Bounded Operator]
    \label{lmm:commutant-extension-bounded}
    \lean{Physicslib4.Modular.commutantExtension_bounded}
    \uses{lmm:commutant-extension-well-defined}
    \leanok
    In the situation of Lemma~\ref{lmm:commutant-extension-well-defined}, there is a unique $T \in \mathcal{B}(\mathbf{H})$ that extends $T_0$, and $\|T\| \le c$.
\end{lemma}
\begin{proof}
    \leanok
    \uses{lmm:commutant-extension-well-defined, def:cyclic-separating-vector}
    $\mathcal{M}'\Omega$ is dense because $\Omega$ is cyclic for $\mathcal{M}'$ (\ref{def:cyclic-separating-vector}). A bounded linear map on a dense subspace extends uniquely to a bounded operator on the completion with the same bound (Mathlib's \texttt{ContinuousLinearMap.extend} and \texttt{ContinuousLinearMap.opNorm\_extend\_le}).
\end{proof}

\begin{lemma}[The Extension Lies in $\mathcal{M}$]
    \label{lmm:commutant-extension-in-algebra}
    \lean{Physicslib4.Modular.commutantExtension_mem}
    \uses{lmm:commutant-extension-bounded, def:von-neumann-algebra-commutant-convention}
    \leanok
    The operator $T$ of Lemma~\ref{lmm:commutant-extension-bounded} commutes with every $d' \in \mathcal{M}'$, and hence $T \in \mathcal{M}'' = \mathcal{M}$.
\end{lemma}
\begin{proof}
    \leanok
    \uses{lmm:commutant-extension-bounded, lmm:commutant-von-neumann-algebra, def:cyclic-separating-vector}
    For $b' \in \mathcal{M}'$, $d'b' \in \mathcal{M}'$ (Lemma~\ref{lmm:commutant-von-neumann-algebra}), so $Td'(b'\Omega) = d'b'\psi = d'T(b'\Omega)$. The bounded operators $Td'$ and $d'T$ agree on the dense set $\mathcal{M}'\Omega$, hence are equal. So $T$ lies in the centralizer of $\mathcal{M}'$, which is $\mathcal{M}$ by \ref{def:von-neumann-algebra-commutant-convention}.
\end{proof}

\begin{lemma}[Commutant Extension Lemma]
    \label{lmm:commutant-extension}
    \lean{Physicslib4.Modular.commutantExtension}
    \uses{def:cyclic-separating-vector, def:von-neumann-algebra-commutant-convention}
    \leanok
    Let $\mathcal{M}$ be a von Neumann algebra on $\mathbf{H}$ and $\Omega$ cyclic and separating for $\mathcal{M}$. If $\psi \in \mathbf{H}$ and $c \ge 0$ satisfy $\|b'\psi\| \le c\|b'\Omega\|$ for all $b' \in \mathcal{M}'$, then there is a unique $a \in \mathcal{M}$ with $a\Omega = \psi$, and $\|a\| \le c$.
\end{lemma}
\begin{proof}
    \leanok
    \uses{lmm:commutant-extension-bounded, lmm:commutant-extension-in-algebra, lmm:cyclic-commutant-separating}
    Take $a = T$ (Lemmas~\ref{lmm:commutant-extension-bounded} and~\ref{lmm:commutant-extension-in-algebra}); since $\mathbf{1} \in \mathcal{M}'$, $a\Omega = T(\mathbf{1}\Omega) = \mathbf{1}\psi = \psi$, and $\|a\| \le c$. If also $\tilde{a} \in \mathcal{M}$ and $\tilde{a}\Omega = \psi$, then $(a - \tilde{a})\Omega = 0$ with $a - \tilde{a} \in \mathcal{M}$, so $a = \tilde{a}$ because $\Omega$ is separating for $\mathcal{M}$ (Lemma~\ref{lmm:cyclic-commutant-separating}).
\end{proof}

\subsubsection{The Resolvent of the Modular Operator}

For $z \in \mathbb{C} \setminus [0, \infty)$ let $d(z) = \inf_{t \ge 0} |z - t|$, and let $r_z : \mathbb{R} \to \mathbb{C}$ be $r_z(\lambda) = (z - \lambda)^{-1}$ for $\lambda \ge 0$ and $r_z(\lambda) = 0$ for $\lambda < 0$. Write $\iota(\lambda) = \lambda$, and write $z - \iota$ for the function $\lambda \mapsto z - \lambda$.

\begin{lemma}[Properties of the Scalar Function $r_z$]
    \label{lmm:resolvent-function-props}
    \lean{Physicslib4.Modular.resolventFunction_properties}
    \leanok
    Let $z \in \mathbb{C} \setminus [0, \infty)$. Then $d(z) > 0$, $r_z$ is measurable, $|r_z(\lambda)| \le 1/d(z)$ for every $\lambda \in \mathbb{R}$, $\overline{r_z(\lambda)} = r_{\overline{z}}(\lambda)$ for every $\lambda \in \mathbb{R}$, and $(z - \lambda)\,r_z(\lambda) = 1_{[0, \infty)}(\lambda)$ for every $\lambda \in \mathbb{R}$.
\end{lemma}
\begin{proof}
    \leanok
    $[0, \infty)$ is closed, nonempty and does not contain $z$, so $d(z) > 0$ (Mathlib's \texttt{IsClosed.notMem\_iff\_infDist\_pos}). $r_z$ is the indicator of the measurable set $[0, \infty)$ times the measurable function $\lambda \mapsto (z - \lambda)^{-1}$ (Mathlib's \texttt{Measurable.inv} and \texttt{Measurable.indicator}). For $\lambda \ge 0$, $|z - \lambda| \ge d(z)$ (Mathlib's \texttt{Metric.infDist\_le\_dist\_of\_mem}), so $|r_z(\lambda)| \le 1/d(z)$; for $\lambda < 0$, $r_z(\lambda) = 0$. Since $\lambda$ is real, $\overline{(z - \lambda)^{-1}} = (\overline{z} - \lambda)^{-1}$ (Mathlib's \texttt{Complex.conj\_ofReal}), and $\overline{z} \notin [0, \infty)$. For $\lambda \ge 0$, $z \ne \lambda$, so $(z - \lambda)(z - \lambda)^{-1} = 1$; for $\lambda < 0$, both sides are $0$.
\end{proof}

\begin{lemma}[The Bounded Operator $r_z(\Delta)$]
    \label{lmm:modular-resolvent-bounded}
    \lean{Physicslib4.Modular.modularResolvent_bounded}
    \uses{def:modular-operator, def:hall-10.5, lmm:resolvent-function-props}
    \leanok
    Let $z \in \mathbb{C} \setminus [0, \infty)$. Then $r_z(\Delta)$ is a bounded operator defined on all of $\mathbf{H}$ with $\|r_z(\Delta)\| \le 1/d(z)$ and $r_z(\Delta)^* = r_{\overline{z}}(\Delta)$.
\end{lemma}
\begin{proof}
    \leanok
    \uses{lmm:resolvent-function-props, thrm:modular-operator-properties, lmm:functional-calculus-bounded, prpstn:integral-norm-bound, prpstn:integral-conjugation}
    $\Delta$ is self-adjoint (Theorem~\ref{thrm:modular-operator-properties}) and $r_z$ is bounded and measurable (Lemma~\ref{lmm:resolvent-function-props}), so Lemma~\ref{lmm:functional-calculus-bounded} gives the bounded operator, and Proposition~\ref{prpstn:integral-norm-bound} with $|r_z| \le 1/d(z)$ gives the norm bound. Since $\overline{r_z} = r_{\overline{z}}$, Proposition~\ref{prpstn:integral-conjugation} gives the adjoint.
\end{proof}

\begin{lemma}[The Function $z - \iota$ of $\Delta$]
    \label{lmm:modular-z-minus-fc}
    \lean{Physicslib4.Modular.modularOperator_functionalCalculus_const_sub}
    \uses{def:modular-operator, def:hall-10.5}
    \leanok
    Let $z \in \mathbb{C}$. Then $W_{z - \iota} = \text{Dom}(\Delta)$, and $(z - \iota)(\Delta)\psi = z\psi - \Delta\psi$ for every $\psi \in \text{Dom}(\Delta)$.
\end{lemma}
\begin{proof}
    \leanok
    \uses{thrm:modular-operator-properties, thrm:hall-10.4, lmm:functional-calculus-sub-bounded, lmm:functional-calculus-bounded, thrm:operator-valued-integration, prpstn:integral-of-indicator, def:projection-valued-measure}
    $\Delta$ is self-adjoint (Theorem~\ref{thrm:modular-operator-properties}). Lemma~\ref{lmm:functional-calculus-sub-bounded}, with the bounded constant function $z$ and $f = \iota$, gives $W_{z - \iota} = W_\iota$ and $(z - \iota)(\Delta)\psi = z(\Delta)\psi - \iota(\Delta)\psi$ for $\psi \in W_\iota$. By Theorem~\ref{thrm:hall-10.4}, $\iota(\Delta) = \int_{\mathbb{R}} \lambda \, d\mu^\Delta(\lambda) = \Delta$ as operators with domains, so $W_\iota = \text{Dom}(\Delta)$ and $\iota(\Delta)\psi = \Delta\psi$. The constant function $z$ is bounded, so $z(\Delta)$ is the bounded integral (Lemma~\ref{lmm:functional-calculus-bounded}), and by linearity of that integral (Theorem~\ref{thrm:operator-valued-integration}) and Proposition~\ref{prpstn:integral-of-indicator}, $z(\Delta) = z\,\mu^\Delta(\mathbb{R}) = z\mathbf{1}$ (Definition~\ref{def:projection-valued-measure}).
\end{proof}

\begin{lemma}[$1_{[0, \infty)}(\Delta) = \mathbf{1}$]
    \label{lmm:modular-indicator-nonneg-one}
    \lean{Physicslib4.Modular.modularOperator_functionalCalculus_indicator_Ici}
    \uses{def:modular-operator, def:hall-10.5}
    \leanok
    $1_{[0, \infty)}(\Delta) = \mathbf{1}$.
\end{lemma}
\begin{proof}
    \leanok
    \uses{thrm:modular-operator-properties, lmm:positive-spectral-support, lmm:integral-ignores-null-sets, prpstn:integral-of-indicator, def:projection-valued-measure}
    $\Delta$ is positive and self-adjoint (Theorem~\ref{thrm:modular-operator-properties}), so $\mu^\Delta((-\infty, 0)) = 0$ (Lemma~\ref{lmm:positive-spectral-support}). The functions $1_{[0, \infty)}$ and $1_{\mathbb{R}}$ agree off $(-\infty, 0)$, so $1_{[0, \infty)}(\Delta) = 1_{\mathbb{R}}(\Delta)$ (Lemma~\ref{lmm:integral-ignores-null-sets}), and $1_{\mathbb{R}}(\Delta) = \mu^\Delta(\mathbb{R}) = \mathbf{1}$ (Proposition~\ref{prpstn:integral-of-indicator} and Definition~\ref{def:projection-valued-measure}).
\end{proof}

\begin{lemma}[$r_z(\Delta)$ is a Right Inverse of $z - \Delta$]
    \label{lmm:modular-resolvent-right-inverse}
    \lean{Physicslib4.Modular.modularResolvent_rightInverse}
    \uses{lmm:modular-resolvent-bounded}
    \leanok
    Let $z \in \mathbb{C} \setminus [0, \infty)$ and $\xi \in \mathbf{H}$. Then $r_z(\Delta)\xi \in \text{Dom}(\Delta)$ and $z\, r_z(\Delta)\xi - \Delta r_z(\Delta)\xi = \xi$.
\end{lemma}
\begin{proof}
    \leanok
    \uses{lmm:modular-resolvent-bounded, lmm:resolvent-function-props, lmm:modular-z-minus-fc, lmm:modular-indicator-nonneg-one, lmm:fc-composition-domain, thrm:fc-composition, lmm:functional-calculus-bounded}
    Let $g = z - \iota$. Then $g\,r_z = 1_{[0, \infty)}$ is bounded and measurable, so $W_{g r_z} = W_{r_z} = \mathbf{H}$ (Lemmas~\ref{lmm:resolvent-function-props} and~\ref{lmm:functional-calculus-bounded}). Lemma~\ref{lmm:fc-composition-domain} gives $r_z(\Delta)\xi \in W_g = \text{Dom}(\Delta)$ (Lemma~\ref{lmm:modular-z-minus-fc}). Theorem~\ref{thrm:fc-composition} and Lemma~\ref{lmm:modular-indicator-nonneg-one} give $g(\Delta)r_z(\Delta)\xi = 1_{[0, \infty)}(\Delta)\xi = \xi$, and $g(\Delta)r_z(\Delta)\xi = z\,r_z(\Delta)\xi - \Delta r_z(\Delta)\xi$ by Lemma~\ref{lmm:modular-z-minus-fc}.
\end{proof}

\begin{lemma}[$r_z(\Delta)$ is a Left Inverse of $z - \Delta$]
    \label{lmm:modular-resolvent-left-inverse}
    \lean{Physicslib4.Modular.modularResolvent_leftInverse}
    \uses{lmm:modular-resolvent-bounded}
    \leanok
    Let $z \in \mathbb{C} \setminus [0, \infty)$ and $\psi \in \text{Dom}(\Delta)$. Then $r_z(\Delta)(z\psi - \Delta\psi) = \psi$.
\end{lemma}
\begin{proof}
    \leanok
    \uses{lmm:resolvent-function-props, lmm:modular-z-minus-fc, lmm:modular-indicator-nonneg-one, lmm:fc-bounded-left-composition}
    With $g = z - \iota$, $\psi \in W_g$ and $z\psi - \Delta\psi = g(\Delta)\psi$ by Lemma~\ref{lmm:modular-z-minus-fc}. Lemma~\ref{lmm:fc-bounded-left-composition} with the bounded measurable $u = r_z$ (Lemma~\ref{lmm:resolvent-function-props}) gives $r_z(\Delta)g(\Delta)\psi = (r_z g)(\Delta)\psi = 1_{[0, \infty)}(\Delta)\psi$, which is $\psi$ by Lemma~\ref{lmm:modular-indicator-nonneg-one}.
\end{proof}

\begin{corollary}[The Resolvent of $\Delta$]
    \label{crllr:modular-resolvent}
    \lean{Physicslib4.Modular.modularResolvent_isResolvent}
    \uses{def:hall-9.16, lmm:uniqueness-of-resolvent, lmm:modular-resolvent-bounded}
    \leanok
    Every $z \in \mathbb{C} \setminus [0, \infty)$ lies in the resolvent set of $\Delta$ (\ref{def:hall-9.16}), and $(\Delta - z\mathbf{1})^{-1} = -r_z(\Delta)$. We write $(z - \Delta)^{-1} = r_z(\Delta)$. Thus $(z - \Delta)^{-1} \in \mathcal{B}(\mathbf{H})$, $\|(z - \Delta)^{-1}\| \le 1/d(z)$ and $\big( (z - \Delta)^{-1} \big)^* = (\overline{z} - \Delta)^{-1}$.
\end{corollary}
\begin{proof}
    \leanok
    \uses{lmm:modular-resolvent-right-inverse, lmm:modular-resolvent-left-inverse, lmm:modular-resolvent-bounded, lmm:uniqueness-of-resolvent}
    $B = -r_z(\Delta)$ is bounded, and properties 1 and 2 of Definition~\ref{def:hall-9.16} for $B$ are Lemmas~\ref{lmm:modular-resolvent-right-inverse} and~\ref{lmm:modular-resolvent-left-inverse}, multiplied by $-1$. Uniqueness is Lemma~\ref{lmm:uniqueness-of-resolvent}; the norm bound and the adjoint are Lemma~\ref{lmm:modular-resolvent-bounded}.
\end{proof}

\subsubsection{Takesaki's Lemma}

\begin{lemma}[An Elementary Inequality]
    \label{lmm:takesaki-elementary-inequality}
    \lean{Physicslib4.Modular.takesaki_elementary_inequality}
    \leanok
    Let $z \in \mathbb{C} \setminus [0, \infty)$. Then $|z| - \operatorname{Re} z > 0$, and for all real $p, q \ge 0$,
    \begin{align}
        |zp - q| \ge \sqrt{2(|z| - \operatorname{Re} z)\, pq}.
    \end{align}
\end{lemma}
\begin{proof}
    \leanok
    $\operatorname{Re} z \le |z|$, with equality only for $z \in [0, \infty)$ (Mathlib's \texttt{Complex.re\_le\_norm} and \texttt{Complex.re\_eq\_norm}). Expanding, $|zp - q|^2 = |z|^2p^2 - 2(\operatorname{Re} z)pq + q^2$, and $|z|^2p^2 + q^2 \ge 2|z|pq$ because $(|z|p - q)^2 \ge 0$; so $|zp - q|^2 \ge 2(|z| - \operatorname{Re} z)pq$. Take square roots.
\end{proof}

For the rest of the subsection fix $z \in \mathbb{C} \setminus [0, \infty)$ and $a' \in \mathcal{M}'$, and let
\begin{align}
    c = \frac{\|a'\|}{\sqrt{2(|z| - \operatorname{Re} z)}}, \qquad \psi = (z - \Delta)^{-1}a'\Omega, \qquad \mathcal{B}_c = \{ a \in \mathcal{M} \mid \|a\| \le c \},
\end{align}
and let $\mathcal{K}_c$ be the closure of $\mathcal{B}_c\Omega = \{ a\Omega \mid a \in \mathcal{B}_c \}$ in $\mathbf{H}$. For $u, v \in \mathbf{H}$ and $a \in \mathcal{M}$ write
\begin{align}
    G(u, v; a) = \operatorname{Re} \left< u, a'\Omega \right> - \operatorname{Re}\big( z\left< u, a\Omega \right> - \left< a^*\Omega, v \right> \big).
\end{align}

\begin{lemma}[$\mathcal{K}_c$ is Closed and Convex]
    \label{lmm:takesaki-ball-closed-convex}
    \lean{Physicslib4.Modular.takesakiBall_closure_isClosed_convex}
    \uses{lmm:commutant-von-neumann-algebra}
    \leanok
    $\mathcal{K}_c$ is closed and convex, and $0 \in \mathcal{K}_c$.
\end{lemma}
\begin{proof}
    \leanok
    \uses{lmm:commutant-von-neumann-algebra, def:von-neumann-algebra-commutant-convention}
    $\mathcal{M}$ is a linear subspace (Lemma~\ref{lmm:commutant-von-neumann-algebra} and \ref{def:von-neumann-algebra-commutant-convention}), so $\mathcal{B}_c$, its intersection with a ball, is convex, and so is the image $\mathcal{B}_c\Omega$ under the linear map $a \mapsto a\Omega$. Its closure $\mathcal{K}_c$ is closed and convex (Mathlib's \texttt{isClosed\_closure} and \texttt{Convex.closure}). Since $c \ge 0$, $0 \in \mathcal{B}_c$, so $0 = 0\Omega \in \mathcal{K}_c$.
\end{proof}

\begin{lemma}[Vectors in $\mathcal{K}_c$ Satisfy the Commutant Bound]
    \label{lmm:takesaki-ball-closure-bound}
    \lean{Physicslib4.Modular.takesakiBall_closure_bound}
    \uses{lmm:commutant-von-neumann-algebra}
    \leanok
    Every $\xi \in \mathcal{K}_c$ satisfies $\|b'\xi\| \le c\|b'\Omega\|$ for all $b' \in \mathcal{M}'$.
\end{lemma}
\begin{proof}
    \leanok
    \uses{def:von-neumann-algebra-commutant-convention, thrm:sequential-closedness}
    For $a \in \mathcal{B}_c$, $b'a = ab'$ (\ref{def:von-neumann-algebra-commutant-convention}), so $\|b'a\Omega\| = \|ab'\Omega\| \le c\|b'\Omega\|$. If $a_n\Omega \to \xi$ with $a_n \in \mathcal{B}_c$ (Theorem~\ref{thrm:sequential-closedness}), then $b'a_n\Omega \to b'\xi$ by continuity of $b'$, and the bound passes to the limit.
\end{proof}

\begin{lemma}[Separation from a Closed Convex Set]
    \label{lmm:takesaki-separation}
    \lean{Physicslib4.Modular.exists_re_inner_add_le_of_not_mem_convex}
    \leanok
    Let $\mathcal{K} \subset \mathbf{H}$ be nonempty, closed and convex, and $\xi \notin \mathcal{K}$. Then there are $\phi \in \mathbf{H}$ and $\gamma > 0$ with $\operatorname{Re}\left< \phi, k \right> + \gamma \le \operatorname{Re}\left< \phi, \xi \right>$ for every $k \in \mathcal{K}$.
\end{lemma}
\begin{proof}
    \leanok
    Regard $\mathbf{H}$ as a real inner product space with $\left< u, v \right>_{\mathbb{R}} = \operatorname{Re}\left< u, v \right>$ (Mathlib's \texttt{InnerProductSpace.rclikeToReal} and \texttt{real\_inner\_eq\_re\_inner}). Let $v \in \mathcal{K}$ be the nearest point to $\xi$ (Mathlib's \texttt{exists\_norm\_eq\_iInf\_of\_complete\_convex}) and $\phi = \xi - v \ne 0$. Then $\operatorname{Re}\left< \phi, k - v \right> \le 0$ for $k \in \mathcal{K}$ (Mathlib's \texttt{norm\_eq\_iInf\_iff\_real\_inner\_le\_zero}), and $\operatorname{Re}\left< \phi, \xi - v \right> = \|\phi\|^2$; take $\gamma = \|\phi\|^2$.
\end{proof}

\begin{lemma}[The Resolvent Vector $\chi$]
    \label{lmm:takesaki-transfer-resolvent}
    \lean{Physicslib4.Modular.takesaki_transfer_resolvent}
    \uses{crllr:modular-resolvent, def:tomita-operator}
    \leanok
    Let $\phi \in \mathbf{H}$ and $\chi = (\overline{z} - \Delta)^{-1}\phi$. Then $\chi \in \text{Dom}(\Delta) \subset \text{Dom}(S)$, $\overline{z}\chi - \Delta\chi = \phi$, and $\left< \phi, \psi \right> = \left< \chi, a'\Omega \right>$.
\end{lemma}
\begin{proof}
    \leanok
    \uses{crllr:modular-resolvent, lmm:modular-resolvent-right-inverse, def:adjoint-product-operator}
    $\overline{z} \notin [0, \infty)$, so Lemma~\ref{lmm:modular-resolvent-right-inverse} (for $\overline{z}$) gives $\chi \in \text{Dom}(\Delta)$ and $\overline{z}\chi - \Delta\chi = \phi$, and $\text{Dom}(\Delta) \subset \text{Dom}(S)$ by Definition~\ref{def:adjoint-product-operator}. By the adjoint formula of Corollary~\ref{crllr:modular-resolvent}, $\left< \phi, (z - \Delta)^{-1}a'\Omega \right> = \left< (\overline{z} - \Delta)^{-1}\phi, a'\Omega \right>$.
\end{proof}

\begin{lemma}[Transfer through the Resolvent]
    \label{lmm:takesaki-transfer}
    \lean{Physicslib4.Modular.takesaki_transfer}
    \uses{lmm:takesaki-transfer-resolvent, def:tomita-operator}
    \leanok
    Let $\phi \in \mathbf{H}$ and $\chi = (\overline{z} - \Delta)^{-1}\phi$. Then for every $a \in \mathcal{M}$,
    \begin{align}
        \left< \phi, a\Omega \right> = z\left< \chi, a\Omega \right> - \left< a^*\Omega, S\chi \right>.
    \end{align}
    In particular $\operatorname{Re}\left< \phi, \psi \right> - \operatorname{Re}\left< \phi, a\Omega \right> = G(\chi, S\chi; a)$.
\end{lemma}
\begin{proof}
    \leanok
    \uses{lmm:takesaki-transfer-resolvent, lmm:modular-operator-intrinsic, lmm:tomita-operator-independent-of-c, lmm:tomita-pre-operators-apply}
    By Lemma~\ref{lmm:takesaki-transfer-resolvent}, $\phi = \overline{z}\chi - \Delta\chi$, and the inner product is conjugate-linear in its first slot, so $\left< \phi, a\Omega \right> = z\left< \chi, a\Omega \right> - \left< \Delta\chi, a\Omega \right>$. Since $(a\Omega, a^*\Omega) \in \Gamma(S_0) \subset \Gamma(S)$ (Lemmas~\ref{lmm:tomita-pre-operators-apply} and~\ref{lmm:tomita-operator-independent-of-c}), $a\Omega \in \text{Dom}(S)$ and $S(a\Omega) = a^*\Omega$. The forward direction of Lemma~\ref{lmm:modular-operator-intrinsic}, for the vector $\chi \in \text{Dom}(\Delta)$ and its image $\Delta\chi$, tested against the vector $a\Omega \in \text{Dom}(S)$, gives $\left< a^*\Omega, S\chi \right> = \left< \Delta\chi, a\Omega \right>$. The last claim follows by taking real parts and using $\left< \phi, \psi \right> = \left< \chi, a'\Omega \right>$ (Lemma~\ref{lmm:takesaki-transfer-resolvent}).
\end{proof}

\begin{lemma}[$\mathcal{M}\Omega$ is a Core for $S$]
    \label{lmm:tomita-core}
    \lean{Physicslib4.Modular.tomitaOperator_exists_seq_tendsto}
    \uses{def:tomita-operator}
    \leanok
    For every $\chi \in \text{Dom}(S)$ there are $h_n \in \mathcal{M}$ with $h_n\Omega \to \chi$ and $h_n^*\Omega \to S\chi$.
\end{lemma}
\begin{proof}
    \leanok
    \uses{lmm:tomita-operator-independent-of-c, lmm:tomita-pre-operators-apply, thrm:sequential-closedness, lmm:componentwise-convergence}
    $(\chi, S\chi) \in \Gamma(S)$, the closure of $\Gamma(S_0) = \{ (h\Omega, h^*\Omega) \mid h \in \mathcal{M} \}$ (Lemmas~\ref{lmm:tomita-operator-independent-of-c} and~\ref{lmm:tomita-pre-operators-apply}). Take a sequence in $\Gamma(S_0)$ converging to it (Theorem~\ref{thrm:sequential-closedness}) and use Lemma~\ref{lmm:componentwise-convergence}.
\end{proof}

\begin{lemma}[$G$ is Lipschitz, Uniformly on $\mathcal{B}_c$]
    \label{lmm:takesaki-G-lipschitz}
    \lean{Physicslib4.Modular.takesakiG_lipschitz}
    \leanok
    For all $u, v, u_1, v_1 \in \mathbf{H}$ and $a \in \mathcal{B}_c$,
    \begin{align}
        |G(u, v; a) - G(u_1, v_1; a)| \le \big( \|a'\Omega\| + |z|c\|\Omega\| \big)\|u - u_1\| + c\|\Omega\|\,\|v - v_1\|.
    \end{align}
\end{lemma}
\begin{proof}
    \leanok
    By linearity, $G(u, v; a) - G(u_1, v_1; a) = \operatorname{Re}\left< u - u_1, a'\Omega \right> - \operatorname{Re}\big( z\left< u - u_1, a\Omega \right> \big) + \operatorname{Re}\left< a^*\Omega, v - v_1 \right>$. Since $\|a\| \le c$ and $\|a^*\| = \|a\|$ (Mathlib's \texttt{norm\_star}), $\|a\Omega\| \le c\|\Omega\|$ and $\|a^*\Omega\| \le c\|\Omega\|$. Bound each term by the Cauchy--Schwarz inequality (Mathlib's \texttt{norm\_inner\_le\_norm}) and $|\operatorname{Re} w| \le |w|$.
\end{proof}

\begin{lemma}[Uniform Approximation of $G$ from $\mathcal{M}\Omega$]
    \label{lmm:takesaki-uniform-approximation}
    \lean{Physicslib4.Modular.takesakiG_uniform_approximation}
    \uses{def:tomita-operator}
    \leanok
    Let $\chi \in \text{Dom}(S)$ and $\delta > 0$. Then there is $h \in \mathcal{M}$ with $|G(\chi, S\chi; a) - G(h\Omega, h^*\Omega; a)| < \delta$ for every $a \in \mathcal{B}_c$.
\end{lemma}
\begin{proof}
    \leanok
    \uses{lmm:tomita-core, lmm:takesaki-G-lipschitz}
    Let $L = \|a'\Omega\| + |z|c\|\Omega\| + c\|\Omega\| + 1 > 0$. Take $h_n$ from Lemma~\ref{lmm:tomita-core} and, by Mathlib's \texttt{Metric.tendsto\_atTop}, $n$ with $\|h_n\Omega - \chi\| < \delta/L$ and $\|h_n^*\Omega - S\chi\| < \delta/L$. For $h = h_n$, Lemma~\ref{lmm:takesaki-G-lipschitz} bounds the difference by $(L - 1)\delta/L < \delta$, for every $a \in \mathcal{B}_c$.
\end{proof}

The next lemmas bound $G(h\Omega, h^*\Omega; \cdot)$ from below. For $h \in \mathcal{M}$ and $\varepsilon > 0$ let $P = h^*h$ and, for $s \in \mathbb{R}$, $P_\varepsilon^s = f_s(P)$ in the continuous functional calculus of Mathlib, with $f_s(t) = (\varepsilon + \max(t, 0))^s$. Put
\begin{align}
    u_\varepsilon = hP_\varepsilon^{-1/2}, \qquad x_\varepsilon = P_\varepsilon^{-1/4}h^*, \qquad p_\varepsilon = \|x_\varepsilon^*\Omega\|, \qquad q_\varepsilon = \|x_\varepsilon\Omega\|.
\end{align}

\begin{lemma}[Functional Calculus Stays in $\mathcal{M}$]
    \label{lmm:cfc-mem-von-neumann}
    \lean{Physicslib4.Modular.cfc_mem_of_isVonNeumannSet}
    \uses{def:von-neumann-algebra-commutant-convention}
    \leanok
    Let $x \in \mathcal{M}$ be self-adjoint and $f : \mathbb{R} \to \mathbb{R}$. Then $f(x) \in \mathcal{M}$, where $f(x)$ is the continuous functional calculus of the C*-algebra $\mathcal{B}(\mathbf{H})$ (Mathlib's \texttt{cfc}).
\end{lemma}
\begin{proof}
    \leanok
    \uses{lmm:commutant-von-neumann-algebra}
    Every $d' \in \mathcal{M}'$ commutes with $x$, so it commutes with $f(x)$ (Mathlib's \texttt{Commute.cfc\_real}). Hence $f(x)$ lies in the centralizer of $\mathcal{M}'$, which is $\mathcal{M}$.
\end{proof}

\begin{lemma}[Shifted Powers of $h^*h$]
    \label{lmm:takesaki-shifted-powers}
    \lean{Physicslib4.Modular.takesakiShiftedPow_mem}
    \uses{lmm:cfc-mem-von-neumann}
    \leanok
    Let $h \in \mathcal{M}$, $\varepsilon > 0$ and $s, t \in \mathbb{R}$. Then $P_\varepsilon^s \in \mathcal{M}$, $P_\varepsilon^s$ is self-adjoint, $P_\varepsilon^sP_\varepsilon^t = P_\varepsilon^{s+t}$, $P_\varepsilon^0 = \mathbf{1}$, $P_\varepsilon^sP = PP_\varepsilon^s$, and $\|P_\varepsilon^s\| \le \varepsilon^s$ if $s \le 0$.
\end{lemma}
\begin{proof}
    \leanok
    \uses{lmm:cfc-mem-von-neumann, lmm:commutant-von-neumann-algebra}
    $P = h^*h \in \mathcal{M}$ is self-adjoint (Lemma~\ref{lmm:commutant-von-neumann-algebra}), and $f_s$ is real-valued and continuous because $\varepsilon + \max(t, 0) \ge \varepsilon > 0$. Lemma~\ref{lmm:cfc-mem-von-neumann} gives $P_\varepsilon^s \in \mathcal{M}$, and Mathlib's \texttt{IsSelfAdjoint.cfc} gives self-adjointness. Since $f_sf_t = f_{s+t}$ (Mathlib's \texttt{Real.rpow\_add}) and $f_0 = 1$, Mathlib's \texttt{cfc\_mul} and \texttt{cfc\_one} give the product rule and $P_\varepsilon^0 = \mathbf{1}$. $P = \iota(P)$ (Mathlib's \texttt{cfc\_id}), so $P_\varepsilon^s$ commutes with $P$ (Mathlib's \texttt{cfc\_commute\_cfc}). For $s \le 0$, $|f_s(t)| \le \varepsilon^s$ for all $t$ (Mathlib's \texttt{Real.rpow\_le\_rpow\_of\_nonpos}), so $\|P_\varepsilon^s\| \le \varepsilon^s$ (Mathlib's \texttt{norm\_cfc\_le}).
\end{proof}

\begin{lemma}[An Approximate Polar Decomposition]
    \label{lmm:takesaki-approximate-polar}
    \lean{Physicslib4.Modular.takesaki_approximate_polar}
    \uses{lmm:takesaki-shifted-powers}
    \leanok
    Let $h \in \mathcal{M}$ and $\varepsilon > 0$. Then $u_\varepsilon, x_\varepsilon \in \mathcal{M}$, $x_\varepsilon^* = hP_\varepsilon^{-1/4}$, and
    \begin{align}
        h^* = P_\varepsilon^{1/4}x_\varepsilon, \qquad h^*u_\varepsilon = x_\varepsilon x_\varepsilon^*, \qquad u_\varepsilon h^* = x_\varepsilon^*x_\varepsilon.
    \end{align}
\end{lemma}
\begin{proof}
    \leanok
    \uses{lmm:takesaki-shifted-powers, lmm:commutant-von-neumann-algebra}
    Membership follows from Lemma~\ref{lmm:takesaki-shifted-powers} and closure of $\mathcal{M}$ under products and adjoints (Lemma~\ref{lmm:commutant-von-neumann-algebra}); $x_\varepsilon^* = hP_\varepsilon^{-1/4}$ because $P_\varepsilon^{-1/4}$ is self-adjoint. By the product rule and $P_\varepsilon^0 = \mathbf{1}$ (Lemma~\ref{lmm:takesaki-shifted-powers}): $P_\varepsilon^{1/4}x_\varepsilon = P_\varepsilon^0h^* = h^*$; $u_\varepsilon h^* = hP_\varepsilon^{-1/4}P_\varepsilon^{-1/4}h^* = x_\varepsilon^*x_\varepsilon$; and, since $P$ commutes with $P_\varepsilon^{-1/4}$, $h^*u_\varepsilon = PP_\varepsilon^{-1/4}P_\varepsilon^{-1/4} = P_\varepsilon^{-1/4}h^*hP_\varepsilon^{-1/4} = x_\varepsilon x_\varepsilon^*$.
\end{proof}

\begin{lemma}[Norms in the Approximate Polar Decomposition]
    \label{lmm:takesaki-approx-polar-norms}
    \lean{Physicslib4.Modular.takesaki_approx_polar_norms}
    \uses{lmm:takesaki-approximate-polar}
    \leanok
    Let $h \in \mathcal{M}$ and $\varepsilon > 0$. Then $\|u_\varepsilon\| \le 1$ and $\|x_\varepsilon\|^2 \le \|h\|$.
\end{lemma}
\begin{proof}
    \leanok
    \uses{lmm:takesaki-shifted-powers, lmm:takesaki-approximate-polar}
    $P \ge 0$ (Mathlib's \texttt{star\_mul\_self\_nonneg}), so $\sigma(P) \subset [0, \|P\|]$ (Mathlib's \texttt{spectrum\_nonneg\_of\_nonneg} and \texttt{spectrum.norm\_le\_norm\_of\_mem}), and $\|P\| = \|h\|^2$ (Mathlib's \texttt{CStarRing.norm\_star\_mul\_self}). By the C*-identity, $\|u_\varepsilon\|^2 = \|u_\varepsilon^*u_\varepsilon\| = \|P_\varepsilon^{-1/2}PP_\varepsilon^{-1/2}\|$ and $\|x_\varepsilon\|^2 = \|x_\varepsilon x_\varepsilon^*\| = \|P_\varepsilon^{-1/4}PP_\varepsilon^{-1/4}\|$ (Mathlib's \texttt{CStarRing.norm\_self\_mul\_star} and Lemma~\ref{lmm:takesaki-approximate-polar}). By \texttt{cfc\_mul} and \texttt{cfc\_id} these operators are $g(P)$ and $k(P)$ with $g(t) = t(\varepsilon + \max(t, 0))^{-1}$ and $k(t) = t(\varepsilon + \max(t, 0))^{-1/2}$. On $[0, \|P\|]$, $0 \le g \le 1$ and $0 \le k(t) \le \sqrt{t} \le \|h\|$, so Mathlib's \texttt{norm\_cfc\_le} gives both bounds.
\end{proof}

\begin{lemma}[A Bound for $\|P_\varepsilon^{1/4}\Omega\|$]
    \label{lmm:takesaki-shifted-root-bound}
    \lean{Physicslib4.Modular.takesaki_shifted_root_bound}
    \uses{lmm:takesaki-approximate-polar}
    \leanok
    Let $h \in \mathcal{M}$ and $\varepsilon > 0$. Then $\|P_\varepsilon^{1/4}\Omega\|^2 \le p_\varepsilon^2 + \sqrt{\varepsilon}\|\Omega\|^2$.
\end{lemma}
\begin{proof}
    \leanok
    \uses{lmm:takesaki-shifted-powers, lmm:takesaki-approximate-polar}
    By Lemma~\ref{lmm:takesaki-shifted-powers}, $P_\varepsilon^{1/4}$ is self-adjoint and $P_\varepsilon^{1/4}P_\varepsilon^{1/4} = P_\varepsilon^{1/2}$, so $\|P_\varepsilon^{1/4}\Omega\|^2 = \left< \Omega, P_\varepsilon^{1/2}\Omega \right>$. Since $P = h^*h \ge 0$, $\sigma(P) \subset [0, \infty)$ (Mathlib's \texttt{star\_mul\_self\_nonneg} and \texttt{spectrum\_nonneg\_of\_nonneg}), and on $\sigma(P)$, $f_1(t) = \varepsilon + t$ (Mathlib's \texttt{Real.rpow\_one}), so $P_\varepsilon^1 = \varepsilon\mathbf{1} + P$ (Mathlib's \texttt{cfc\_congr}, \texttt{cfc\_add}, \texttt{cfc\_const} and \texttt{cfc\_id}), and $P_\varepsilon^{1/2} = P_\varepsilon^1P_\varepsilon^{-1/2} = PP_\varepsilon^{-1/2} + \varepsilon P_\varepsilon^{-1/2}$. Here $PP_\varepsilon^{-1/2} = h^*u_\varepsilon = x_\varepsilon x_\varepsilon^*$ (Lemma~\ref{lmm:takesaki-approximate-polar}), so $\left< \Omega, PP_\varepsilon^{-1/2}\Omega \right> = \|x_\varepsilon^*\Omega\|^2 = p_\varepsilon^2$. By the Cauchy--Schwarz inequality and $\|P_\varepsilon^{-1/2}\| \le \varepsilon^{-1/2}$, $|\left< \Omega, \varepsilon P_\varepsilon^{-1/2}\Omega \right>| \le \sqrt{\varepsilon}\|\Omega\|^2$. Take real parts.
\end{proof}

\begin{lemma}[An Upper Bound for $\left< h\Omega, a'\Omega \right>$]
    \label{lmm:takesaki-numerator-bound}
    \lean{Physicslib4.Modular.takesaki_numerator_bound}
    \uses{lmm:takesaki-approximate-polar}
    \leanok
    For $h \in \mathcal{M}$ and $\varepsilon > 0$, $|\left< h\Omega, a'\Omega \right>| \le \|a'\|\, \big( p_\varepsilon + \varepsilon^{1/4}\|\Omega\| \big)\, q_\varepsilon$.
\end{lemma}
\begin{proof}
    \leanok
    \uses{lmm:takesaki-shifted-powers, lmm:takesaki-approximate-polar, lmm:takesaki-shifted-root-bound, lmm:commutant-von-neumann-algebra, def:von-neumann-algebra-commutant-convention}
    By Lemma~\ref{lmm:takesaki-approximate-polar}, $h^* = P_\varepsilon^{1/4}x_\varepsilon$ with $x_\varepsilon \in \mathcal{M}$, so $a' \in \mathcal{M}'$ commutes with $x_\varepsilon$ (\ref{def:von-neumann-algebra-commutant-convention} and Lemma~\ref{lmm:commutant-von-neumann-algebra}), and $P_\varepsilon^{1/4}$ is self-adjoint (Lemma~\ref{lmm:takesaki-shifted-powers}). Hence
    \begin{align}
        \left< h\Omega, a'\Omega \right> = \left< \Omega, h^*a'\Omega \right> = \left< P_\varepsilon^{1/4}\Omega, a'x_\varepsilon\Omega \right>.
    \end{align}
    By the Cauchy--Schwarz inequality its modulus is at most $\|P_\varepsilon^{1/4}\Omega\|\,\|a'\|\,q_\varepsilon$, and $\|P_\varepsilon^{1/4}\Omega\| \le \sqrt{p_\varepsilon^2 + \sqrt{\varepsilon}\|\Omega\|^2} \le p_\varepsilon + \varepsilon^{1/4}\|\Omega\|$ (Lemma~\ref{lmm:takesaki-shifted-root-bound}).
\end{proof}

\begin{lemma}[A Complex Rotation]
    \label{lmm:complex-rotation-choice}
    \lean{Physicslib4.Modular.exists_norm_le_re_mul_eq_mul_norm}
    \leanok
    For $w \in \mathbb{C}$ and $c \ge 0$ there is $s \in \mathbb{C}$ with $|s| \le c$ and $\operatorname{Re}(sw) = c|w|$.
\end{lemma}
\begin{proof}
    \leanok
    If $w = 0$, take $s = 0$. Otherwise take $s = c\overline{w}/|w|$: then $|s| = c$ (Mathlib's \texttt{Complex.norm\_conj}) and $sw = c|w|^2/|w| = c|w|$ (Mathlib's \texttt{Complex.conj\_mul'}), which is real.
\end{proof}

\begin{lemma}[Inner Products of the Test Elements]
    \label{lmm:takesaki-test-inner-products}
    \lean{Physicslib4.Modular.takesaki_test_inner_products}
    \uses{lmm:takesaki-approximate-polar}
    \leanok
    For $h \in \mathcal{M}$, $\varepsilon > 0$ and $s \in \mathbb{C}$, $\left< h\Omega, s u_\varepsilon\Omega \right> = sp_\varepsilon^2$ and $\left< (s u_\varepsilon)^*\Omega, h^*\Omega \right> = sq_\varepsilon^2$.
\end{lemma}
\begin{proof}
    \leanok
    \uses{lmm:takesaki-approximate-polar}
    By Lemma~\ref{lmm:takesaki-approximate-polar}, $\left< h\Omega, su_\varepsilon\Omega \right> = s\left< \Omega, h^*u_\varepsilon\Omega \right> = s\left< \Omega, x_\varepsilon x_\varepsilon^*\Omega \right> = s\|x_\varepsilon^*\Omega\|^2$. Since $(su_\varepsilon)^* = \overline{s}u_\varepsilon^*$ and the inner product is conjugate-linear in its first slot, $\left< (su_\varepsilon)^*\Omega, h^*\Omega \right> = s\left< \Omega, u_\varepsilon h^*\Omega \right> = s\left< \Omega, x_\varepsilon^*x_\varepsilon\Omega \right> = s\|x_\varepsilon\Omega\|^2$.
\end{proof}

\begin{lemma}[A Test Element in $\mathcal{B}_c$]
    \label{lmm:takesaki-test-element}
    \lean{Physicslib4.Modular.takesaki_test_element}
    \uses{lmm:takesaki-approximate-polar}
    \leanok
    For $h \in \mathcal{M}$ and $\varepsilon > 0$ there is $a \in \mathcal{B}_c$ with
    \begin{align}
        \operatorname{Re}\big( z\left< h\Omega, a\Omega \right> - \left< a^*\Omega, h^*\Omega \right> \big) \ge \|a'\|\,p_\varepsilon q_\varepsilon.
    \end{align}
\end{lemma}
\begin{proof}
    \leanok
    \uses{lmm:complex-rotation-choice, lmm:takesaki-test-inner-products, lmm:takesaki-approximate-polar, lmm:takesaki-approx-polar-norms, lmm:takesaki-elementary-inequality, lmm:commutant-von-neumann-algebra}
    Let $w = zp_\varepsilon^2 - q_\varepsilon^2$, take $s$ from Lemma~\ref{lmm:complex-rotation-choice}, and let $a = su_\varepsilon$. Then $a \in \mathcal{M}$ (Lemmas~\ref{lmm:takesaki-approximate-polar} and~\ref{lmm:commutant-von-neumann-algebra}) and $\|a\| \le |s|\|u_\varepsilon\| \le c$ (Lemma~\ref{lmm:takesaki-approx-polar-norms}), so $a \in \mathcal{B}_c$. By Lemma~\ref{lmm:takesaki-test-inner-products} the left side is $\operatorname{Re}(sw) = c|w|$. Lemma~\ref{lmm:takesaki-elementary-inequality} with $p_\varepsilon^2, q_\varepsilon^2$ gives $|w| \ge \sqrt{2(|z| - \operatorname{Re} z)}\,p_\varepsilon q_\varepsilon$, and $c\sqrt{2(|z| - \operatorname{Re} z)} = \|a'\|$.
\end{proof}

\begin{lemma}[The Key Estimate]
    \label{lmm:takesaki-key-estimate}
    \lean{Physicslib4.Modular.takesaki_key_estimate}
    \uses{lmm:takesaki-test-element}
    \leanok
    For every $h \in \mathcal{M}$ and $\delta > 0$ there is $a \in \mathcal{B}_c$ with $G(h\Omega, h^*\Omega; a) \le \delta$.
\end{lemma}
\begin{proof}
    \leanok
    \uses{lmm:takesaki-numerator-bound, lmm:takesaki-test-element, lmm:takesaki-approx-polar-norms}
    Choose $\varepsilon > 0$ with $\varepsilon^{1/4}\|a'\|\,\|h\|^{1/2}\|\Omega\|^2 \le \delta$ and take $a$ from Lemma~\ref{lmm:takesaki-test-element}. By Lemmas~\ref{lmm:takesaki-numerator-bound} and~\ref{lmm:takesaki-approx-polar-norms} ($q_\varepsilon \le \|x_\varepsilon\|\|\Omega\| \le \|h\|^{1/2}\|\Omega\|$),
    \begin{align}
        G(h\Omega, h^*\Omega; a) \le \|a'\|(p_\varepsilon + \varepsilon^{1/4}\|\Omega\|)q_\varepsilon - \|a'\|p_\varepsilon q_\varepsilon \le \varepsilon^{1/4}\|a'\|\,\|h\|^{1/2}\|\Omega\|^2 \le \delta.
    \end{align}
\end{proof}

\begin{lemma}[The Resolvent Vector Lies in $\mathcal{K}_c$]
    \label{lmm:takesaki-resolvent-in-ball-closure}
    \lean{Physicslib4.Modular.takesaki_resolvent_mem_ball_closure}
    \uses{crllr:modular-resolvent, lmm:takesaki-ball-closed-convex, lmm:takesaki-ball-closure-bound}
    \leanok
    $\psi \in \mathcal{K}_c$.
\end{lemma}
\begin{proof}
    \leanok
    \uses{lmm:takesaki-ball-closed-convex, lmm:takesaki-separation, lmm:takesaki-transfer-resolvent, lmm:takesaki-transfer, lmm:takesaki-uniform-approximation, lmm:takesaki-key-estimate}
    Suppose not. $\mathcal{K}_c$ is nonempty, closed and convex (Lemma~\ref{lmm:takesaki-ball-closed-convex}), so Lemma~\ref{lmm:takesaki-separation} gives $\phi$ and $\gamma > 0$ with $\operatorname{Re}\left< \phi, \psi \right> - \operatorname{Re}\left< \phi, a\Omega \right> \ge \gamma$ for all $a \in \mathcal{B}_c$. With $\chi = (\overline{z} - \Delta)^{-1}\phi \in \text{Dom}(S)$ (Lemma~\ref{lmm:takesaki-transfer-resolvent}), Lemma~\ref{lmm:takesaki-transfer} gives $G(\chi, S\chi; a) \ge \gamma$ for all $a \in \mathcal{B}_c$. Lemma~\ref{lmm:takesaki-uniform-approximation} with $\delta = \gamma/2$ gives $h \in \mathcal{M}$ with $G(h\Omega, h^*\Omega; a) > \gamma/2$ for all $a \in \mathcal{B}_c$, contradicting Lemma~\ref{lmm:takesaki-key-estimate} with $\delta = \gamma/2$.
\end{proof}

\begin{theorem}[Takesaki's Lemma]
    \label{thrm:takesaki-lemma}
    \lean{Physicslib4.Modular.takesaki_lemma}
    \uses{crllr:modular-resolvent, def:cyclic-separating-vector, def:von-neumann-algebra-commutant-convention}
    \leanok
    Let $\mathcal{M}$ be a von Neumann algebra on $\mathbf{H}$, $\Omega$ cyclic and separating for $\mathcal{M}$, and $\Delta$ its modular operator. For every $z \in \mathbb{C} \setminus [0, \infty)$ and $a' \in \mathcal{M}'$ there is a unique $a \in \mathcal{M}$ with
    \begin{align}
        a\Omega = (z - \Delta)^{-1}a'\Omega, \qquad \text{and} \qquad \|a\| \le \frac{\|a'\|}{\sqrt{2(|z| - \operatorname{Re} z)}}.
    \end{align}
\end{theorem}
\begin{proof}
    \leanok
    \uses{lmm:takesaki-elementary-inequality, lmm:takesaki-resolvent-in-ball-closure, lmm:takesaki-ball-closure-bound, lmm:commutant-extension}
    The constant $c$ is well defined by Lemma~\ref{lmm:takesaki-elementary-inequality}. By Lemmas~\ref{lmm:takesaki-resolvent-in-ball-closure} and~\ref{lmm:takesaki-ball-closure-bound}, $\|b'\psi\| \le c\|b'\Omega\|$ for all $b' \in \mathcal{M}'$, and the commutant extension lemma (Lemma~\ref{lmm:commutant-extension}) gives the unique $a \in \mathcal{M}$ with $a\Omega = \psi$ and $\|a\| \le c$.
\end{proof}

\subsection{The Adjoint of the Tomita Operator}\label{sbsctn:tomita-adjoint}

This subsection proves that the antilinear adjoint of $S$ is the closure of $F_0$, and hence that the Tomita operator of $(\mathcal{M}', \Omega)$ is $S^*$, with modular operator $\Delta^{-1}$ and modular conjugation $J$. This is Theorem~13.1.3 and Corollary~13.1.4 of V.~F.~R.~Jones, \emph{von Neumann algebras}, Berkeley lecture notes, Chapter~13, a result of Takesaki, in the presentation of van Daele's proof of Tomita's theorem (A.~van Daele, \emph{J. Funct. Anal.} 15 (1974) 378--393) that Jones follows. The inclusion of the closure of $F_0$ in $S^*$ is the adjoint relation of Lemma~\ref{lmm:tomita-adjoint-relation}. For the converse, Jones takes, for $\xi \in \text{Dom}(S^*)$, the closure $c$ of $x\Omega \mapsto x\xi$ ($x \in \mathcal{M}$), an operator affiliated with $\mathcal{M}'$, and cuts it off with its polar decomposition. The partial isometry of that polar decomposition need not be unitary; the proof below uses instead the bounded operators $c(\mathbf{1} + |c|/n)^{-1}$; these commute with the unitaries of $\mathcal{M}$, so they lie in $\mathcal{M}'$.

The antilinear adjoint of $S$ is expressed through the linear adjoint $(S')^*$ of $S' = CS$. Write $S^*$ for the antilinear operator $(S')^*C$, with $\text{Dom}(S^*) = \{ \xi \in \mathbf{H} \mid C\xi \in \text{Dom}((S')^*) \}$ and $S^*\xi = (S')^*C\xi$, and $\Gamma(S^*)$ for its graph. Lemma~\ref{lmm:tomita-adjoint-intrinsic} shows that $\Gamma(S^*)$ is described by $S$ alone, so $S^*$ does not depend on $C$. Write $\Gamma(F_0) = \{ (b'\Omega, b'^*\Omega) \mid b' \in \mathcal{M}' \}$ for the graph of $F_0$ (Lemma~\ref{lmm:tomita-pre-operators-apply}).

\subsubsection{The Antilinear Adjoint of $S$}

\begin{lemma}[An Intrinsic Description of $S^*$]
    \label{lmm:tomita-adjoint-intrinsic}
    \lean{Physicslib4.Modular.mem_tomitaAdjoint_iff}
    \uses{def:tomita-operator, def:hall-9.1}
    \leanok
    Let $\xi, \eta \in \mathbf{H}$. Then $C\xi \in \text{Dom}((S')^*)$ and $(S')^*C\xi = \eta$ if and only if $\left< \xi, S\zeta \right> = \left< \zeta, \eta \right>$ for every $\zeta \in \text{Dom}(S)$.
\end{lemma}
\begin{proof}
    \leanok
    \uses{lmm:characterizing-adjoint-domain-membership, lmm:reference-conjugation-properties, crllr:linearized-tomita-properties}
    $\text{Dom}(S') = \text{Dom}(S)$, which is dense (Corollary~\ref{crllr:linearized-tomita-properties}), and $S'\zeta = CS\zeta$. By Lemma~\ref{lmm:characterizing-adjoint-domain-membership}, the left side holds if and only if $\left< C\xi, CS\zeta \right> = \left< \eta, \zeta \right>$ for every $\zeta \in \text{Dom}(S)$. By Lemma~\ref{lmm:reference-conjugation-properties}, $\left< C\xi, CS\zeta \right> = \overline{\left< \xi, S\zeta \right>}$, and $\left< \eta, \zeta \right> = \overline{\left< \zeta, \eta \right>}$; take complex conjugates.
\end{proof}

\begin{lemma}[$S^*$ Tested on $\mathcal{M}\Omega$]
    \label{lmm:tomita-adjoint-intrinsic-core}
    \lean{Physicslib4.Modular.mem_tomitaAdjoint_iff_inner_algebra}
    \uses{def:tomita-operator, def:hall-9.1}
    \leanok
    Let $\xi, \eta \in \mathbf{H}$. Then $C\xi \in \text{Dom}((S')^*)$ and $(S')^*C\xi = \eta$ if and only if
    \begin{align}
        \left< \xi, x^*\Omega \right> = \left< x\Omega, \eta \right> \qquad \text{for every } x \in \mathcal{M}.
    \end{align}
\end{lemma}
\begin{proof}
    \leanok
    \uses{lmm:tomita-adjoint-intrinsic, lmm:tomita-pre-operators-apply, lmm:tomita-operator-independent-of-c, lmm:tomita-core, prpstn:continuity-of-norm-and-inner-product}
    For $x \in \mathcal{M}$, $(x\Omega, x^*\Omega) \in \Gamma(S_0) \subset \Gamma(S)$ (Lemmas~\ref{lmm:tomita-pre-operators-apply} and~\ref{lmm:tomita-operator-independent-of-c}), so the condition is the case $\zeta = x\Omega$ of Lemma~\ref{lmm:tomita-adjoint-intrinsic}. Conversely, for $\zeta \in \text{Dom}(S)$ take $h_n \in \mathcal{M}$ with $h_n\Omega \to \zeta$ and $h_n^*\Omega \to S\zeta$ (Lemma~\ref{lmm:tomita-core}); the condition for $x = h_n$ passes to the limit (Proposition~\ref{prpstn:continuity-of-norm-and-inner-product}) and gives $\left< \xi, S\zeta \right> = \left< \zeta, \eta \right>$.
\end{proof}

\begin{lemma}[The Graph of $S^*$ is Closed]
    \label{lmm:tomita-adjoint-graph-closed}
    \lean{Physicslib4.Modular.tomitaAdjoint_graph_isClosed}
    \uses{def:tomita-operator, def:hall-9.1}
    \leanok
    The set $\Gamma(S^*) = \{ (\xi, \eta) \in \mathbf{H} \times \mathbf{H} \mid C\xi \in \text{Dom}((S')^*),\ (S')^*C\xi = \eta \}$ is closed in $\mathbf{H} \times \mathbf{H}$.
\end{lemma}
\begin{proof}
    \leanok
    \uses{lmm:tomita-adjoint-intrinsic-core}
    By Lemma~\ref{lmm:tomita-adjoint-intrinsic-core}, $\Gamma(S^*)$ is the intersection over $x \in \mathcal{M}$ of the sets $\{ (\xi, \eta) \mid \left< \xi, x^*\Omega \right> = \left< x\Omega, \eta \right> \}$. Each is the set where two continuous functions of $(\xi, \eta)$ agree (Mathlib's \texttt{isClosed\_eq}, with \texttt{continuous\_inner}), hence closed, and an intersection of closed sets is closed (Mathlib's \texttt{isClosed\_iInter}).
\end{proof}

\begin{lemma}[$S^*$ Extends the Closure of $F_0$]
    \label{lmm:tomita-adjoint-extends-closure}
    \lean{Physicslib4.Modular.mem_tomitaAdjoint_of_mem_closure_graph}
    \uses{def:tomita-operator, def:tomita-pre-operators, def:hall-9.1}
    \leanok
    The closure of $\Gamma(F_0)$ in $\mathbf{H} \times \mathbf{H}$ is contained in $\Gamma(S^*)$: if $(\xi, \eta)$ lies in the closure of $\{ (b'\Omega, b'^*\Omega) \mid b' \in \mathcal{M}' \}$, then $C\xi \in \text{Dom}((S')^*)$ and $(S')^*C\xi = \eta$.
\end{lemma}
\begin{proof}
    \leanok
    \uses{lmm:tomita-adjoint-relation, lmm:tomita-pre-operators-apply, lmm:tomita-adjoint-intrinsic-core, lmm:tomita-adjoint-graph-closed}
    For $b' \in \mathcal{M}'$ and $x \in \mathcal{M}$, the adjoint relation (Lemma~\ref{lmm:tomita-adjoint-relation}, with Lemma~\ref{lmm:tomita-pre-operators-apply}) reads $\left< b'\Omega, x^*\Omega \right> = \left< x\Omega, b'^*\Omega \right>$, so $(b'\Omega, b'^*\Omega) \in \Gamma(S^*)$ by Lemma~\ref{lmm:tomita-adjoint-intrinsic-core}. $\Gamma(S^*)$ is closed (Lemma~\ref{lmm:tomita-adjoint-graph-closed}), so it contains the closure (Mathlib's \texttt{closure\_minimal}).
\end{proof}

\subsubsection{Operators Commuting with a Unitary}

\begin{definition}[A Unitary Commuting with an Operator]
    \label{def:unitary-commutes}
    \lean{Physicslib4.Spectral.Unbounded.UnitaryCommutes}
    \uses{def:unitary-operator}
    \leanok
    Let $U$ be unitary and $A$ an operator on $\mathbf{H}$. $U$ \emph{commutes with} $A$ if, for every $\psi \in \text{Dom}(A)$, the vectors $U\psi$ and $U^{-1}\psi$ lie in $\text{Dom}(A)$ and $AU\psi = UA\psi$, $AU^{-1}\psi = U^{-1}A\psi$.
\end{definition}

For $A \in \mathcal{B}(\mathbf{H})$ the condition of Definition~\ref{def:unitary-commutes} says $UA = AU$. Unitary conjugation enters only through the antiunitaries $C$ and $UC$ (Lemmas~\ref{lmm:reference-conjugation-properties} and~\ref{lmm:antiunitary-composition}): for an operator $B$, $(UC)\big(CBC^{-1}\big)(UC)^{-1}$ denotes the conjugate by the antiunitary $UC$ of the conjugate $CBC^{-1}$ by $C$ (\ref{def:product-of-unbounded-operators}). Its domain is $\{ \phi \mid C^{-1}(UC)^{-1}\phi \in \text{Dom}(B) \}$ and its value at such $\phi$ is $UC\,C\,B\,C^{-1}(UC)^{-1}\phi$.

\begin{lemma}[The Double Conjugate by a Commuting Unitary]
    \label{lmm:commuting-unitary-double-conjugate}
    \lean{Physicslib4.Spectral.Unbounded.antiunitaryConj_trans_antiunitaryConj_of_commute}
    \uses{def:unitary-operator, def:reference-conjugation, def:product-of-unbounded-operators, def:unitary-commutes}
    \leanok
    Let $U$ be a unitary that commutes with an operator $A$. Then $(UC)\big(CAC^{-1}\big)(UC)^{-1} = A$, including domains.
\end{lemma}
\begin{proof}
    \leanok
    \uses{lmm:reference-conjugation-properties, lmm:antiunitary-composition}
    $C^{-1} = C$ because $C(C\phi) = \phi$ (Lemma~\ref{lmm:reference-conjugation-properties}), so $C^{-1}(UC)^{-1}\phi = C(CU^{-1}\phi) = U^{-1}\phi$ and $UC\,C\chi = U\chi$. Hence the left side has domain $\{ \phi \mid U^{-1}\phi \in \text{Dom}(A) \}$ and value $UAU^{-1}\phi$ there. If $\phi \in \text{Dom}(A)$ then $U^{-1}\phi \in \text{Dom}(A)$ and $UAU^{-1}\phi = UU^{-1}A\phi = A\phi$; if $U^{-1}\phi \in \text{Dom}(A)$ then $\phi = U(U^{-1}\phi) \in \text{Dom}(A)$. The two operators have the same graph (Mathlib's \texttt{LinearPMap.ext}).
\end{proof}

\begin{lemma}[Unitary Conjugation and the Functional Calculus]
    \label{lmm:unitary-conjugate-functional-calculus}
    \lean{Physicslib4.Spectral.Unbounded.antiunitaryConj_trans_functionalCalculus_of_commute}
    \uses{def:unitary-operator, def:reference-conjugation, def:product-of-unbounded-operators, def:hall-10.5, def:unitary-commutes}
    \leanok
    Let $A$ be self-adjoint and $U$ a unitary that commutes with $A$. Then $(UC)\big(Cf(A)C^{-1}\big)(UC)^{-1} = f(A)$, including domains, for every real measurable $f$.
\end{lemma}
\begin{proof}
    \leanok
    \uses{lmm:reference-conjugation-properties, lmm:antiunitary-composition, lmm:antiunitary-conjugate-spectral-measure, crllr:antiunitary-conjugate-functional-calculus, lmm:commuting-unitary-double-conjugate}
    $C$ and $UC$ are antiunitary (Lemmas~\ref{lmm:reference-conjugation-properties} and~\ref{lmm:antiunitary-composition}). $B = CAC^{-1}$ is self-adjoint (Lemma~\ref{lmm:antiunitary-conjugate-spectral-measure}), and Corollary~\ref{crllr:antiunitary-conjugate-functional-calculus}, applied to $C$ and $A$ and then to $UC$ and $B$, gives
    \begin{align}
        (UC)\big(Cf(A)C^{-1}\big)(UC)^{-1} = (UC)f(B)(UC)^{-1} = f\big((UC)B(UC)^{-1}\big).
    \end{align}
    $(UC)B(UC)^{-1} = A$ by Lemma~\ref{lmm:commuting-unitary-double-conjugate}.
\end{proof}

\begin{lemma}[One Half of the Adjoint Commutation]
    \label{lmm:commuting-unitary-adjoint-half}
    \lean{Physicslib4.Spectral.Unbounded.adjoint_commute_of_symm_commute}
    \uses{def:unitary-operator, def:hall-9.1, def:unitary-commutes}
    \leanok
    Let $c$ be a densely defined operator and $V$ a unitary such that, for every $\psi \in \text{Dom}(c)$, $V^{-1}\psi \in \text{Dom}(c)$ and $cV^{-1}\psi = V^{-1}c\psi$. Then, for every $\phi \in \text{Dom}(c^*)$, $V\phi \in \text{Dom}(c^*)$ and $c^*V\phi = Vc^*\phi$.
\end{lemma}
\begin{proof}
    \leanok
    \uses{lmm:characterizing-adjoint-domain-membership}
    Let $\phi \in \text{Dom}(c^*)$. For $\psi \in \text{Dom}(c)$, since $V^* = V^{-1}$,
    \begin{align}
        \left< V\phi, c\psi \right> = \left< \phi, V^{-1}c\psi \right> = \left< \phi, cV^{-1}\psi \right> = \left< c^*\phi, V^{-1}\psi \right> = \left< Vc^*\phi, \psi \right>.
    \end{align}
    By Lemma~\ref{lmm:characterizing-adjoint-domain-membership}, $V\phi \in \text{Dom}(c^*)$ and $c^*V\phi = Vc^*\phi$.
\end{proof}

\begin{lemma}[A Commuting Unitary Commutes with the Adjoint]
    \label{lmm:commuting-unitary-adjoint}
    \lean{Physicslib4.Spectral.Unbounded.adjoint_commute_of_commute}
    \uses{def:unitary-operator, def:hall-9.1, def:unitary-commutes}
    \leanok
    Let $c$ be a densely defined operator and $U$ a unitary that commutes with $c$. Then $U$ commutes with $c^*$.
\end{lemma}
\begin{proof}
    \leanok
    \uses{lmm:commuting-unitary-adjoint-half}
    Apply Lemma~\ref{lmm:commuting-unitary-adjoint-half} with $V = U$, whose hypothesis is the $U^{-1}$ half of the commutation of $U$ with $c$, and with $V = U^{-1}$, whose hypothesis is the $U$ half since $(U^{-1})^{-1} = U$.
\end{proof}

\begin{lemma}[A Commuting Unitary Commutes with $c^*c$]
    \label{lmm:commuting-unitary-adjoint-product}
    \lean{Physicslib4.Spectral.Unbounded.adjointMulSelf_commute_of_commute}
    \uses{def:unitary-operator, def:adjoint-product-operator, def:unitary-commutes}
    \leanok
    Let $c$ be a densely defined operator and $U$ a unitary that commutes with $c$. Then $U$ commutes with $c^*c$.
\end{lemma}
\begin{proof}
    \leanok
    \uses{lmm:commuting-unitary-adjoint}
    Let $\psi \in \text{Dom}(c^*c)$, so $\psi \in \text{Dom}(c)$ and $c\psi \in \text{Dom}(c^*)$. Then $U\psi \in \text{Dom}(c)$ with $cU\psi = Uc\psi \in \text{Dom}(c^*)$ and $c^*cU\psi = c^*Uc\psi = Uc^*c\psi$ by Lemma~\ref{lmm:commuting-unitary-adjoint}; likewise for $U^{-1}$.
\end{proof}

\begin{lemma}[A Commuting Unitary Commutes with Bounded Functions of $c^*c$]
    \label{lmm:commuting-unitary-bounded-function}
    \lean{Physicslib4.Spectral.Unbounded.functionalCalculus_adjointMulSelf_commute_of_commute}
    \uses{def:unitary-operator, def:adjoint-product-operator, def:hall-10.5, def:unitary-commutes}
    \leanok
    Let $c$ be a closed, densely defined operator, $D = c^*c$, and $U$ a unitary that commutes with $c$. Then $Uf(D) = f(D)U$ for every bounded real measurable $f$.
\end{lemma}
\begin{proof}
    \leanok
    \uses{lmm:commuting-unitary-adjoint-product, lmm:unitary-conjugate-functional-calculus, thrm:von-neumann-adjoint-product, lmm:functional-calculus-bounded, lmm:reference-conjugation-properties}
    $D$ is self-adjoint (Theorem~\ref{thrm:von-neumann-adjoint-product}) and $U$ commutes with $D$ (Lemma~\ref{lmm:commuting-unitary-adjoint-product}), so $(UC)\big(Cf(D)C^{-1}\big)(UC)^{-1} = f(D)$ by Lemma~\ref{lmm:unitary-conjugate-functional-calculus}. $f(D)$ is defined on all of $\mathbf{H}$ (Lemma~\ref{lmm:functional-calculus-bounded}). Evaluate both sides at $U\psi$: since $C^{-1}(UC)^{-1}U\psi = C(CU^{-1}U\psi) = \psi$ (Lemma~\ref{lmm:reference-conjugation-properties}), the left side is $UC\,C\,f(D)\psi = Uf(D)\psi$, and the right side is $f(D)U\psi$.
\end{proof}

\subsubsection{Bounded Approximants of a Closed Operator}

For the next six results, $c$ is a closed, densely defined operator on $\mathbf{H}$, $D = c^*c$ and $|c| = s(D)$ (Definition~\ref{def:absolute-value}). For an integer $n \ge 1$ let
\begin{align}
    h_n(\lambda) = \frac{n}{n + s(\lambda)}, \qquad k_n(\lambda) = s(\lambda)h_n(\lambda), \qquad R_n = h_n(D), \qquad K_n = k_n(D).
\end{align}
Thus $R_n = (\mathbf{1} + |c|/n)^{-1}$.

\begin{lemma}[The Scalar Functions $h_n$ and $k_n$]
    \label{lmm:approximant-functions}
    \lean{Physicslib4.Spectral.Unbounded.approximant_functions}
    \leanok
    For every $n \ge 1$, $h_n$ and $k_n$ are real, continuous and measurable, $0 < h_n \le 1$, $0 \le k_n \le n$ and $k_n \le s$. For every $\lambda \in \mathbb{R}$, $h_n(\lambda) \to 1$ and $k_n(\lambda) \to s(\lambda)$ as $n \to \infty$.
\end{lemma}
\begin{proof}
    \leanok
    $s$ is continuous and non-negative (Mathlib's \texttt{Real.continuous\_sqrt} and \texttt{continuous\_max}), so $n + s > 0$ and $h_n$, $k_n$ are continuous, hence measurable (Mathlib's \texttt{Continuous.measurable}). $0 < h_n \le 1$ because $0 < n \le n + s$; $k_n = ns/(n + s)$ is at most $n$ and at most $s$. For fixed $\lambda$, $h_n(\lambda) = (1 + s(\lambda)/n)^{-1}$ and $s(\lambda)/n \to 0$ (Mathlib's \texttt{tendsto\_const\_div\_atTop\_nhds\_zero\_nat}), so $h_n(\lambda) \to 1$ by continuity of $t \mapsto (1 + t)^{-1}$ at $0$, and $k_n(\lambda) = s(\lambda)h_n(\lambda) \to s(\lambda)$.
\end{proof}

\begin{lemma}[The Bounded Operators $R_n$ and $K_n$]
    \label{lmm:approximant-bounded}
    \lean{Physicslib4.Spectral.Unbounded.approximant_bounded}
    \uses{def:absolute-value, def:hall-10.5}
    \leanok
    For every $n \ge 1$, $R_n$ and $K_n$ are bounded operators defined on all of $\mathbf{H}$, $R_n$ is self-adjoint, $\|R_n\| \le 1$ and $\|K_n\| \le n$.
\end{lemma}
\begin{proof}
    \leanok
    \uses{lmm:approximant-functions, thrm:von-neumann-adjoint-product, lmm:functional-calculus-bounded, prpstn:integral-norm-bound, prpstn:integral-conjugation}
    $D$ is self-adjoint (Theorem~\ref{thrm:von-neumann-adjoint-product}), and $h_n$, $k_n$ are bounded, real and measurable (Lemma~\ref{lmm:approximant-functions}). Lemma~\ref{lmm:functional-calculus-bounded} identifies $R_n$ and $K_n$ with bounded integrals; Proposition~\ref{prpstn:integral-norm-bound} with $0 < h_n \le 1$ and $0 \le k_n \le n$ gives the norm bounds, and Proposition~\ref{prpstn:integral-conjugation} gives $R_n^* = R_n$.
\end{proof}

\begin{lemma}[$R_n$ Maps into $\text{Dom}(c)$]
    \label{lmm:approximant-range}
    \lean{Physicslib4.Spectral.Unbounded.approximant_range}
    \uses{lmm:approximant-bounded, def:absolute-value}
    \leanok
    For every $n \ge 1$ and $\psi \in \mathbf{H}$, $R_n\psi \in \text{Dom}(|c|) = \text{Dom}(c)$, $|c|R_n\psi = K_n\psi$, and $\|cR_n\psi\| = \|K_n\psi\| \le n\|\psi\|$.
\end{lemma}
\begin{proof}
    \leanok
    \uses{thrm:von-neumann-adjoint-product, thrm:fc-composition, lmm:approximant-functions, lmm:approximant-bounded, lmm:functional-calculus-bounded, thrm:absolute-value-domain}
    $D$ is self-adjoint (Theorem~\ref{thrm:von-neumann-adjoint-product}). Theorem~\ref{thrm:fc-composition} with $f = s$ and $g = h_n$ gives $\text{Dom}(s(D)h_n(D)) = W_{h_n} \cap W_{k_n}$, which is $\mathbf{H}$ because $h_n$ and $k_n$ are bounded (Lemmas~\ref{lmm:approximant-functions} and~\ref{lmm:functional-calculus-bounded}), and $s(D)h_n(D)\psi = k_n(D)\psi$. Theorem~\ref{thrm:absolute-value-domain} gives $\text{Dom}(|c|) = \text{Dom}(c)$ and $\|cR_n\psi\| = \||c|R_n\psi\|$, and $\|K_n\| \le n$ (Lemma~\ref{lmm:approximant-bounded}).
\end{proof}

\begin{lemma}[Convergence of the Approximants]
    \label{lmm:approximant-convergence}
    \lean{Physicslib4.Spectral.Unbounded.approximant_convergence}
    \uses{lmm:approximant-bounded, def:absolute-value}
    \leanok
    $R_n\phi \to \phi$ for every $\phi \in \mathbf{H}$, and $K_n\psi \to |c|\psi$ for every $\psi \in \text{Dom}(c)$, as $n \to \infty$.
\end{lemma}
\begin{proof}
    \leanok
    \uses{lmm:fc-dominated-convergence, lmm:approximant-functions, lmm:functional-calculus-bounded, prpstn:integral-of-indicator, def:projection-valued-measure, thrm:absolute-value-domain}
    Apply Lemma~\ref{lmm:fc-dominated-convergence} to $\mu^D$ twice, with the bounds and limits of Lemma~\ref{lmm:approximant-functions}. First $h = 1$ and the $h_n$: here $W_1 = \mathbf{H}$ and $1(D) = \mu^D(\mathbb{R}) = \mathbf{1}$ (Lemma~\ref{lmm:functional-calculus-bounded}, Proposition~\ref{prpstn:integral-of-indicator} and Definition~\ref{def:projection-valued-measure}). Second $h = s$ and the $k_n$, for $\psi \in W_s = \text{Dom}(|c|) = \text{Dom}(c)$ (Theorem~\ref{thrm:absolute-value-domain}).
\end{proof}

\begin{lemma}[The Bounded Operators $cR_n$]
    \label{lmm:approximant-operator}
    \lean{Physicslib4.Spectral.Unbounded.approximant_operator}
    \uses{lmm:approximant-range}
    \leanok
    For every $n \ge 1$ there is a unique $x_n \in \mathcal{B}(\mathbf{H})$ with $x_n\psi = cR_n\psi$ for every $\psi \in \mathbf{H}$, and $\|x_n\| \le n$. Moreover $x_n^*\phi = R_nc^*\phi$ for every $\phi \in \text{Dom}(c^*)$.
\end{lemma}
\begin{proof}
    \leanok
    \uses{lmm:approximant-range, lmm:approximant-bounded, def:hall-9.1}
    $\psi \mapsto cR_n\psi$ is defined on all of $\mathbf{H}$ (Lemma~\ref{lmm:approximant-range}) and linear, as a composite of the linear $R_n$ with the linear $c$, and $\|cR_n\psi\| \le n\|\psi\|$ (Lemma~\ref{lmm:approximant-range}); Mathlib's \texttt{LinearMap.mkContinuous} gives $x_n$ with $\|x_n\| \le n$. For $\phi \in \text{Dom}(c^*)$ and $\psi \in \mathbf{H}$, the definition of the adjoint (\ref{def:hall-9.1}) and $R_n^* = R_n$ (Lemma~\ref{lmm:approximant-bounded}) give $\left< \phi, x_n\psi \right> = \left< c^*\phi, R_n\psi \right> = \left< R_nc^*\phi, \psi \right>$. Since $\left< x_n^*\phi, \psi \right> = \left< \phi, x_n\psi \right>$ (Mathlib's \texttt{ContinuousLinearMap.adjoint\_inner\_left}) for every $\psi$, $x_n^*\phi = R_nc^*\phi$ (Mathlib's \texttt{ext\_inner\_right}).
\end{proof}

\begin{lemma}[A Commuting Unitary Commutes with $cR_n$]
    \label{lmm:approximant-operator-commutes}
    \lean{Physicslib4.Spectral.Unbounded.approximant_operator_commute}
    \uses{lmm:approximant-operator, def:unitary-operator, def:unitary-commutes}
    \leanok
    Let $U$ be a unitary that commutes with $c$. Then $Ux_n = x_nU$ for every $n \ge 1$.
\end{lemma}
\begin{proof}
    \leanok
    \uses{lmm:commuting-unitary-bounded-function, lmm:approximant-functions, lmm:approximant-range, lmm:approximant-operator}
    $UR_n = R_nU$ by Lemma~\ref{lmm:commuting-unitary-bounded-function} with the bounded real measurable $h_n$ (Lemma~\ref{lmm:approximant-functions}). For $\psi \in \mathbf{H}$, $R_n\psi \in \text{Dom}(c)$ (Lemma~\ref{lmm:approximant-range}), so $Ux_n\psi = UcR_n\psi = cUR_n\psi = cR_nU\psi = x_nU\psi$.
\end{proof}

\subsubsection{Unitaries of a Von Neumann Algebra}

\begin{lemma}[A Unitary from a Self-Adjoint Contraction]
    \label{lmm:von-neumann-unitary-of-contraction}
    \lean{Physicslib4.Modular.self_add_I_smul_cfcSqrt_mem_of_isVonNeumannSet}
    \uses{def:von-neumann-algebra-commutant-convention, def:unitary-operator}
    \leanok
    Let $a \in \mathcal{M}$ be self-adjoint with $\|a\| \le 1$. Then $u = a + i\sqrt{\mathbf{1} - a^2}$, with the square root of the continuous functional calculus of $\mathcal{B}(\mathbf{H})$ (Mathlib's \texttt{CFC.sqrt}), is unitary, and $u, u^* \in \mathcal{M}$.
\end{lemma}
\begin{proof}
    \leanok
    \uses{lmm:cfc-mem-von-neumann, lmm:commutant-von-neumann-algebra}
    $u$ is unitary by Mathlib's \texttt{IsSelfAdjoint.self\_add\_I\_smul\_cfcSqrt\_sub\_sq\_mem\_unitary} (for the Loewner order of $\mathcal{B}(\mathbf{H})$, \texttt{ContinuousLinearMap.instStarOrderedRing}). Since $\|a\| \le 1$, $\mathbf{1} - a^2 \ge 0$, so $\sqrt{\mathbf{1} - a^2}$ is the continuous functional calculus of the self-adjoint $\mathbf{1} - a^2 \in \mathcal{M}$ with the real function $\sqrt{\cdot}$ (Mathlib's \texttt{CFC.sqrt\_eq\_real\_sqrt}), which lies in $\mathcal{M}$ by Lemma~\ref{lmm:cfc-mem-von-neumann}. $\mathcal{M}$ is a self-adjoint subalgebra (Lemma~\ref{lmm:commutant-von-neumann-algebra}), so $u, u^* \in \mathcal{M}$.
\end{proof}

\begin{lemma}[Real and Imaginary Parts in a Von Neumann Algebra]
    \label{lmm:von-neumann-real-imaginary-part}
    \lean{Physicslib4.Modular.realPart_imaginaryPart_mem_of_isVonNeumannSet}
    \uses{def:von-neumann-algebra-commutant-convention}
    \leanok
    Let $y \in \mathcal{M}$. Then the real part $\Re y = \tfrac{1}{2}(y + y^*)$ and the imaginary part $\Im y = -\tfrac{i}{2}(y - y^*)$ of $y$ in $\mathcal{B}(\mathbf{H})$ (Mathlib's \texttt{realPart} and \texttt{imaginaryPart}) lie in $\mathcal{M}$.
\end{lemma}
\begin{proof}
    \leanok
    \uses{lmm:commutant-von-neumann-algebra}
    Mathlib's \texttt{realPart\_apply\_coe} and \texttt{imaginaryPart\_apply\_coe} give the two formulas. $\mathcal{M}$ is a complex linear subspace closed under adjoints (Lemma~\ref{lmm:commutant-von-neumann-algebra}), so $y^* \in \mathcal{M}$ and both combinations lie in $\mathcal{M}$.
\end{proof}

\begin{lemma}[Unitaries Span a Von Neumann Algebra]
    \label{lmm:von-neumann-span-unitaries}
    \lean{Physicslib4.Modular.exists_unitary_decomposition_of_isVonNeumannSet}
    \uses{def:von-neumann-algebra-commutant-convention, def:unitary-operator}
    \leanok
    Let $x \in \mathcal{M}$ with $x \ne 0$. Then there are unitaries $u_1, u_2$ with $u_1, u_1^*, u_2, u_2^* \in \mathcal{M}$ and $x = \tfrac{\|x\|}{2}\big(u_1 + u_1^* + i(u_2 + u_2^*)\big)$.
\end{lemma}
\begin{proof}
    \leanok
    \uses{lmm:von-neumann-unitary-of-contraction, lmm:von-neumann-real-imaginary-part, lmm:commutant-von-neumann-algebra}
    Mathlib's \texttt{CStarAlgebra.norm\_smul\_two\_inv\_smul\_add\_four\_unitary}, in the C*-algebra $\mathcal{B}(\mathbf{H})$, gives the formula with $u_1$ and $u_2$ the unitaries $a + i\sqrt{\mathbf{1} - a^2}$ (Mathlib's \texttt{selfAdjoint.unitarySelfAddISMul}) for the self-adjoint contractions $a = \Re y$ and $a = \Im y$, $y = \|x\|^{-1}x$. Here $y \in \mathcal{M}$ (Lemma~\ref{lmm:commutant-von-neumann-algebra}), so $\Re y, \Im y \in \mathcal{M}$ (Lemma~\ref{lmm:von-neumann-real-imaginary-part}), and Lemma~\ref{lmm:von-neumann-unitary-of-contraction} gives $u_1, u_1^*, u_2, u_2^* \in \mathcal{M}$.
\end{proof}

\begin{lemma}[Commuting with the Unitaries of $\mathcal{M}$]
    \label{lmm:commutant-of-unitaries}
    \lean{Physicslib4.Modular.mem_centralizer_of_commute_unitary}
    \uses{def:von-neumann-algebra-commutant-convention, def:unitary-operator}
    \leanok
    If $b \in \mathcal{B}(\mathbf{H})$ satisfies $ub = bu$ for every unitary $u \in \mathcal{M}$, then $b \in \mathcal{M}'$.
\end{lemma}
\begin{proof}
    \leanok
    \uses{lmm:von-neumann-span-unitaries}
    Let $x \in \mathcal{M}$; for $x = 0$, $xb = 0 = bx$. Otherwise write $x = \tfrac{\|x\|}{2}\big(u_1 + u_1^* + i(u_2 + u_2^*)\big)$ with $u_1, u_1^*, u_2, u_2^*$ unitaries in $\mathcal{M}$ (Lemma~\ref{lmm:von-neumann-span-unitaries}). Each of the four commutes with $b$, so $xb = bx$, and $b$ lies in the centralizer of $\mathcal{M}$.
\end{proof}

\subsubsection{$\mathcal{M}'\Omega$ is a Core for $S^*$}

For the next six results fix $\xi, \eta \in \mathbf{H}$ with $C\xi \in \text{Dom}((S')^*)$ and $(S')^*C\xi = \eta$; by Lemma~\ref{lmm:tomita-adjoint-intrinsic-core},
\begin{align}
    \left< \xi, x^*\Omega \right> = \left< x\Omega, \eta \right> \qquad \text{for every } x \in \mathcal{M}. \label{eq:tomita-adjoint-pair}
\end{align}

\begin{definition}[The Operator $x\Omega \mapsto x\xi$]
    \label{def:tomita-left-map}
    \lean{Physicslib4.Modular.tomitaLeftMap}
    \uses{def:cyclic-separating-vector, def:hall-9.1}
    \leanok
    For $\xi \in \mathbf{H}$, $a_0$ is the operator whose graph is the subspace $\{ (x\Omega, x\xi) \mid x \in \text{span}\,\mathcal{M} \}$ of $\mathbf{H} \times \mathbf{H}$ (Mathlib's \texttt{Submodule.toLinearPMap}, as for $S_0'$ in~\ref{def:linearized-tomita-pre-operator}). Lemma~\ref{lmm:tomita-adjoint-left-map} shows that this subspace is a graph when $\Omega$ is cyclic and separating.
\end{definition}

\begin{lemma}[The Operator $x\Omega \mapsto x\xi$]
    \label{lmm:tomita-adjoint-left-map}
    \lean{Physicslib4.Modular.tomitaLeftMap_properties}
    \uses{def:cyclic-separating-vector, lmm:commutant-von-neumann-algebra, def:tomita-left-map}
    \leanok
    $\text{Dom}(a_0) = \mathcal{M}\Omega$, $a_0(x\Omega) = x\xi$ for every $x \in \mathcal{M}$, and $a_0$ is densely defined.
\end{lemma}
\begin{proof}
    \leanok
    \uses{lmm:cyclic-commutant-separating, lmm:commutant-von-neumann-algebra, def:cyclic-separating-vector}
    $\mathcal{M}$ is a linear subspace (Lemma~\ref{lmm:commutant-von-neumann-algebra}), so $\mathcal{M}\Omega$ is the range of the linear map $x \mapsto x\Omega$, and $x \mapsto x\xi$ is linear. The first map is injective because $\Omega$ is separating for $\mathcal{M}$ (Lemma~\ref{lmm:cyclic-commutant-separating}), so the subspace of~\ref{def:tomita-left-map} is the graph of the second map composed with the inverse of the first on its range (Mathlib's \texttt{LinearEquiv.ofInjective}), and \texttt{Submodule.toLinearPMap} returns that operator. $\mathcal{M}\Omega$ is dense because $\Omega$ is cyclic for $\mathcal{M}$ (\ref{def:cyclic-separating-vector}).
\end{proof}

\begin{lemma}[The Adjoint of $a_0$]
    \label{lmm:tomita-adjoint-left-map-adjoint}
    \lean{Physicslib4.Modular.tomitaAdjoint_leftMap_adjoint}
    \uses{lmm:tomita-adjoint-left-map, def:hall-9.1, def:tomita-left-map}
    \leanok
    For every $y \in \mathcal{M}$, $y\Omega \in \text{Dom}(a_0^*)$ and $a_0^*(y\Omega) = y\eta$. In particular $a_0^*$ is densely defined.
\end{lemma}
\begin{proof}
    \leanok
    \uses{lmm:tomita-adjoint-left-map, lmm:tomita-adjoint-intrinsic-core, lmm:characterizing-adjoint-domain-membership, lmm:commutant-von-neumann-algebra}
    Let $x \in \mathcal{M}$. Then $x^*y \in \mathcal{M}$ and $(y^*x)^* = x^*y$ (Lemma~\ref{lmm:commutant-von-neumann-algebra}), so \eqref{eq:tomita-adjoint-pair} for $y^*x$ gives
    \begin{align}
        \left< a_0(x\Omega), y\Omega \right> = \left< x\xi, y\Omega \right> = \left< \xi, x^*y\Omega \right> = \left< y^*x\Omega, \eta \right> = \left< x\Omega, y\eta \right>.
    \end{align}
    Taking complex conjugates, $\left< y\Omega, a_0(x\Omega) \right> = \left< y\eta, x\Omega \right>$ for every $x \in \mathcal{M}$, so Lemma~\ref{lmm:characterizing-adjoint-domain-membership} gives the claim. $\mathcal{M}\Omega$ is dense (Lemma~\ref{lmm:tomita-adjoint-left-map}).
\end{proof}

\begin{lemma}[The Closure $c$ of $a_0$]
    \label{lmm:tomita-adjoint-closure-props}
    \lean{Physicslib4.Modular.tomitaAdjoint_leftMap_closure_props}
    \uses{lmm:tomita-adjoint-left-map, def:hall-9.6, def:hall-9.1, def:tomita-left-map}
    \leanok
    $a_0$ is closable. Its closure $c = a_0^{\text{cl}}$ is closed and densely defined, $\mathcal{M}\Omega \subset \text{Dom}(c)$ with $c(x\Omega) = x\xi$ for $x \in \mathcal{M}$, so $c\Omega = \xi$; and $\Omega \in \text{Dom}(c^*)$ with $c^*\Omega = \eta$.
\end{lemma}
\begin{proof}
    \leanok
    \uses{lmm:tomita-adjoint-left-map-adjoint, lmm:tomita-adjoint-left-map, lmm:closable-of-dense-adjoint, prpstn:hall-9.10}
    $a_0^*$ is densely defined (Lemma~\ref{lmm:tomita-adjoint-left-map-adjoint}), so $a_0$ is closable (Lemma~\ref{lmm:closable-of-dense-adjoint}); its closure is closed, extends $a_0$, and so is densely defined (Lemma~\ref{lmm:tomita-adjoint-left-map}), and $\mathbf{1} \in \mathcal{M}$ gives $c\Omega = \xi$. By Proposition~\ref{prpstn:hall-9.10}, $c^* = a_0^*$, so $c^*\Omega = a_0^*(\mathbf{1}\Omega) = \eta$ by Lemma~\ref{lmm:tomita-adjoint-left-map-adjoint} with $y = \mathbf{1}$.
\end{proof}

\begin{lemma}[The Unitaries of $\mathcal{M}$ Commute with $c$]
    \label{lmm:tomita-adjoint-closure-commutes}
    \lean{Physicslib4.Modular.tomitaAdjoint_leftMap_closure_commute}
    \uses{lmm:tomita-adjoint-closure-props, def:unitary-operator, def:unitary-commutes}
    \leanok
    Every unitary $u \in \mathcal{M}$ commutes with $c$.
\end{lemma}
\begin{proof}
    \leanok
    \uses{lmm:tomita-adjoint-closure-props, prpstn:closure-linearity-and-sequential-description, lmm:commutant-von-neumann-algebra}
    $u^{-1} = u^*$ is also a unitary in $\mathcal{M}$ (Lemma~\ref{lmm:commutant-von-neumann-algebra}), so it suffices to show, for every unitary $w \in \mathcal{M}$ and $\psi \in \text{Dom}(c)$, that $w\psi \in \text{Dom}(c)$ and $cw\psi = wc\psi$. Take $y_n \in \mathcal{M}$ with $y_n\Omega \to \psi$ and $y_n\xi \to c\psi$ (Proposition~\ref{prpstn:closure-linearity-and-sequential-description} and Lemma~\ref{lmm:tomita-adjoint-closure-props}). Then $wy_n \in \mathcal{M}$, $wy_n\Omega \to w\psi$ and $c(wy_n\Omega) = wy_n\xi \to wc\psi$ by continuity of $w$; since $c$ is closed, $w\psi \in \text{Dom}(c)$ and $cw\psi = wc\psi$.
\end{proof}

\begin{lemma}[The Approximants Lie in $\mathcal{M}'$]
    \label{lmm:tomita-adjoint-approximants-commutant}
    \lean{Physicslib4.Modular.tomitaAdjoint_approximant_mem_centralizer}
    \uses{lmm:tomita-adjoint-closure-props, lmm:approximant-operator}
    \leanok
    With $c$ as in Lemma~\ref{lmm:tomita-adjoint-closure-props} and $x_n$ the operators of Lemma~\ref{lmm:approximant-operator} for this $c$, $x_n \in \mathcal{M}'$ for every $n \ge 1$.
\end{lemma}
\begin{proof}
    \leanok
    \uses{lmm:tomita-adjoint-closure-commutes, lmm:approximant-operator-commutes, lmm:commutant-of-unitaries}
    Every unitary $u \in \mathcal{M}$ commutes with $c$ (Lemma~\ref{lmm:tomita-adjoint-closure-commutes}), so $ux_n = x_nu$ (Lemma~\ref{lmm:approximant-operator-commutes}); apply Lemma~\ref{lmm:commutant-of-unitaries}.
\end{proof}

\begin{lemma}[Convergence of the Approximants on $\Omega$]
    \label{lmm:tomita-adjoint-approximants-converge}
    \lean{Physicslib4.Modular.tomitaAdjoint_approximant_tendsto}
    \uses{lmm:tomita-adjoint-closure-props, lmm:approximant-operator}
    \leanok
    With $c$ and $x_n$ as in Lemma~\ref{lmm:tomita-adjoint-approximants-commutant}, $x_n\Omega \to \xi$ and $x_n^*\Omega \to \eta$ as $n \to \infty$.
\end{lemma}
\begin{proof}
    \leanok
    \uses{lmm:tomita-adjoint-closure-props, lmm:approximant-operator, lmm:approximant-range, lmm:approximant-convergence, thrm:absolute-value-domain}
    $\Omega \in \text{Dom}(c)$ with $c\Omega = \xi$, and $\Omega \in \text{Dom}(c^*)$ with $c^*\Omega = \eta$ (Lemma~\ref{lmm:tomita-adjoint-closure-props}). Since $R_n\Omega - \Omega \in \text{Dom}(c)$ (Lemma~\ref{lmm:approximant-range}), Theorem~\ref{thrm:absolute-value-domain} and Lemma~\ref{lmm:approximant-range} give
    \begin{align}
        \|x_n\Omega - \xi\| = \|c(R_n\Omega - \Omega)\| = \||c|(R_n\Omega - \Omega)\| = \|K_n\Omega - |c|\Omega\|,
    \end{align}
    which tends to $0$ by Lemma~\ref{lmm:approximant-convergence}. By Lemma~\ref{lmm:approximant-operator}, $x_n^*\Omega = R_nc^*\Omega = R_n\eta$, which tends to $\eta$ by Lemma~\ref{lmm:approximant-convergence}.
\end{proof}

\begin{theorem}[$S^* = \overline{F_0}$]
    \label{thrm:tomita-adjoint-eq-closure}
    \lean{Physicslib4.Modular.mem_tomitaAdjoint_iff_mem_closure_graph}
    \uses{def:tomita-operator, def:tomita-pre-operators, def:hall-9.1}
    \leanok
    Let $\xi, \eta \in \mathbf{H}$. Then $C\xi \in \text{Dom}((S')^*)$ and $(S')^*C\xi = \eta$ if and only if $(\xi, \eta)$ lies in the closure of $\Gamma(F_0) = \{ (b'\Omega, b'^*\Omega) \mid b' \in \mathcal{M}' \}$ in $\mathbf{H} \times \mathbf{H}$. That is, the antilinear adjoint $S^*$ of $S$ is the closure of $F_0$, and $\mathcal{M}'\Omega$ is a core for $S^*$.
\end{theorem}
\begin{proof}
    \leanok
    \uses{lmm:tomita-adjoint-extends-closure, lmm:tomita-adjoint-approximants-commutant, lmm:tomita-adjoint-approximants-converge, lmm:componentwise-convergence, thrm:sequential-closedness}
    The reverse implication is Lemma~\ref{lmm:tomita-adjoint-extends-closure}. For the forward one, the operators $x_n$ of Lemma~\ref{lmm:tomita-adjoint-approximants-commutant} lie in $\mathcal{M}'$, so $(x_n\Omega, x_n^*\Omega) \in \Gamma(F_0)$, and $(x_n\Omega, x_n^*\Omega) \to (\xi, \eta)$ by Lemmas~\ref{lmm:tomita-adjoint-approximants-converge} and~\ref{lmm:componentwise-convergence}. Hence $(\xi, \eta)$ lies in the closure of $\Gamma(F_0)$ (Theorem~\ref{thrm:sequential-closedness}).
\end{proof}

\subsubsection{The Polar Decomposition of $S^*$}

\begin{lemma}[$S^* = \Delta^{1/2}J$]
    \label{lmm:tomita-adjoint-sqrt-conjugation}
    \lean{Physicslib4.Modular.mem_tomitaAdjoint_iff_modularOperatorSqrt}
    \uses{def:tomita-operator, def:modular-operator, def:modular-conjugation}
    \leanok
    Let $\xi, \eta \in \mathbf{H}$. Then $C\xi \in \text{Dom}((S')^*)$ and $(S')^*C\xi = \eta$ if and only if $J\xi \in \text{Dom}(\Delta^{1/2})$ and $\Delta^{1/2}J\xi = \eta$.
\end{lemma}
\begin{proof}
    \leanok
    \uses{lmm:tomita-adjoint-intrinsic, thrm:modular-conjugation-polar, thrm:modular-conjugation-involution, lmm:antilinear-isometry-antiunitary, lmm:square-root-positive, thrm:modular-operator-properties, lmm:characterizing-adjoint-domain-membership}
    By Theorem~\ref{thrm:modular-conjugation-polar}, $\text{Dom}(S) = \text{Dom}(\Delta^{1/2})$ and $S\zeta = J\Delta^{1/2}\zeta$. Since $J(J\xi) = \xi$ (Theorem~\ref{thrm:modular-conjugation-involution}) and $J$ is antiunitary (Lemma~\ref{lmm:antilinear-isometry-antiunitary}), $\left< \xi, J\Delta^{1/2}\zeta \right> = \overline{\left< J\xi, \Delta^{1/2}\zeta \right>}$. So the condition of Lemma~\ref{lmm:tomita-adjoint-intrinsic} becomes $\left< J\xi, \Delta^{1/2}\zeta \right> = \left< \eta, \zeta \right>$ for all $\zeta \in \text{Dom}(\Delta^{1/2})$, which by Lemma~\ref{lmm:characterizing-adjoint-domain-membership} says $J\xi \in \text{Dom}((\Delta^{1/2})^*)$ with $(\Delta^{1/2})^*J\xi = \eta$; and $(\Delta^{1/2})^* = \Delta^{1/2}$ (Lemma~\ref{lmm:square-root-positive} and Theorem~\ref{thrm:modular-operator-properties}).
\end{proof}

\begin{corollary}[$S^* = J\Delta^{-1/2}$]
    \label{crllr:tomita-adjoint-polar}
    \lean{Physicslib4.Modular.mem_tomitaAdjoint_iff_modularOperatorInvSqrt}
    \uses{def:tomita-operator, def:modular-operator, def:modular-conjugation, lmm:inverse-square-root-positive}
    \leanok
    Let $\xi, \eta \in \mathbf{H}$. Then $C\xi \in \text{Dom}((S')^*)$ and $(S')^*C\xi = \eta$ if and only if $\xi \in \text{Dom}(\Delta^{-1/2})$ and $\eta = J\Delta^{-1/2}\xi$. In particular the closure of $F_0$ is $J\Delta^{-1/2}$ (Theorem~\ref{thrm:tomita-adjoint-eq-closure}).
\end{corollary}
\begin{proof}
    \leanok
    \uses{lmm:tomita-adjoint-sqrt-conjugation, thrm:modular-conjugation-involution, def:product-of-unbounded-operators}
    By Theorem~\ref{thrm:modular-conjugation-involution}, $J\Delta^{1/2}J = \Delta^{-1/2}$ including domains (\ref{def:product-of-unbounded-operators}): $J\xi \in \text{Dom}(\Delta^{1/2})$ if and only if $\xi \in \text{Dom}(\Delta^{-1/2})$, and then $J\Delta^{1/2}J\xi = \Delta^{-1/2}\xi$, so $\Delta^{1/2}J\xi = J\Delta^{-1/2}\xi$ because $J(J\phi) = \phi$. Combine with Lemma~\ref{lmm:tomita-adjoint-sqrt-conjugation}.
\end{proof}

\subsubsection{The Modular Objects of the Commutant}

Write $S_{\mathcal{M}'}$, $\Delta_{\mathcal{M}'}$, $\Delta_{\mathcal{M}'}^{1/2}$ and $J_{\mathcal{M}'}$ for the Tomita operator, the modular operator, its square root and the modular conjugation of $(\mathcal{M}', \Omega)$ (Definitions~\ref{def:tomita-operator}, \ref{def:modular-operator} and~\ref{def:modular-conjugation}, with the same reference conjugation $C$; by Lemma~\ref{lmm:tomita-operator-independent-of-c} and Theorem~\ref{thrm:modular-objects-independent-of-c} the choice does not matter). They are defined by Lemma~\ref{lmm:commutant-cyclic-separating}.

\begin{lemma}[$\Omega$ is Cyclic and Separating for $\mathcal{M}'$]
    \label{lmm:commutant-cyclic-separating}
    \lean{Physicslib4.Modular.isCyclicSeparatingVector_centralizer}
    \uses{def:cyclic-separating-vector, def:von-neumann-algebra-commutant-convention}
    \leanok
    $\mathcal{M}'$ is a von Neumann algebra on $\mathbf{H}$, and $\Omega$ is cyclic and separating for $\mathcal{M}'$.
\end{lemma}
\begin{proof}
    \leanok
    \uses{lmm:commutant-von-neumann-algebra}
    $\mathcal{M}'$ is a von Neumann algebra and $(\mathcal{M}')' = \mathcal{M}$ (Lemma~\ref{lmm:commutant-von-neumann-algebra}). So the two density conditions of \ref{def:cyclic-separating-vector} for $\mathcal{M}'$, density of $\mathcal{M}'\Omega$ and of $(\mathcal{M}')'\Omega = \mathcal{M}\Omega$, are those for $\mathcal{M}$ in the other order.
\end{proof}

\begin{lemma}[The Tomita Operator of $\mathcal{M}'$ is $S^*$]
    \label{lmm:commutant-tomita-operator}
    \lean{Physicslib4.Modular.tomitaOperator_centralizer_iff}
    \uses{lmm:commutant-cyclic-separating, def:tomita-operator}
    \leanok
    Let $\xi, \eta \in \mathbf{H}$. Then $\xi \in \text{Dom}(S_{\mathcal{M}'})$ and $S_{\mathcal{M}'}\xi = \eta$ if and only if $\xi \in \text{Dom}(\Delta^{-1/2})$ and $\eta = J\Delta^{-1/2}\xi$, if and only if $C\xi \in \text{Dom}((S')^*)$ and $(S')^*C\xi = \eta$.
\end{lemma}
\begin{proof}
    \leanok
    \uses{lmm:commutant-cyclic-separating, lmm:tomita-operator-independent-of-c, lmm:tomita-pre-operators-apply, thrm:tomita-adjoint-eq-closure, crllr:tomita-adjoint-polar}
    By Lemma~\ref{lmm:tomita-pre-operators-apply} for $(\mathcal{M}', \Omega)$, the graph of its pre-operator $S_0$ is $\{ (b'\Omega, b'^*\Omega) \mid b' \in \mathcal{M}' \} = \Gamma(F_0)$, so Lemma~\ref{lmm:tomita-operator-independent-of-c} for $(\mathcal{M}', \Omega)$ identifies $\Gamma(S_{\mathcal{M}'})$ with the closure of $\Gamma(F_0)$. Theorem~\ref{thrm:tomita-adjoint-eq-closure} and Corollary~\ref{crllr:tomita-adjoint-polar} give the two equivalences.
\end{proof}

\begin{lemma}[An Intrinsic Description of $\Delta_{\mathcal{M}'}$]
    \label{lmm:commutant-modular-operator-intrinsic}
    \lean{Physicslib4.Modular.mem_modularOperator_centralizer_iff}
    \uses{lmm:commutant-cyclic-separating, def:modular-operator, lmm:inverse-square-root-positive}
    \leanok
    $\psi \in \text{Dom}(\Delta_{\mathcal{M}'})$ with $\Delta_{\mathcal{M}'}\psi = \chi$ if and only if $\psi \in \text{Dom}(\Delta^{-1/2})$ and $\left< \Delta^{-1/2}\psi, \Delta^{-1/2}\phi \right> = \left< \chi, \phi \right>$ for every $\phi \in \text{Dom}(\Delta^{-1/2})$.
\end{lemma}
\begin{proof}
    \leanok
    \uses{lmm:modular-operator-intrinsic, lmm:commutant-tomita-operator, lmm:antilinear-isometry-antiunitary, thrm:modular-conjugation-polar}
    Lemma~\ref{lmm:modular-operator-intrinsic} for $(\mathcal{M}', \Omega)$ describes $\Delta_{\mathcal{M}'}$ by $\psi \in \text{Dom}(S_{\mathcal{M}'})$ and $\left< S_{\mathcal{M}'}\phi, S_{\mathcal{M}'}\psi \right> = \left< \chi, \phi \right>$ for $\phi \in \text{Dom}(S_{\mathcal{M}'})$. By Lemma~\ref{lmm:commutant-tomita-operator}, $\text{Dom}(S_{\mathcal{M}'}) = \text{Dom}(\Delta^{-1/2})$ and $S_{\mathcal{M}'} = J\Delta^{-1/2}$ there, and $J$ is antiunitary (Theorem~\ref{thrm:modular-conjugation-polar} and Lemma~\ref{lmm:antilinear-isometry-antiunitary}), so $\left< J\Delta^{-1/2}\phi, J\Delta^{-1/2}\psi \right> = \left< \Delta^{-1/2}\psi, \Delta^{-1/2}\phi \right>$.
\end{proof}

\begin{lemma}[The Domain of $f(A)f(A)$]
    \label{lmm:fc-domain-square}
    \lean{Physicslib4.Spectral.Unbounded.pmapMul_functionalCalculus_self}
    \uses{def:hall-10.5, def:product-of-unbounded-operators}
    \leanok
    Let $A$ be self-adjoint and $f : \mathbb{R} \to \mathbb{R}$ measurable. Then $W_{f^2} \subset W_f$, and $f(A)f(A) = (f^2)(A)$, including domains.
\end{lemma}
\begin{proof}
    \leanok
    \uses{prpstn:hall-10.2, def:projection-valued-measure, thrm:fc-composition}
    Let $\psi \in W_{f^2}$. Since $|f|^2 \le 1 + |f|^4$ and $\mu^A_\psi(\mathbb{R}) = \|\psi\|^2 < \infty$ (Definition~\ref{def:projection-valued-measure}), $\int |f|^2 \, d\mu^A_\psi \le \|\psi\|^2 + \int |f^2|^2 \, d\mu^A_\psi < \infty$, so $\psi \in W_f$ (Proposition~\ref{prpstn:hall-10.2}). Theorem~\ref{thrm:fc-composition} with $g = f$ gives $f(A)f(A) \subset (f^2)(A)$ with $\text{Dom}(f(A)f(A)) = W_f \cap W_{f^2} = W_{f^2} = \text{Dom}((f^2)(A))$, so the two operators are equal.
\end{proof}

\begin{lemma}[$A^{-1/2}A^{-1/2} = A^{-1}$]
    \label{lmm:inverse-square-root-product}
    \lean{Physicslib4.Spectral.Unbounded.pmapMul_inv_sqrt_functionalCalculus}
    \uses{lmm:inverse-square-root-positive, def:inverse-of-operator, def:product-of-unbounded-operators}
    \leanok
    Let $A$ be positive, self-adjoint and injective. Then the product $A^{-1/2}A^{-1/2}$ equals $A^{-1}$, including domains.
\end{lemma}
\begin{proof}
    \leanok
    \uses{lmm:fc-domain-square, crllr:fc-of-fc, lmm:square-of-inverse-square-root}
    With $r$ as in Lemma~\ref{lmm:inverse-square-root-positive}, real and measurable, $A^{-1/2} = r(A)$, and Lemma~\ref{lmm:fc-domain-square} gives $r(A)r(A) = (r^2)(A)$. With $q(\lambda) = \lambda^2$, $(r^2)(A) = (q \circ r)(A) = q(r(A))$ by Corollary~\ref{crllr:fc-of-fc}, and $q(A^{-1/2}) = A^{-1}$ by Lemma~\ref{lmm:square-of-inverse-square-root}.
\end{proof}

\begin{lemma}[$\Delta_{\mathcal{M}'} = \Delta^{-1}$]
    \label{lmm:commutant-modular-operator}
    \lean{Physicslib4.Modular.modularOperator_centralizer}
    \uses{lmm:commutant-cyclic-separating, def:modular-operator, def:inverse-of-operator}
    \leanok
    $\Delta_{\mathcal{M}'} = \Delta^{-1}$, including domains.
\end{lemma}
\begin{proof}
    \leanok
    \uses{lmm:commutant-modular-operator-intrinsic, lmm:characterizing-adjoint-domain-membership, lmm:inverse-square-root-positive, lmm:inverse-square-root-product, thrm:modular-operator-properties}
    Write $P = \Delta^{-1/2}$, which is self-adjoint (Lemma~\ref{lmm:inverse-square-root-positive}). By Lemma~\ref{lmm:characterizing-adjoint-domain-membership} and $P^* = P$, the condition of Lemma~\ref{lmm:commutant-modular-operator-intrinsic} says $\psi \in \text{Dom}(P)$, $P\psi \in \text{Dom}(P)$ and $PP\psi = \chi$, that is $\psi \in \text{Dom}(PP)$ and $PP\psi = \chi$. $PP = \Delta^{-1}$ by Lemma~\ref{lmm:inverse-square-root-product} for the positive, self-adjoint, injective $\Delta$ (Theorem~\ref{thrm:modular-operator-properties}).
\end{proof}

\begin{lemma}[The Square Root of the Inverse]
    \label{lmm:sqrt-of-inverse}
    \lean{Physicslib4.Spectral.Unbounded.sqrt_functionalCalculus_inverse}
    \uses{lmm:inverse-square-root-positive, def:inverse-of-operator}
    \leanok
    Let $A$ be positive, self-adjoint and injective. Then $A^{-1}$ is self-adjoint and $(A^{-1})^{1/2} = s(A^{-1}) = A^{-1/2}$.
\end{lemma}
\begin{proof}
    \leanok
    \uses{lmm:inverse-functional-calculus, crllr:fc-of-fc}
    $A^{-1} = \iota^-(A)$ (Lemma~\ref{lmm:inverse-functional-calculus}) with $\iota^-$ real and measurable (Mathlib's \texttt{measurable\_inv}), so Corollary~\ref{crllr:fc-of-fc} gives that $A^{-1}$ is self-adjoint and $s(\iota^-(A)) = (s \circ \iota^-)(A)$. Here $\iota^-$ is Mathlib's \texttt{Inv.inv} on $\mathbb{R}$, with $0^{-1} = 0$. For $\lambda > 0$, $s(\lambda^{-1}) = (\sqrt{\lambda})^{-1} = (\lambda^{1/2})^{-1} = \lambda^{-1/2} = r(\lambda)$ (Mathlib's \texttt{Real.sqrt\_inv}, \texttt{Real.sqrt\_eq\_rpow} and \texttt{Real.rpow\_neg}); for $\lambda \le 0$, $\iota^-(\lambda) \le 0$, so $s(\iota^-(\lambda)) = 0 = r(\lambda)$. Hence $s \circ \iota^- = r$.
\end{proof}

\begin{lemma}[$J_{\mathcal{M}'} = J$]
    \label{lmm:commutant-modular-conjugation}
    \lean{Physicslib4.Modular.modularOperatorSqrt_modularConjugation_centralizer}
    \uses{lmm:commutant-cyclic-separating, def:modular-conjugation}
    \leanok
    $\Delta_{\mathcal{M}'}^{1/2} = \Delta^{-1/2}$ and $J_{\mathcal{M}'} = J$.
\end{lemma}
\begin{proof}
    \uses{lmm:commutant-modular-operator, lmm:sqrt-of-inverse, thrm:modular-operator-properties, lmm:modular-conjugation-intrinsic, lmm:commutant-tomita-operator, lmm:antiunitary-isometric}
    $\Delta_{\mathcal{M}'}^{1/2} = s(\Delta_{\mathcal{M}'}) = s(\Delta^{-1}) = \Delta^{-1/2}$ by Lemmas~\ref{lmm:commutant-modular-operator} and~\ref{lmm:sqrt-of-inverse} (Theorem~\ref{thrm:modular-operator-properties}). By Lemma~\ref{lmm:modular-conjugation-intrinsic} for $(\mathcal{M}', \Omega)$, $J_{\mathcal{M}'}$ is the only continuous map with $J_{\mathcal{M}'}\Delta_{\mathcal{M}'}^{1/2}\psi = S_{\mathcal{M}'}\psi$ for $\psi \in \text{Dom}(\Delta_{\mathcal{M}'}^{1/2})$. $J$ is continuous (Lemma~\ref{lmm:antiunitary-isometric}) and $J\Delta^{-1/2}\psi = S_{\mathcal{M}'}\psi$ for $\psi \in \text{Dom}(\Delta^{-1/2})$ (Lemma~\ref{lmm:commutant-tomita-operator}), so $J_{\mathcal{M}'} = J$.
\end{proof}

\begin{theorem}[The Modular Objects of $(\mathcal{M}', \Omega)$]
    \label{thrm:commutant-modular-objects}
    \lean{Physicslib4.Modular.commutant_modular_objects}
    \uses{lmm:commutant-cyclic-separating, def:tomita-operator, def:modular-operator, def:modular-conjugation, def:inverse-of-operator}
    \leanok
    Let $\mathcal{M}$ be a von Neumann algebra on $\mathbf{H}$ and $\Omega$ cyclic and separating for $\mathcal{M}$, with Tomita operator $S$, modular operator $\Delta$ and modular conjugation $J$. Then $\Omega$ is cyclic and separating for $\mathcal{M}'$, and the Tomita operator, the modular operator and the modular conjugation of $(\mathcal{M}', \Omega)$ are
    \begin{align}
        S_{\mathcal{M}'} = S^* = J\Delta^{-1/2}, \qquad \Delta_{\mathcal{M}'} = \Delta^{-1}, \qquad \Delta_{\mathcal{M}'}^{1/2} = \Delta^{-1/2}, \qquad J_{\mathcal{M}'} = J,
    \end{align}
    where $S_{\mathcal{M}'} = S^*$ means $\Gamma(S_{\mathcal{M}'}) = \Gamma(S^*)$, and the equalities of operators include domains.
\end{theorem}
\begin{proof}
    \uses{lmm:commutant-cyclic-separating, lmm:commutant-tomita-operator, lmm:commutant-modular-operator, lmm:commutant-modular-conjugation}
    Lemma~\ref{lmm:commutant-cyclic-separating} gives the hypotheses for $(\mathcal{M}', \Omega)$. The four identities are Lemmas~\ref{lmm:commutant-tomita-operator}, \ref{lmm:commutant-modular-operator} and~\ref{lmm:commutant-modular-conjugation}.
\end{proof}
